% Môn: Vật lý
% Lớp: 11
% Chương: 
% Bài: 
% Số câu: 178
% App đọc file này trực tiếp — sửa rồi Commit trên GitHub, bấm Đồng bộ GitHub .tex

% _Chua_phan_loai/de.tex

% ===== Câu 1 =====
% ID: L11CXBX-00-DS
% Mức: TH
\begin{ex}
% ID: L11CXBX-00-DS
% Mức: TH
Xét các mệnh đề sau về công và công suất:
\choiceTF

\loigiai{
A. Đúng, vì công suất $P = \frac{A}{t}$.
B. Sai, công suất có đơn vị là watt.
C. Đúng, vì công suất lớn nghĩa là công thực hiện nhanh.
D. Đúng, công âm khi lực ngược chiều chuyển động.
}
\end{ex}

% ===== Câu 2 =====
% ID: L11CXBX-00-TLN
% Mức: VD
\dangbt{AI tạo tự động}
\begin{ex}
% ID: L11CXBX-00-TLN
% Mức: VD
Một lò xo có độ cứng $150\,\text{N}$ /m bị giãn $0,1\,\text{m}$. Tính thế năng đàn hồi của lò xo.
\shortans{0,75}
\loigiai{
Thế năng đàn hồi của lò xo được tính bằng công thức: 
 $W = \frac{1}{2} k x^2$
 Trong đó, $k = 150 \, \text{N/m}$ là độ cứng của lò xo và $x = 0,1 \, \text{m}$ là độ giãn của lò xo. Thay các giá trị vào công thức, ta có:
 $W = \frac{1}{2} \times 150 \times (0,1)^2W = \frac{1}{2} \times 150 \times 0,01W = \frac{1}{2} \times 1,5W = 0,75 \, \text{J}$
 
 Kiểm tra từng bước tính toán:
 - Tính $(0,1)^2 = 0,01$.
 - Nhân $150 \times 0,01 = 1,5$.
 - Tính $\frac{1}{2} \times 1,5 = 0,75$.
}
\end{ex}

% ===== Câu 3 =====
% ID: L11CXBX-00-TLN
% Mức: VD
\begin{ex}
% ID: L11CXBX-00-TLN
% Mức: VD
Một vật có khối lượng $3\,\text{kg}$ được thả rơi tự do từ độ cao $20\,\text{m}$. Tính vận tốc của vật khi chạm đất. (lấy $g = 9,8 \text{ m/s}^2$)
\shortans{19,8}
\loigiai{
Chưa có lời giải. ADMIN cần kiểm tra và bổ sung.
}
\end{ex}

% ===== Câu 4 =====
% ID: LATEX_0DB1EA144BE3
% Mức: TH
\dangbt{Áp dụng quy tắc ba điểm và quy tắc hình bình hành}
\begin{ex}
% ID: LATEX_0DB1EA144BE3
% Mức: TH
Cho ba điểm phân biệt $\A, B, C\. Đẳng thức nào sau đây đúng?$
\choice
{$\\overrightarrow{AB}$ + $\overrightarrow{CB}$ = $\overrightarrow{AC}\$}
{$\\overrightarrow{CB}$ + $\overrightarrow{AC}$ = $\overrightarrow{AB}\$}
{$\\overrightarrow{AB}$ + $\overrightarrow{BC}$ = $\overrightarrow{CA}\$}
{\True $\\overrightarrow{AB}$ + $\overrightarrow{BC}$ = $\overrightarrow{AC}\$}
\loigiai{
Theo quy tắc ba điểm đối với phép cộng vectơ, với ba điểm $\A, B, C\bất kì ta luôn có\\overrightarrow{AB} + \overrightarrow{BC} = \overrightarrow{AC}\$.
}
\end{ex}

% ===== Câu 5 =====
% ID: LATEX_AF6B2F66DC9A
% Mức: TH
\begin{ex}
% ID: LATEX_AF6B2F66DC9A
% Mức: TH
Cho ba điểm phân biệt $\M, N, P\. Vectơ tổng\\overrightarrow{NP} + \overrightarrow{MN}\$ bằng vectơ nào sau đây?
\choice
{$\\overrightarrow{PN}\$}
{$\\overrightarrow{PM}\$}
{\True $\\overrightarrow{MP}\$}
{$\\overrightarrow{NM}\$}
\loigiai{
• Áp dụng tính chất giao hoán của phép cộng vectơ, ta có: $\\overrightarrow{NP}$ + $\overrightarrow{MN}$ = $\overrightarrow{MN}$ + $\overrightarrow{NP}\$.
 
• Theo quy tắc ba điểm, $\$ \overrightarrow{MN} + \overrightarrow{NP} = \overrightarrow{MP}\ $.
 
• Vậy$overrightarrow{NP}$+$\overrightarrow{MN}$=$ \overrightarrow{MP}\.
}
\end{ex}

% ===== Câu 6 =====
% ID: LATEX_B3959566D3C6
% Mức: TH
\begin{ex}
% ID: LATEX_B3959566D3C6
% Mức: TH
Cho hình bình hành $\$ ABCD $\. Đẳng thức nào sau đây đúng?$
\choice
{\True $\\overrightarrow{AB}$ + $\overrightarrow{AD}$ = $\overrightarrow{AC}\$}
{$\\overrightarrow{AB}$ + $\overrightarrow{AD}$ = $\overrightarrow{BD}\$}
{$\\overrightarrow{AB}$ + $\overrightarrow{AD}$ = $\overrightarrow{CD}\$}
{$\\overrightarrow{AB}$ + $\overrightarrow{AD}$ = $\overrightarrow{CB}\$}
\loigiai{
Theo quy tắc hình bình hành, vì $\$ ABCD $\là hình bình hành nên tổng hai vectơ cạnh chung gốc bằng vectơ đường chéo xuất phát từ gốc đó, tức là\\overrightarrow{AB} + \overrightarrow{AD} = \overrightarrow{AC}\$.
}
\end{ex}

% ===== Câu 7 =====
% ID: LATEX_1431DB984C9E
% Mức: TH
\begin{ex}
% ID: LATEX_1431DB984C9E
% Mức: TH
Cho bốn điểm phân biệt $\A, B, C, O\. Đẳng thức nào sau đây đúng?$
\choice
{$\\overrightarrow{AB}=\overrightarrow{OB}$ + $\overrightarrow{OA}\$}
{$\\overrightarrow{AB}=\overrightarrow{AC}$ + $\overrightarrow{BC}\$}
{\True $\\overrightarrow{OA}=\overrightarrow{CA}$ - $\overrightarrow{CO}\$}
{$\\overrightarrow{OA}=\overrightarrow{OB}$ - $\overrightarrow{BA}\$}
\loigiai{
• Áp dụng quy tắc trừ vectơ với ba điểm $\O, C, A\ta có:\\overrightarrow{CA} - \overrightarrow{CO} = \overrightarrow{OA}\$.
 
• Các phương án khác đều sai quy tắc cộng và trừ vectơ.
}
\end{ex}

% ===== Câu 8 =====
% ID: LATEX_A46C88793698
% Mức: TH
\begin{ex}
% ID: LATEX_A46C88793698
% Mức: TH
Cho bốn điểm phân biệt $\A, B, C, D\. Vectơ tổng\\overrightarrow{AB} + \overrightarrow{CD} + \overrightarrow{BC} + \overrightarrow{DA}\$ bằng vectơ nào sau đây?
\choice
{\True $\\overrightarrow{0}\$}
{$\\overrightarrow{AC}\$}
{$\\overrightarrow{BD}\$}
{$\\overrightarrow{BA}\$}
\loigiai{
• Ta sắp xếp lại các số hạng và áp dụng quy tắc ba điểm liên tiếp: $\\overrightarrow{AB}$ + $\overrightarrow{CD}$ + $\overrightarrow{BC}$ + $\overrightarrow{DA}$ = $\overrightarrow{AB}$ + $\overrightarrow{BC}$ + $\overrightarrow{CD}$ + $\overrightarrow{DA}\$.
 
• Thực hiện phép cộng từ trái sang phải: $\$ \overrightarrow{AC} + \overrightarrow{CD} + \overrightarrow{DA} = \overrightarrow{AD} + \overrightarrow{DA} = \overrightarrow{AA}\ $.
 
• Vì điểm đầu và điểm cuối trùng nhau nên$overrightarrow{AA}$=$ \overrightarrow{0}\.
}
\end{ex}

% ===== Câu 9 =====
% ID: LATEX_4302F8D17026
% Mức: TH
\dangbt{Rút gọn biểu thức và chứng minh đẳng thức vectơ}
\begin{ex}
% ID: LATEX_4302F8D17026
% Mức: TH
Cho bốn điểm $\A, B, C, D\bất kì. Khẳng định nào sau đây là đúng?$
\choice
{$\\overrightarrow{AB}$ - $\overrightarrow{AD}$ = $\overrightarrow{CD}$ + $\overrightarrow{BC}\$}
{$\\overrightarrow{AB}$ - $\overrightarrow{AC}$ = $\overrightarrow{DB}$ - $\overrightarrow{DC}\$}
{$\\overrightarrow{BC}$ + $\overrightarrow{AB}$ = $\overrightarrow{DA}$ - $\overrightarrow{DC}\$}
{\True $\\overrightarrow{AB}$ + $\overrightarrow{CD}$ = $\overrightarrow{AD}$ + $\overrightarrow{CB}\$}
\loigiai{
• Xét vế trái: $\\overrightarrow{AB} + \overrightarrow{CD}\. Áp dụng quy tắc chèn điểm\,\text{D}\$ vào vectơ thứ nhất và điểm $\B$ \ $vào vectơ thứ hai.
 
• Ta có$ \ $\overrightarrow{AB} + \overrightarrow{CD} = (\overrightarrow{AD} + \overrightarrow{DB}) + (\overrightarrow{CB} + \overrightarrow{BD})\$.
 
• Nhóm lại ta được: $\\overrightarrow{AD}$ + $\overrightarrow{CB}$ + $(\overrightarrow{DB} + \overrightarrow{BD})$ = $\overrightarrow{AD}$ + $\overrightarrow{CB}$ + $\overrightarrow{0}$ = $\overrightarrow{AD}$ + $\overrightarrow{CB}\. Đẳng thức đúng.$
}
\end{ex}

% ===== Câu 10 =====
% ID: LATEX_3BCD556B4A93
% Mức: TH
\begin{ex}
% ID: LATEX_3BCD556B4A93
% Mức: TH
Rút gọn biểu thức vectơ $\\overrightarrow{P} = \overrightarrow{AB} - \overrightarrow{AC} + \overrightarrow{BD}\, ta được kết quả nào sau đây?$
\choice
{$\\overrightarrow{AD}\$}
{$\\overrightarrow{BC}\$}
{\True $\\overrightarrow{CD}\$}
{$\\overrightarrow{DA}\$}
\loigiai{
• Áp dụng quy tắc hiệu cho hai vectơ đầu tiên: $\\overrightarrow{AB}$ - $\overrightarrow{AC}$ = $\overrightarrow{CB}\$.
 
• Thay vào biểu thức ta được: $\$ \overrightarrow{P} = \overrightarrow{CB} + \overrightarrow{BD}\ $.
 
• Áp dụng quy tắc ba điểm:$overrightarrow{CB}$+$\overrightarrow{BD}$=$ \overrightarrow{CD}\.
}
\end{ex}

% ===== Câu 11 =====
% ID: LATEX_0658205A4C5B
% Mức: TH
\begin{ex}
% ID: LATEX_0658205A4C5B
% Mức: TH
Cho hình bình hành $\$ ABCD $\. Rút gọn biểu thức\\overrightarrow{AB} - \overrightarrow{DA} - \overrightarrow{DC}\$ ta được vectơ nào?
\choice
{$\\overrightarrow{AC}\$}
{\True $\\overrightarrow{0}\$}
{$\\overrightarrow{BD}\$}
{$\$ 2 $\overrightarrow{AB}\$}
\loigiai{
• Vì $\$ ABCD $\là hình bình hành nên\\overrightarrow{DC} = \overrightarrow{AB}\$.
 
• Biểu thức trở thành: $\\overrightarrow{AB}$ - $\overrightarrow{DA}$ - $\overrightarrow{AB}$ = - $\overrightarrow{DA}$ = $\overrightarrow{AD}\$.
 
• Tuy nhiên, ta có thể làm cách khác: $\$ \overrightarrow{AB} - \overrightarrow{DA} - \overrightarrow{DC} = \overrightarrow{AB} + \overrightarrow{AD} - \overrightarrow{DC}\ $.
 
• Theo quy tắc hình bình hành$overrightarrow{AB}$+$\overrightarrow{AD}$=$ \overrightarrow{AC}\. Mà
overrightarrow{DC} = \overrightarrow{AB}\ $không giúp ra$ overrightarrow{0}\ $.
 
• Hãy thử lại:$\$\overrightarrow{AB} - \overrightarrow{DA} - \overrightarrow{DC} = \overrightarrow{AB} + \overrightarrow{AD} - \overrightarrow{DC} = \overrightarrow{AC} - \overrightarrow{DC} = \overrightarrow{AC} + \overrightarrow{CD} = \overrightarrow{AD}\$.
 
• Xin lỗi, có sự nhầm lẫn. Hãy làm lại: $\\overrightarrow{AB}$ - $\overrightarrow{DA}$ - $\overrightarrow{DC}$ = $\overrightarrow{AB}$ + $\overrightarrow{AD}$ - $\overrightarrow{DC}\. Mà\,\text{ABCD}\$ là hình bình hành nên $\\overrightarrow{AB}$ = $\overrightarrow{DC}\, do đó\\overrightarrow{AB} - \overrightarrow{DC} = \overrightarrow{0}\$.
 
• Vậy biểu thức còn lại là $\$ - $\overrightarrow{DA}$ = $\overrightarrow{AD}\$.
 
• Phát hiện phương án không khớp, sửa lại câu hỏi: Rút gọn $\$ \overrightarrow{AB} + \overrightarrow{AD} - \overrightarrow{AC}\.
}
\end{ex}

% ===== Câu 12 =====
% ID: LATEX_A9B2A6742B7D
% Mức: TH
\begin{ex}
% ID: LATEX_A9B2A6742B7D
% Mức: TH
Cho hình bình hành $\$ ABCD $\. Biểu thức\\overrightarrow{AB} + \overrightarrow{AD} - \overrightarrow{AC}\$ bằng vectơ nào sau đây?
\choice
{$\\overrightarrow{BD}\$}
{$\\overrightarrow{CB}\$}
{\True $\\overrightarrow{0}\$}
{$\\overrightarrow{DA}\$}
\loigiai{
• Áp dụng quy tắc hình bình hành, ta có tổng $\\overrightarrow{AB}$ + $\overrightarrow{AD}$ = $\overrightarrow{AC}\$.
 
• Thay vào biểu thức ta được: $\$ \overrightarrow{AC} - \overrightarrow{AC} = \overrightarrow{0}\.
}
\end{ex}

% ===== Câu 13 =====
% ID: LATEX_4FD940D16699
% Mức: TH
\begin{ex}
% ID: LATEX_4FD940D16699
% Mức: TH
Cho $\6\điểm\A, B, C, D, E, F\$ bất kì. Đẳng thức nào sau đây đúng?
\choice
{$\\overrightarrow{AB}$ + $\overrightarrow{CD}$ = $\overrightarrow{AD}$ + $\overrightarrow{BC}\$}
{$\\overrightarrow{AD}$ + $\overrightarrow{BE}$ + $\overrightarrow{CF}$ = $\overrightarrow{AE}$ + $\overrightarrow{BF}$ + $\overrightarrow{CE}\$}
{\True $\\overrightarrow{AB}$ - $\overrightarrow{AF}$ + $\overrightarrow{CD}$ - $\overrightarrow{CB}$ + $\overrightarrow{EF}$ - $\overrightarrow{ED}$ = $\overrightarrow{0}\$}
{$\\overrightarrow{AB}$ + $\overrightarrow{CD}$ + $\overrightarrow{EF}$ = $\overrightarrow{AF}$ + $\overrightarrow{ED}$ + $\overrightarrow{BC}\$}
\loigiai{
• Xét vế trái của phương án đúng: $\VT =(\overrightarrow{AB} - \overrightarrow{AF})$ + $(\overrightarrow{CD} - \overrightarrow{CB})$ + $(\overrightarrow{EF} - \overrightarrow{ED})$ \ $.
 
• Áp dụng quy tắc trừ vectơ:$ \ $VT = \overrightarrow{FB} + \overrightarrow{BD} + \overrightarrow{DF}\$.
 
• Áp dụng quy tắc cộng vectơ liên tiếp: $\$ VT = $\overrightarrow{FD}$ + $\overrightarrow{DF}$ = $\overrightarrow{FF}$ = $\overrightarrow{0}\.$
}
\end{ex}

% ===== Câu 14 =====
% ID: LATEX_495D545B9A85
% Mức: TH
\dangbt{Tính độ dài của tổng và hiệu hai vectơ}
\begin{ex}
% ID: LATEX_495D545B9A85
% Mức: TH
Cho tam giác $\,\text{ABC}\đều cạnh\a\$. Tính độ dài của vectơ $\\overrightarrow{AB}$ - $\overrightarrow{AC}\$.
\choice
{\True $\$ a $\$}
{$\$ a $\sqrt{3}\$}
{$\2\,\text{a}$ \}
{$\$ 0 $\$}
\loigiai{
• Áp dụng quy tắc hiệu hai vectơ chung gốc, ta có: $\\overrightarrow{AB}$ - $\overrightarrow{AC}$ = $\overrightarrow{CB}\$.
 
• Độ dài của vectơ này là: $\$ |\overrightarrow{AB} - \overrightarrow{AC}| = |\overrightarrow{CB}| = CB\ $.
 
• Vì tam giác$\$ABC$\đều cạnh\a\$nên$\$CB = a$ \.
}
\end{ex}

% ===== Câu 15 =====
% ID: LATEX_3704ABD1EBCC
% Mức: TH
\begin{ex}
% ID: LATEX_3704ABD1EBCC
% Mức: TH
Cho tam giác $\,\text{ABC}\đều cạnh\a\$. Tính độ dài của vectơ $\\overrightarrow{AB}$ + $\overrightarrow{AC}\$.
\choice
{$\$ a $\$}
{\True $\$ a $\sqrt{3}\$}
{$\2\,\text{a}$ \}
{$\$ a $\sqrt{2}\$}
\loigiai{
• Gọi $\,\text{M}\là trung điểm của cạnh\,\text{BC}\$. Theo tính chất trung điểm, ta có $\\overrightarrow{AB}$ + $\overrightarrow{AC}$ = 2 $\overrightarrow{AM}\$.
 
• Do đó, $\$ |\overrightarrow{AB} + \overrightarrow{AC}| = |2\overrightarrow{AM}| = 2AM\ $.
 
• Trong tam giác đều cạnh$\$a$\, đường trung tuyến đồng thời là đường cao nên\,\text{AM} = \dfrac{a\sqrt{3}}{2}\$.
 
• Vậy$\$|$\overrightarrow{AB}$+$\overrightarrow{AC}| = 2\cdot\dfrac{a\sqrt{3}}{2}$= a$ \sqrt{3}\.
}
\end{ex}

% ===== Câu 16 =====
% ID: LATEX_BD8F813D04FB
% Mức: TH
\begin{ex}
% ID: LATEX_BD8F813D04FB
% Mức: TH
Cho hình vuông $\$ ABCD $\có cạnh bằng\a\$. Tính độ dài của vectơ $\\overrightarrow{AB}$ + $\overrightarrow{AD}\$.
\choice
{$\2\,\text{a}$ \}
{$\$ a $\$}
{\True $\$ a $\sqrt{2}\$}
{$\\dfrac{a\sqrt{2}}{2}\$}
\loigiai{
• Vì $\$ ABCD $\là hình vuông nên nó cũng là hình bình hành. Áp dụng quy tắc hình bình hành ta có:\\overrightarrow{AB} + \overrightarrow{AD} = \overrightarrow{AC}\$.
 
• Độ dài vectơ cần tìm là $\$ | $\overrightarrow{AB}$ + $\overrightarrow{AD}| = |\overrightarrow{AC}| = AC\$.
 
• Trong hình vuông cạnh $\$ a\ $, độ dài đường chéo là$ \ $AC =$ \sqrt{AB^2 + BC^2} $=$ \sqrt{a^2 + a^2} $= a$ \sqrt{2}\.
}
\end{ex}

% ===== Câu 17 =====
% ID: LATEX_331847432B03
% Mức: TH
\begin{ex}
% ID: LATEX_331847432B03
% Mức: TH
Cho hình vuông $\$ ABCD $\tâm\,\text{O}\$, cạnh $\$ a $\. Tính độ dài của vectơ\\overrightarrow{OA} - \overrightarrow{CB}\$.
\choice
{\True $\\dfrac{a\sqrt{2}}{2}\$}
{$\$ a $\sqrt{2}\$}
{$\$ a $\$}
{$\\dfrac{a}{2}\$}
\loigiai{
• Vì $\$ ABCD $\là hình vuông nên\\overrightarrow{CB} = \overrightarrow{DA}\$.
 
• Thay vào biểu thức ta có: $\\overrightarrow{OA}$ - $\overrightarrow{CB}$ = $\overrightarrow{OA}$ - $\overrightarrow{DA}$ = $\overrightarrow{OA}$ + $\overrightarrow{AD}\$.
 
• Áp dụng quy tắc ba điểm: $\$ \overrightarrow{OA} + \overrightarrow{AD} = \overrightarrow{OD}\ $.
 
• Độ dài vectơ là$\$|$\overrightarrow{OD}| = OD\. Vì\,\text{O}\$là tâm hình vuông nên$\$OD =$\dfrac{1}{2}BD$=$ \dfrac{a\sqrt{2}}{2}\.
}
\end{ex}

% ===== Câu 18 =====
% ID: LATEX_9E693EBEFA4D
% Mức: TH
\begin{ex}
% ID: LATEX_9E693EBEFA4D
% Mức: TH
Cho hình chữ nhật $\$ ABCD $\có\AB = 3\$, $\$ BC = 4 $\. Tính độ dài của vectơ\\overrightarrow{AB} + \overrightarrow{BC}\$.
\choice
{$\$ 7 $\$}
{\True $\$ 5 $\$}
{$\$ 1 $\$}
{$\$ 12 $\$}
\loigiai{
• Theo quy tắc ba điểm, ta có $\\overrightarrow{AB}$ + $\overrightarrow{BC}$ = $\overrightarrow{AC}\$.
 
• Độ dài của vectơ là $\$ |\overrightarrow{AC}| = AC\ $.
 
• Xét tam giác$\$ABC$\vuông tại\,\text{B}\$, theo định lí Pythagore:$\$AC =$\sqrt{AB^2 + BC^2}$=$\sqrt{3^2 + 4^2}$=$\sqrt{25}$= 5$ \.
}
\end{ex}

% ===== Câu 19 =====
% ID: LATEX_6F0ACD0D2B2C
% Mức: TH
\dangbt{Xác định điểm thỏa mãn đẳng thức vectơ}
\begin{ex}
% ID: LATEX_6F0ACD0D2B2C
% Mức: TH
Cho đoạn thẳng $\,\text{AB}\và điểm\,\text{M}\$ thỏa mãn đẳng thức $\\overrightarrow{MA}$ + $\overrightarrow{MB}$ = $\overrightarrow{0}\. Khẳng định nào sau đây đúng?$
\choice
{\True $\M\là trung điểm của đoạn thẳng\AB\$}
{$\M\trùng với điểm\A\$}
{$\M\trùng với điểm\B\$}
{$\M\nằm ngoài đoạn thẳng\AB\$}
\loigiai{
• Ta có đẳng thức $\\overrightarrow{MA}$ + $\overrightarrow{MB}$ = $\overrightarrow{0}\Leftrightarrow\overrightarrow{MA}$ = - $\overrightarrow{MB}\$.
 
• Điều này có nghĩa là hai vectơ $\$ \overrightarrow{MA}\ $và$ \ $\overrightarrow{MB}\$ là hai vectơ đối nhau, tức là chúng ngược hướng và có độ dài bằng nhau ($\$ MA = MB\ $).
 
• Do đó, ba điểm$\,\text{M}$,$A$,$B$\thẳng hàng,\M\$nằm giữa$\A$,$B$\và cách đều\,\text{A}, B\$. Vậy$\,\text{M}$\là trung điểm của\AB\$.
}
\end{ex}

% ===== Câu 20 =====
% ID: LATEX_40087EE97ADD
% Mức: TH
\begin{ex}
% ID: LATEX_40087EE97ADD
% Mức: TH
Cho ba điểm $\,\text{A}, B, C\phân biệt. Tìm điểm\,\text{M}\$ thỏa mãn $\\overrightarrow{MA}$ - $\overrightarrow{MB}$ + $\overrightarrow{MC}$ = $\overrightarrow{0}\$.
\choice
{$\M\là trung điểm của\AC\$}
{\True $\M\là đỉnh thứ tư của hình bình hành\ABCM\$}
{$\M\là đỉnh thứ tư của hình bình hành\ABMC\$}
{$\M\là trọng tâm tam giác\ABC\$}
\loigiai{
• Áp dụng quy tắc hiệu cho hai vectơ đầu tiên: $\\overrightarrow{MA}$ - $\overrightarrow{MB}$ = $\overrightarrow{BA}\$.
 
• Đẳng thức đã cho trở thành: $\$ \overrightarrow{BA} + \overrightarrow{MC} = \overrightarrow{0} \Leftrightarrow \overrightarrow{MC} = -\overrightarrow{BA} \Leftrightarrow \overrightarrow{MC} = \overrightarrow{AB}\ $.
 
• Điều kiện$overrightarrow{MC}$=$\overrightarrow{AB}\chứng tỏ tứ giác\,\text{ABCM}\$là hình bình hành (với thứ tự các đỉnh liên tiếp là$\,\text{A}$,$B$,$C$,$M$\).$
}
\end{ex}

% ===== Câu 21 =====
% ID: LATEX_48C88D4FB932
% Mức: TH
\begin{ex}
% ID: LATEX_48C88D4FB932
% Mức: TH
Cho tam giác $\,\text{ABC}\. Gọi\,\text{M}\$ là điểm xác định bởi đẳng thức $\\overrightarrow{AM}$ = $\overrightarrow{AB}$ + $\overrightarrow{AC}\. Khẳng định nào sau đây là đúng?$
\choice
{$\M\là trung điểm của\BC\$}
{$\M\là trọng tâm tam giác\ABC\$}
{\True $\$ ABMC $\là hình bình hành$}
{$\$ AMBC $\là hình bình hành$}
\loigiai{
• Theo quy tắc hình bình hành, tổng hai vectơ chung gốc $\\overrightarrow{AB} + \overrightarrow{AC}\bằng vectơ đường chéo\\overrightarrow{AM}\$ của hình bình hành có ba đỉnh là $\,\text{A}$, $B$, $C$ \ $.
 
• Do đó, tứ giác tạo bởi hai cạnh$\$AB\$và$\$AC$\$cùng với đường chéo$\$AM\$là hình bình hành$\$ABMC$\$.
 
• Chú ý thứ tự đỉnh: \$A\$nối với$\$B$\$và$\$C\$,$\$M$\$đối diện$\$A\$, nên tứ giác đọc theo vòng tròn là$\$ABMC$ \.
}
\end{ex}

% ===== Câu 22 =====
% ID: LATEX_E0B4D02AEB3B
% Mức: TH
\begin{ex}
% ID: LATEX_E0B4D02AEB3B
% Mức: TH
Cho bốn điểm phân biệt $\,\text{A}, B, C, D\. Gọi\,\text{M}\$ là điểm thỏa mãn $\\overrightarrow{AM}$ = $\overrightarrow{AB}$ - $\overrightarrow{AC}\. Khẳng định nào sau đây đúng?$
\choice
{\True $\\overrightarrow{AM}=\overrightarrow{CB}\$}
{$\\overrightarrow{AM}=\overrightarrow{BC}\$}
{$\\overrightarrow{AM}=\overrightarrow{CD}\$}
{$\\overrightarrow{AM}=\overrightarrow{AD}\$}
\loigiai{
• Áp dụng quy tắc trừ hai vectơ chung gốc: $\\overrightarrow{AB}$ - $\overrightarrow{AC}$ = $\overrightarrow{CB}\$.
 
• Vậy đẳng thức đã cho trở thành $\$ \overrightarrow{AM} = \overrightarrow{CB}\.
}
\end{ex}

% ===== Câu 23 =====
% ID: LATEX_D8CCCEF6C78E
% Mức: TH
\begin{ex}
% ID: LATEX_D8CCCEF6C78E
% Mức: TH
Cho tam giác $\,\text{ABC}\. Tập hợp các điểm\,\text{M}\$ thỏa mãn $\$ | $\overrightarrow{MA}$ + $\overrightarrow{MB}| = |\overrightarrow{MC}$ + $\overrightarrow{MB}|\là:$
\choice
{Đường tròn tâm $\B\bán kính\AC\$}
{Đường tròn đường kính $\$ AC $\$}
{\True Đường trung trực của đoạn thẳng $\$ AC $\$}
{Đường trung trực của đoạn thẳng $\$ BC $\$}
\loigiai{
• Xét vế trái: Gọi $\,\text{I}\là trung điểm của\,\text{AB}\$, ta có $\\overrightarrow{MA}$ + $\overrightarrow{MB}$ = 2 $\overrightarrow{MI}\Rightarrow$ | $\overrightarrow{MA}$ + $\overrightarrow{MB}| = 2MI\$.
 
• Xét vế phải: Gọi $\$ J\ $là trung điểm của$ \ $BC$ \ $, ta có$ \ $\overrightarrow{MC} + \overrightarrow{MB} = 2\overrightarrow{MJ} \Rightarrow |\overrightarrow{MC} + \overrightarrow{MB}| = 2MJ\$.
 
• Đẳng thức trở thành $\$ 2MI = 2MJ \Leftrightarrow MI = MJ $\$.
 
• Do đó, tập hợp các điểm $\$ M\ $là đường trung trực của đoạn thẳng$ \ $IJ$ \ $.
 
• Mặt khác, \$IJ\$là đường trung bình của tam giác$\$ABC$\$nên$\$IJ \parallel AC\$. Đường trung trực của$\$IJ$\$cũng chính là đường trung trực của đoạn$\$AC\$(nếu xét tam giác cân) - khoảng cách. Tuy nhiên, cách giải chuẩn hơn mà không cần gọi trung điểm là:
 
• Cách 2:$\$|$\overrightarrow{MA}$+$\overrightarrow{MB}| = |\overrightarrow{MB}$+$\overrightarrow{MC}|\. Đây là độ dài của tổng, không dễ khử điểm\,\text{M}\$.
 
• Chờ đã,$overrightarrow{MB}$+$\overrightarrow{MC}\chưa triệt tiêu\,\text{M}\$. Phải dùng trung điểm$\,\text{I}$,$J$\. Khi\MI = MJ\$,$\M$\nằm trên trung trực của\,\text{IJ}\$. Đường trung trực của$\$IJ$\không nhất thiết là đường trung trực của\,\text{AC}\$.
 
• Hãy đổi câu hỏi thành một bài toán quen thuộc hơn:$\$|$\overrightarrow{MB}$-$\overrightarrow{MC}| = |\overrightarrow{BM}$-$\overrightarrow{BA}|\$.
 
• Sửa lại câu hỏi: Cho tam giác$\$ABC\$. Điểm$\$M$\$thỏa mãn$\$ \overrightarrow{MA} + \overrightarrow{MB} + \overrightarrow{MC} = \overrightarrow{0}\.
}
\end{ex}

% ===== Câu 24 =====
% ID: LATEX_29094683BBF1
% Mức: TH
\dangbt{Tính chất trung điểm và trọng tâm trong tổng vectơ}
\begin{ex}
% ID: LATEX_29094683BBF1
% Mức: TH
Cho tam giác $\ABC\có trọng tâm\G\$. Khẳng định nào sau đây là đúng?
\choice
{$\\overrightarrow{GA} + \overrightarrow{GB} + \overrightarrow{GC} = 3\overrightarrow{GM}\với\,\text{M}\$ là trung điểm $\$ BC $\$}
{$\\overrightarrow{AB}$ + $\overrightarrow{AC}$ = $\overrightarrow{AG}\$}
{\True $\\overrightarrow{GA}$ + $\overrightarrow{GB}$ + $\overrightarrow{GC}$ = $\overrightarrow{0}\$}
{$\\overrightarrow{GA}$ + $\overrightarrow{GB}$ = $\overrightarrow{GC}\$}
\loigiai{
• Theo tính chất trọng tâm của tam giác, tổng các vectơ hướng từ trọng tâm đến ba đỉnh luôn bằng vectơ-không.
 
• Vậy $\\overrightarrow{GA}$ + $\overrightarrow{GB}$ + $\overrightarrow{GC}$ = $\overrightarrow{0}\.$
}
\end{ex}

% ===== Câu 25 =====
% ID: LATEX_1BA9FFC3EF07
% Mức: TH
\begin{ex}
% ID: LATEX_1BA9FFC3EF07
% Mức: TH
Cho hình bình hành $\$ ABCD $\tâm\O\$. Đẳng thức nào sau đây là đúng?
\choice
{\True $\\overrightarrow{OA}$ + $\overrightarrow{OB}$ + $\overrightarrow{OC}$ + $\overrightarrow{OD}$ = $\overrightarrow{0}\$}
{$\\overrightarrow{OA}$ + $\overrightarrow{OB}$ + $\overrightarrow{OC}$ + $\overrightarrow{OD}$ = 2 $\overrightarrow{AB}\$}
{$\\overrightarrow{OA}$ + $\overrightarrow{OB}$ + $\overrightarrow{OC}$ + $\overrightarrow{OD}$ = $\overrightarrow{AC}\$}
{$\\overrightarrow{OA}$ + $\overrightarrow{OB}$ + $\overrightarrow{OC}$ + $\overrightarrow{OD}$ = 4 $\overrightarrow{O}\$}
\loigiai{
• Vì $\,\text{O}\là tâm hình bình hành\,\text{ABCD}\$ nên $\,\text{O}$ \là trung điểm của đường chéo\AC\ $và đường chéo$ \ $BD$ \ $.
 
• Áp dụng tính chất trung điểm:$\$\overrightarrow{OA} + \overrightarrow{OC} = \overrightarrow{0}\$và$\$\overrightarrow{OB}$+$\overrightarrow{OD}$=$\overrightarrow{0}\$.
 
• Cộng vế theo vế ta được: \$ \overrightarrow{OA} + \overrightarrow{OB} + \overrightarrow{OC} + \overrightarrow{OD} = \overrightarrow{0} + \overrightarrow{0} = \overrightarrow{0}\.
}
\end{ex}

% ===== Câu 26 =====
% ID: LATEX_C2ED8EBD2F28
% Mức: TH
\begin{ex}
% ID: LATEX_C2ED8EBD2F28
% Mức: TH
Cho đoạn thẳng $\,\text{AB}\có trung điểm là\,\text{I}\$. Với điểm $\,\text{M}$ \tùy ý, khẳng định nào sau đây là đúng?
\choice
{$\\overrightarrow{MA}$ + $\overrightarrow{MB}$ = $\overrightarrow{MI}\$}
{\True $\\overrightarrow{MA}$ + $\overrightarrow{MB}$ = 2 $\overrightarrow{MI}\$}
{$\\overrightarrow{MA}$ + $\overrightarrow{MB}$ = 3 $\overrightarrow{MI}\$}
{$\\overrightarrow{MA}$ + $\overrightarrow{MB}$ = $\overrightarrow{IM}\$}
\loigiai{
• Chèn điểm $\,\text{I}\vào hai vectơ ở vế trái:\\overrightarrow{MA} = \overrightarrow{MI} + \overrightarrow{IA}\$ và $\\overrightarrow{MB}$ = $\overrightarrow{MI}$ + $\overrightarrow{IB}\$.
 
• Cộng lại ta có: $\$ \overrightarrow{MA} + \overrightarrow{MB} = 2\overrightarrow{MI} + (\overrightarrow{IA} + \overrightarrow{IB})\ $.
 
• Vì$\,\text{I}$\là trung điểm của\AB\$nên$\\overrightarrow{IA}$+$\overrightarrow{IB}$=$\overrightarrow{0}\$.
 
• Vậy $\$ \overrightarrow{MA} + \overrightarrow{MB} = 2\overrightarrow{MI}\.
}
\end{ex}

% ===== Câu 27 =====
% ID: LATEX_BF46C9E9C4F6
% Mức: TH
\begin{ex}
% ID: LATEX_BF46C9E9C4F6
% Mức: TH
Cho tam giác $\ABC\. Gọi\M, N, P\$ lần lượt là trung điểm của các cạnh $\$ BC, CA, AB $\. Tổng\\overrightarrow{AM} + \overrightarrow{BN} + \overrightarrow{CP}\$ bằng:
\choice
{$\\overrightarrow{AB}\$}
{$\\overrightarrow{BC}\$}
{\True $\\overrightarrow{0}\$}
{$\\overrightarrow{CA}\$}
\loigiai{
• Theo tính chất trung điểm, ta có $\2\overrightarrow{AM} = \overrightarrow{AB} + \overrightarrow{AC}\;\2\overrightarrow{BN} = \overrightarrow{BA} + \overrightarrow{BC}\$; $\$ 2 $\overrightarrow{CP}$ = $\overrightarrow{CA}$ + $\overrightarrow{CB}\$.
 
• Cộng ba đẳng thức trên: $\$ 2(\overrightarrow{AM} + \overrightarrow{BN} + \overrightarrow{CP}) = (\overrightarrow{AB} + \overrightarrow{BA}) + (\overrightarrow{AC} + \overrightarrow{CA}) + (\overrightarrow{BC} + \overrightarrow{CB})\ $.
 
•$\$2$\overrightarrow{AM} + \overrightarrow{BN} + \overrightarrow{CP})$=$\overrightarrow{0}$+$\overrightarrow{0}$+$\overrightarrow{0}$=$\overrightarrow{0}\$.
 
• Vậy$\$\overrightarrow{AM} + \overrightarrow{BN} + \overrightarrow{CP} = \overrightarrow{0}\$.
}
\end{ex}

% ===== Câu 28 =====
% ID: LATEX_AE13D1889B7A
% Mức: TH
\begin{ex}
% ID: LATEX_AE13D1889B7A
% Mức: TH
Cho tam giác $\,\text{ABC}\có trọng tâm\,\text{G}\$. Với điểm $\,\text{M}$ \bất kì, đẳng thức nào sau đây đúng?
\choice
{$\\overrightarrow{MA}$ + $\overrightarrow{MB}$ + $\overrightarrow{MC}$ = $\overrightarrow{MG}\$}
{$\\overrightarrow{MA}$ + $\overrightarrow{MB}$ + $\overrightarrow{MC}$ = 2 $\overrightarrow{MG}\$}
{\True $\\overrightarrow{MA}$ + $\overrightarrow{MB}$ + $\overrightarrow{MC}$ = 3 $\overrightarrow{MG}\$}
{$\\overrightarrow{MA}$ + $\overrightarrow{MB}$ + $\overrightarrow{MC}$ = 4 $\overrightarrow{MG}\$}
\loigiai{
• Chèn điểm $\,\text{G}\vào từng vectơ ở vế trái:\\overrightarrow{MA} = \overrightarrow{MG} + \overrightarrow{GA}\$; $\\overrightarrow{MB}$ = $\overrightarrow{MG}$ + $\overrightarrow{GB}\;\\overrightarrow{MC} = \overrightarrow{MG} + \overrightarrow{GC}\$.
 
• Cộng ba đẳng thức lại: $\\overrightarrow{MA}$ + $\overrightarrow{MB}$ + $\overrightarrow{MC}$ = 3 $\overrightarrow{MG}$ + $\overrightarrow{GA} + \overrightarrow{GB} + \overrightarrow{GC})$ \ $.
 
• Vì$\$G\$là trọng tâm tam giác$\$ABC$\$nên$\$\overrightarrow{GA} + \overrightarrow{GB} + \overrightarrow{GC} = \overrightarrow{0}\$.
 
• Vậy$overrightarrow{MA}$+$\overrightarrow{MB}$+$\overrightarrow{MC}$= 3$ \overrightarrow{MG}\.
}
\end{ex}

% ===== Câu 29 =====
% ID: LATEX_0DB1EA144BE3
% Mức: TH
\dangbt{Áp dụng quy tắc ba điểm và quy tắc hình bình hành}
\begin{ex}
% ID: LATEX_0DB1EA144BE3
% Mức: TH
Cho ba điểm phân biệt $\A, B, C\. Đẳng thức nào sau đây đúng?$
\choice
{$\\overrightarrow{AB}$ + $\overrightarrow{CB}$ = $\overrightarrow{AC}\$}
{$\\overrightarrow{CB}$ + $\overrightarrow{AC}$ = $\overrightarrow{AB}\$}
{$\\overrightarrow{AB}$ + $\overrightarrow{BC}$ = $\overrightarrow{CA}\$}
{\True $\\overrightarrow{AB}$ + $\overrightarrow{BC}$ = $\overrightarrow{AC}\$}
\loigiai{
Theo quy tắc ba điểm đối với phép cộng vectơ, với ba điểm $\A, B, C\bất kì ta luôn có\\overrightarrow{AB} + \overrightarrow{BC} = \overrightarrow{AC}\$.
}
\end{ex}

% ===== Câu 30 =====
% ID: LATEX_AF6B2F66DC9A
% Mức: TH
\begin{ex}
% ID: LATEX_AF6B2F66DC9A
% Mức: TH
Cho ba điểm phân biệt $\M, N, P\. Vectơ tổng\\overrightarrow{NP} + \overrightarrow{MN}\$ bằng vectơ nào sau đây?
\choice
{$\\overrightarrow{PN}\$}
{$\\overrightarrow{PM}\$}
{\True $\\overrightarrow{MP}\$}
{$\\overrightarrow{NM}\$}
\loigiai{
• Áp dụng tính chất giao hoán của phép cộng vectơ, ta có: $\\overrightarrow{NP}$ + $\overrightarrow{MN}$ = $\overrightarrow{MN}$ + $\overrightarrow{NP}\$.
 
• Theo quy tắc ba điểm, $\$ \overrightarrow{MN} + \overrightarrow{NP} = \overrightarrow{MP}\ $.
 
• Vậy$overrightarrow{NP}$+$\overrightarrow{MN}$=$ \overrightarrow{MP}\.
}
\end{ex}

% ===== Câu 31 =====
% ID: LATEX_B3959566D3C6
% Mức: TH
\begin{ex}
% ID: LATEX_B3959566D3C6
% Mức: TH
Cho hình bình hành $\$ ABCD $\. Đẳng thức nào sau đây đúng?$
\choice
{\True $\\overrightarrow{AB}$ + $\overrightarrow{AD}$ = $\overrightarrow{AC}\$}
{$\\overrightarrow{AB}$ + $\overrightarrow{AD}$ = $\overrightarrow{BD}\$}
{$\\overrightarrow{AB}$ + $\overrightarrow{AD}$ = $\overrightarrow{CD}\$}
{$\\overrightarrow{AB}$ + $\overrightarrow{AD}$ = $\overrightarrow{CB}\$}
\loigiai{
Theo quy tắc hình bình hành, vì $\$ ABCD $\là hình bình hành nên tổng hai vectơ cạnh chung gốc bằng vectơ đường chéo xuất phát từ gốc đó, tức là\\overrightarrow{AB} + \overrightarrow{AD} = \overrightarrow{AC}\$.
}
\end{ex}

% ===== Câu 32 =====
% ID: LATEX_1431DB984C9E
% Mức: TH
\begin{ex}
% ID: LATEX_1431DB984C9E
% Mức: TH
Cho bốn điểm phân biệt $\A, B, C, O\. Đẳng thức nào sau đây đúng?$
\choice
{$\\overrightarrow{AB}=\overrightarrow{OB}$ + $\overrightarrow{OA}\$}
{$\\overrightarrow{AB}=\overrightarrow{AC}$ + $\overrightarrow{BC}\$}
{\True $\\overrightarrow{OA}=\overrightarrow{CA}$ - $\overrightarrow{CO}\$}
{$\\overrightarrow{OA}=\overrightarrow{OB}$ - $\overrightarrow{BA}\$}
\loigiai{
• Áp dụng quy tắc trừ vectơ với ba điểm $\O, C, A\ta có:\\overrightarrow{CA} - \overrightarrow{CO} = \overrightarrow{OA}\$.
 
• Các phương án khác đều sai quy tắc cộng và trừ vectơ.
}
\end{ex}

% ===== Câu 33 =====
% ID: LATEX_A46C88793698
% Mức: TH
\begin{ex}
% ID: LATEX_A46C88793698
% Mức: TH
Cho bốn điểm phân biệt $\A, B, C, D\. Vectơ tổng\\overrightarrow{AB} + \overrightarrow{CD} + \overrightarrow{BC} + \overrightarrow{DA}\$ bằng vectơ nào sau đây?
\choice
{\True $\\overrightarrow{0}\$}
{$\\overrightarrow{AC}\$}
{$\\overrightarrow{BD}\$}
{$\\overrightarrow{BA}\$}
\loigiai{
• Ta sắp xếp lại các số hạng và áp dụng quy tắc ba điểm liên tiếp: $\\overrightarrow{AB}$ + $\overrightarrow{CD}$ + $\overrightarrow{BC}$ + $\overrightarrow{DA}$ = $\overrightarrow{AB}$ + $\overrightarrow{BC}$ + $\overrightarrow{CD}$ + $\overrightarrow{DA}\$.
 
• Thực hiện phép cộng từ trái sang phải: $\$ \overrightarrow{AC} + \overrightarrow{CD} + \overrightarrow{DA} = \overrightarrow{AD} + \overrightarrow{DA} = \overrightarrow{AA}\ $.
 
• Vì điểm đầu và điểm cuối trùng nhau nên$overrightarrow{AA}$=$ \overrightarrow{0}\.
}
\end{ex}

% ===== Câu 34 =====
% ID: LATEX_4302F8D17026
% Mức: TH
\dangbt{Rút gọn biểu thức và chứng minh đẳng thức vectơ}
\begin{ex}
% ID: LATEX_4302F8D17026
% Mức: TH
Cho bốn điểm $\A, B, C, D\bất kì. Khẳng định nào sau đây là đúng?$
\choice
{$\\overrightarrow{AB}$ - $\overrightarrow{AD}$ = $\overrightarrow{CD}$ + $\overrightarrow{BC}\$}
{$\\overrightarrow{AB}$ - $\overrightarrow{AC}$ = $\overrightarrow{DB}$ - $\overrightarrow{DC}\$}
{$\\overrightarrow{BC}$ + $\overrightarrow{AB}$ = $\overrightarrow{DA}$ - $\overrightarrow{DC}\$}
{\True $\\overrightarrow{AB}$ + $\overrightarrow{CD}$ = $\overrightarrow{AD}$ + $\overrightarrow{CB}\$}
\loigiai{
• Xét vế trái: $\\overrightarrow{AB} + \overrightarrow{CD}\. Áp dụng quy tắc chèn điểm\,\text{D}\$ vào vectơ thứ nhất và điểm $\B$ \ $vào vectơ thứ hai.
 
• Ta có$ \ $\overrightarrow{AB} + \overrightarrow{CD} = (\overrightarrow{AD} + \overrightarrow{DB}) + (\overrightarrow{CB} + \overrightarrow{BD})\$.
 
• Nhóm lại ta được: $\\overrightarrow{AD}$ + $\overrightarrow{CB}$ + $(\overrightarrow{DB} + \overrightarrow{BD})$ = $\overrightarrow{AD}$ + $\overrightarrow{CB}$ + $\overrightarrow{0}$ = $\overrightarrow{AD}$ + $\overrightarrow{CB}\. Đẳng thức đúng.$
}
\end{ex}

% ===== Câu 35 =====
% ID: LATEX_3BCD556B4A93
% Mức: TH
\begin{ex}
% ID: LATEX_3BCD556B4A93
% Mức: TH
Rút gọn biểu thức vectơ $\\overrightarrow{P} = \overrightarrow{AB} - \overrightarrow{AC} + \overrightarrow{BD}\, ta được kết quả nào sau đây?$
\choice
{$\\overrightarrow{AD}\$}
{$\\overrightarrow{BC}\$}
{\True $\\overrightarrow{CD}\$}
{$\\overrightarrow{DA}\$}
\loigiai{
• Áp dụng quy tắc hiệu cho hai vectơ đầu tiên: $\\overrightarrow{AB}$ - $\overrightarrow{AC}$ = $\overrightarrow{CB}\$.
 
• Thay vào biểu thức ta được: $\$ \overrightarrow{P} = \overrightarrow{CB} + \overrightarrow{BD}\ $.
 
• Áp dụng quy tắc ba điểm:$overrightarrow{CB}$+$\overrightarrow{BD}$=$ \overrightarrow{CD}\.
}
\end{ex}

% ===== Câu 36 =====
% ID: LATEX_0658205A4C5B
% Mức: TH
\begin{ex}
% ID: LATEX_0658205A4C5B
% Mức: TH
Cho hình bình hành $\$ ABCD $\. Rút gọn biểu thức\\overrightarrow{AB} - \overrightarrow{DA} - \overrightarrow{DC}\$ ta được vectơ nào?
\choice
{$\\overrightarrow{AC}\$}
{\True $\\overrightarrow{0}\$}
{$\\overrightarrow{BD}\$}
{$\$ 2 $\overrightarrow{AB}\$}
\loigiai{
• Vì $\$ ABCD $\là hình bình hành nên\\overrightarrow{DC} = \overrightarrow{AB}\$.
 
• Biểu thức trở thành: $\\overrightarrow{AB}$ - $\overrightarrow{DA}$ - $\overrightarrow{AB}$ = - $\overrightarrow{DA}$ = $\overrightarrow{AD}\$.
 
• Tuy nhiên, ta có thể làm cách khác: $\$ \overrightarrow{AB} - \overrightarrow{DA} - \overrightarrow{DC} = \overrightarrow{AB} + \overrightarrow{AD} - \overrightarrow{DC}\ $.
 
• Theo quy tắc hình bình hành$overrightarrow{AB}$+$\overrightarrow{AD}$=$ \overrightarrow{AC}\. Mà
overrightarrow{DC} = \overrightarrow{AB}\ $không giúp ra$ overrightarrow{0}\ $.
 
• Hãy thử lại:$\$\overrightarrow{AB} - \overrightarrow{DA} - \overrightarrow{DC} = \overrightarrow{AB} + \overrightarrow{AD} - \overrightarrow{DC} = \overrightarrow{AC} - \overrightarrow{DC} = \overrightarrow{AC} + \overrightarrow{CD} = \overrightarrow{AD}\$.
 
• Xin lỗi, có sự nhầm lẫn. Hãy làm lại: $\\overrightarrow{AB}$ - $\overrightarrow{DA}$ - $\overrightarrow{DC}$ = $\overrightarrow{AB}$ + $\overrightarrow{AD}$ - $\overrightarrow{DC}\. Mà\,\text{ABCD}\$ là hình bình hành nên $\\overrightarrow{AB}$ = $\overrightarrow{DC}\, do đó\\overrightarrow{AB} - \overrightarrow{DC} = \overrightarrow{0}\$.
 
• Vậy biểu thức còn lại là $\$ - $\overrightarrow{DA}$ = $\overrightarrow{AD}\$.
 
• Phát hiện phương án không khớp, sửa lại câu hỏi: Rút gọn $\$ \overrightarrow{AB} + \overrightarrow{AD} - \overrightarrow{AC}\.
}
\end{ex}

% ===== Câu 37 =====
% ID: LATEX_A9B2A6742B7D
% Mức: TH
\begin{ex}
% ID: LATEX_A9B2A6742B7D
% Mức: TH
Cho hình bình hành $\$ ABCD $\. Biểu thức\\overrightarrow{AB} + \overrightarrow{AD} - \overrightarrow{AC}\$ bằng vectơ nào sau đây?
\choice
{$\\overrightarrow{BD}\$}
{$\\overrightarrow{CB}\$}
{\True $\\overrightarrow{0}\$}
{$\\overrightarrow{DA}\$}
\loigiai{
• Áp dụng quy tắc hình bình hành, ta có tổng $\\overrightarrow{AB}$ + $\overrightarrow{AD}$ = $\overrightarrow{AC}\$.
 
• Thay vào biểu thức ta được: $\$ \overrightarrow{AC} - \overrightarrow{AC} = \overrightarrow{0}\.
}
\end{ex}

% ===== Câu 38 =====
% ID: LATEX_4FD940D16699
% Mức: TH
\begin{ex}
% ID: LATEX_4FD940D16699
% Mức: TH
Cho $\6\điểm\A, B, C, D, E, F\$ bất kì. Đẳng thức nào sau đây đúng?
\choice
{$\\overrightarrow{AB}$ + $\overrightarrow{CD}$ = $\overrightarrow{AD}$ + $\overrightarrow{BC}\$}
{$\\overrightarrow{AD}$ + $\overrightarrow{BE}$ + $\overrightarrow{CF}$ = $\overrightarrow{AE}$ + $\overrightarrow{BF}$ + $\overrightarrow{CE}\$}
{\True $\\overrightarrow{AB}$ - $\overrightarrow{AF}$ + $\overrightarrow{CD}$ - $\overrightarrow{CB}$ + $\overrightarrow{EF}$ - $\overrightarrow{ED}$ = $\overrightarrow{0}\$}
{$\\overrightarrow{AB}$ + $\overrightarrow{CD}$ + $\overrightarrow{EF}$ = $\overrightarrow{AF}$ + $\overrightarrow{ED}$ + $\overrightarrow{BC}\$}
\loigiai{
• Xét vế trái của phương án đúng: $\VT =(\overrightarrow{AB} - \overrightarrow{AF})$ + $(\overrightarrow{CD} - \overrightarrow{CB})$ + $(\overrightarrow{EF} - \overrightarrow{ED})$ \ $.
 
• Áp dụng quy tắc trừ vectơ:$ \ $VT = \overrightarrow{FB} + \overrightarrow{BD} + \overrightarrow{DF}\$.
 
• Áp dụng quy tắc cộng vectơ liên tiếp: $\$ VT = $\overrightarrow{FD}$ + $\overrightarrow{DF}$ = $\overrightarrow{FF}$ = $\overrightarrow{0}\.$
}
\end{ex}

% ===== Câu 39 =====
% ID: LATEX_495D545B9A85
% Mức: TH
\dangbt{Tính độ dài của tổng và hiệu hai vectơ}
\begin{ex}
% ID: LATEX_495D545B9A85
% Mức: TH
Cho tam giác $\,\text{ABC}\đều cạnh\a\$. Tính độ dài của vectơ $\\overrightarrow{AB}$ - $\overrightarrow{AC}\$.
\choice
{\True $\$ a $\$}
{$\$ a $\sqrt{3}\$}
{$\2\,\text{a}$ \}
{$\$ 0 $\$}
\loigiai{
• Áp dụng quy tắc hiệu hai vectơ chung gốc, ta có: $\\overrightarrow{AB}$ - $\overrightarrow{AC}$ = $\overrightarrow{CB}\$.
 
• Độ dài của vectơ này là: $\$ |\overrightarrow{AB} - \overrightarrow{AC}| = |\overrightarrow{CB}| = CB\ $.
 
• Vì tam giác$\$ABC$\đều cạnh\a\$nên$\$CB = a$ \.
}
\end{ex}

% ===== Câu 40 =====
% ID: LATEX_3704ABD1EBCC
% Mức: TH
\begin{ex}
% ID: LATEX_3704ABD1EBCC
% Mức: TH
Cho tam giác $\,\text{ABC}\đều cạnh\a\$. Tính độ dài của vectơ $\\overrightarrow{AB}$ + $\overrightarrow{AC}\$.
\choice
{$\$ a $\$}
{\True $\$ a $\sqrt{3}\$}
{$\2\,\text{a}$ \}
{$\$ a $\sqrt{2}\$}
\loigiai{
• Gọi $\,\text{M}\là trung điểm của cạnh\,\text{BC}\$. Theo tính chất trung điểm, ta có $\\overrightarrow{AB}$ + $\overrightarrow{AC}$ = 2 $\overrightarrow{AM}\$.
 
• Do đó, $\$ |\overrightarrow{AB} + \overrightarrow{AC}| = |2\overrightarrow{AM}| = 2AM\ $.
 
• Trong tam giác đều cạnh$\$a$\, đường trung tuyến đồng thời là đường cao nên\,\text{AM} = \dfrac{a\sqrt{3}}{2}\$.
 
• Vậy$\$|$\overrightarrow{AB}$+$\overrightarrow{AC}| = 2\cdot\dfrac{a\sqrt{3}}{2}$= a$ \sqrt{3}\.
}
\end{ex}

% ===== Câu 41 =====
% ID: LATEX_BD8F813D04FB
% Mức: TH
\begin{ex}
% ID: LATEX_BD8F813D04FB
% Mức: TH
Cho hình vuông $\$ ABCD $\có cạnh bằng\a\$. Tính độ dài của vectơ $\\overrightarrow{AB}$ + $\overrightarrow{AD}\$.
\choice
{$\2\,\text{a}$ \}
{$\$ a $\$}
{\True $\$ a $\sqrt{2}\$}
{$\\dfrac{a\sqrt{2}}{2}\$}
\loigiai{
• Vì $\$ ABCD $\là hình vuông nên nó cũng là hình bình hành. Áp dụng quy tắc hình bình hành ta có:\\overrightarrow{AB} + \overrightarrow{AD} = \overrightarrow{AC}\$.
 
• Độ dài vectơ cần tìm là $\$ | $\overrightarrow{AB}$ + $\overrightarrow{AD}| = |\overrightarrow{AC}| = AC\$.
 
• Trong hình vuông cạnh $\$ a\ $, độ dài đường chéo là$ \ $AC =$ \sqrt{AB^2 + BC^2} $=$ \sqrt{a^2 + a^2} $= a$ \sqrt{2}\.
}
\end{ex}

% ===== Câu 42 =====
% ID: LATEX_331847432B03
% Mức: TH
\begin{ex}
% ID: LATEX_331847432B03
% Mức: TH
Cho hình vuông $\$ ABCD $\tâm\,\text{O}\$, cạnh $\$ a $\. Tính độ dài của vectơ\\overrightarrow{OA} - \overrightarrow{CB}\$.
\choice
{\True $\\dfrac{a\sqrt{2}}{2}\$}
{$\$ a $\sqrt{2}\$}
{$\$ a $\$}
{$\\dfrac{a}{2}\$}
\loigiai{
• Vì $\$ ABCD $\là hình vuông nên\\overrightarrow{CB} = \overrightarrow{DA}\$.
 
• Thay vào biểu thức ta có: $\\overrightarrow{OA}$ - $\overrightarrow{CB}$ = $\overrightarrow{OA}$ - $\overrightarrow{DA}$ = $\overrightarrow{OA}$ + $\overrightarrow{AD}\$.
 
• Áp dụng quy tắc ba điểm: $\$ \overrightarrow{OA} + \overrightarrow{AD} = \overrightarrow{OD}\ $.
 
• Độ dài vectơ là$\$|$\overrightarrow{OD}| = OD\. Vì\,\text{O}\$là tâm hình vuông nên$\$OD =$\dfrac{1}{2}BD$=$ \dfrac{a\sqrt{2}}{2}\.
}
\end{ex}

% ===== Câu 43 =====
% ID: LATEX_9E693EBEFA4D
% Mức: TH
\begin{ex}
% ID: LATEX_9E693EBEFA4D
% Mức: TH
Cho hình chữ nhật $\$ ABCD $\có\AB = 3\$, $\$ BC = 4 $\. Tính độ dài của vectơ\\overrightarrow{AB} + \overrightarrow{BC}\$.
\choice
{$\$ 7 $\$}
{\True $\$ 5 $\$}
{$\$ 1 $\$}
{$\$ 12 $\$}
\loigiai{
• Theo quy tắc ba điểm, ta có $\\overrightarrow{AB}$ + $\overrightarrow{BC}$ = $\overrightarrow{AC}\$.
 
• Độ dài của vectơ là $\$ |\overrightarrow{AC}| = AC\ $.
 
• Xét tam giác$\$ABC$\vuông tại\,\text{B}\$, theo định lí Pythagore:$\$AC =$\sqrt{AB^2 + BC^2}$=$\sqrt{3^2 + 4^2}$=$\sqrt{25}$= 5$ \.
}
\end{ex}

% ===== Câu 44 =====
% ID: LATEX_6F0ACD0D2B2C
% Mức: TH
\dangbt{Xác định điểm thỏa mãn đẳng thức vectơ}
\begin{ex}
% ID: LATEX_6F0ACD0D2B2C
% Mức: TH
Cho đoạn thẳng $\,\text{AB}\và điểm\,\text{M}\$ thỏa mãn đẳng thức $\\overrightarrow{MA}$ + $\overrightarrow{MB}$ = $\overrightarrow{0}\. Khẳng định nào sau đây đúng?$
\choice
{\True $\M\là trung điểm của đoạn thẳng\AB\$}
{$\M\trùng với điểm\A\$}
{$\M\trùng với điểm\B\$}
{$\M\nằm ngoài đoạn thẳng\AB\$}
\loigiai{
• Ta có đẳng thức $\\overrightarrow{MA}$ + $\overrightarrow{MB}$ = $\overrightarrow{0}\Leftrightarrow\overrightarrow{MA}$ = - $\overrightarrow{MB}\$.
 
• Điều này có nghĩa là hai vectơ $\$ \overrightarrow{MA}\ $và$ \ $\overrightarrow{MB}\$ là hai vectơ đối nhau, tức là chúng ngược hướng và có độ dài bằng nhau ($\$ MA = MB\ $).
 
• Do đó, ba điểm$\,\text{M}$,$A$,$B$\thẳng hàng,\M\$nằm giữa$\A$,$B$\và cách đều\,\text{A}, B\$. Vậy$\,\text{M}$\là trung điểm của\AB\$.
}
\end{ex}

% ===== Câu 45 =====
% ID: LATEX_40087EE97ADD
% Mức: TH
\begin{ex}
% ID: LATEX_40087EE97ADD
% Mức: TH
Cho ba điểm $\,\text{A}, B, C\phân biệt. Tìm điểm\,\text{M}\$ thỏa mãn $\\overrightarrow{MA}$ - $\overrightarrow{MB}$ + $\overrightarrow{MC}$ = $\overrightarrow{0}\$.
\choice
{$\M\là trung điểm của\AC\$}
{\True $\M\là đỉnh thứ tư của hình bình hành\ABCM\$}
{$\M\là đỉnh thứ tư của hình bình hành\ABMC\$}
{$\M\là trọng tâm tam giác\ABC\$}
\loigiai{
• Áp dụng quy tắc hiệu cho hai vectơ đầu tiên: $\\overrightarrow{MA}$ - $\overrightarrow{MB}$ = $\overrightarrow{BA}\$.
 
• Đẳng thức đã cho trở thành: $\$ \overrightarrow{BA} + \overrightarrow{MC} = \overrightarrow{0} \Leftrightarrow \overrightarrow{MC} = -\overrightarrow{BA} \Leftrightarrow \overrightarrow{MC} = \overrightarrow{AB}\ $.
 
• Điều kiện$overrightarrow{MC}$=$\overrightarrow{AB}\chứng tỏ tứ giác\,\text{ABCM}\$là hình bình hành (với thứ tự các đỉnh liên tiếp là$\,\text{A}$,$B$,$C$,$M$\).$
}
\end{ex}

% ===== Câu 46 =====
% ID: LATEX_48C88D4FB932
% Mức: TH
\begin{ex}
% ID: LATEX_48C88D4FB932
% Mức: TH
Cho tam giác $\,\text{ABC}\. Gọi\,\text{M}\$ là điểm xác định bởi đẳng thức $\\overrightarrow{AM}$ = $\overrightarrow{AB}$ + $\overrightarrow{AC}\. Khẳng định nào sau đây là đúng?$
\choice
{$\M\là trung điểm của\BC\$}
{$\M\là trọng tâm tam giác\ABC\$}
{\True $\$ ABMC $\là hình bình hành$}
{$\$ AMBC $\là hình bình hành$}
\loigiai{
• Theo quy tắc hình bình hành, tổng hai vectơ chung gốc $\\overrightarrow{AB} + \overrightarrow{AC}\bằng vectơ đường chéo\\overrightarrow{AM}\$ của hình bình hành có ba đỉnh là $\,\text{A}$, $B$, $C$ \ $.
 
• Do đó, tứ giác tạo bởi hai cạnh$\$AB\$và$\$AC$\$cùng với đường chéo$\$AM\$là hình bình hành$\$ABMC$\$.
 
• Chú ý thứ tự đỉnh: \$A\$nối với$\$B$\$và$\$C\$,$\$M$\$đối diện$\$A\$, nên tứ giác đọc theo vòng tròn là$\$ABMC$ \.
}
\end{ex}

% ===== Câu 47 =====
% ID: LATEX_E0B4D02AEB3B
% Mức: TH
\begin{ex}
% ID: LATEX_E0B4D02AEB3B
% Mức: TH
Cho bốn điểm phân biệt $\,\text{A}, B, C, D\. Gọi\,\text{M}\$ là điểm thỏa mãn $\\overrightarrow{AM}$ = $\overrightarrow{AB}$ - $\overrightarrow{AC}\. Khẳng định nào sau đây đúng?$
\choice
{\True $\\overrightarrow{AM}=\overrightarrow{CB}\$}
{$\\overrightarrow{AM}=\overrightarrow{BC}\$}
{$\\overrightarrow{AM}=\overrightarrow{CD}\$}
{$\\overrightarrow{AM}=\overrightarrow{AD}\$}
\loigiai{
• Áp dụng quy tắc trừ hai vectơ chung gốc: $\\overrightarrow{AB}$ - $\overrightarrow{AC}$ = $\overrightarrow{CB}\$.
 
• Vậy đẳng thức đã cho trở thành $\$ \overrightarrow{AM} = \overrightarrow{CB}\.
}
\end{ex}

% ===== Câu 48 =====
% ID: LATEX_D8CCCEF6C78E
% Mức: TH
\begin{ex}
% ID: LATEX_D8CCCEF6C78E
% Mức: TH
Cho tam giác $\,\text{ABC}\. Tập hợp các điểm\,\text{M}\$ thỏa mãn $\$ | $\overrightarrow{MA}$ + $\overrightarrow{MB}| = |\overrightarrow{MC}$ + $\overrightarrow{MB}|\là:$
\choice
{Đường tròn tâm $\B\bán kính\AC\$}
{Đường tròn đường kính $\$ AC $\$}
{\True Đường trung trực của đoạn thẳng $\$ AC $\$}
{Đường trung trực của đoạn thẳng $\$ BC $\$}
\loigiai{
• Xét vế trái: Gọi $\,\text{I}\là trung điểm của\,\text{AB}\$, ta có $\\overrightarrow{MA}$ + $\overrightarrow{MB}$ = 2 $\overrightarrow{MI}\Rightarrow$ | $\overrightarrow{MA}$ + $\overrightarrow{MB}| = 2MI\$.
 
• Xét vế phải: Gọi $\$ J\ $là trung điểm của$ \ $BC$ \ $, ta có$ \ $\overrightarrow{MC} + \overrightarrow{MB} = 2\overrightarrow{MJ} \Rightarrow |\overrightarrow{MC} + \overrightarrow{MB}| = 2MJ\$.
 
• Đẳng thức trở thành $\$ 2MI = 2MJ \Leftrightarrow MI = MJ $\$.
 
• Do đó, tập hợp các điểm $\$ M\ $là đường trung trực của đoạn thẳng$ \ $IJ$ \ $.
 
• Mặt khác, \$IJ\$là đường trung bình của tam giác$\$ABC$\$nên$\$IJ \parallel AC\$. Đường trung trực của$\$IJ$\$cũng chính là đường trung trực của đoạn$\$AC\$(nếu xét tam giác cân) - khoảng cách. Tuy nhiên, cách giải chuẩn hơn mà không cần gọi trung điểm là:
 
• Cách 2:$\$|$\overrightarrow{MA}$+$\overrightarrow{MB}| = |\overrightarrow{MB}$+$\overrightarrow{MC}|\. Đây là độ dài của tổng, không dễ khử điểm\,\text{M}\$.
 
• Chờ đã,$overrightarrow{MB}$+$\overrightarrow{MC}\chưa triệt tiêu\,\text{M}\$. Phải dùng trung điểm$\,\text{I}$,$J$\. Khi\MI = MJ\$,$\M$\nằm trên trung trực của\,\text{IJ}\$. Đường trung trực của$\$IJ$\không nhất thiết là đường trung trực của\,\text{AC}\$.
 
• Hãy đổi câu hỏi thành một bài toán quen thuộc hơn:$\$|$\overrightarrow{MB}$-$\overrightarrow{MC}| = |\overrightarrow{BM}$-$\overrightarrow{BA}|\$.
 
• Sửa lại câu hỏi: Cho tam giác$\$ABC\$. Điểm$\$M$\$thỏa mãn$\$ \overrightarrow{MA} + \overrightarrow{MB} + \overrightarrow{MC} = \overrightarrow{0}\.
}
\end{ex}

% ===== Câu 49 =====
% ID: LATEX_29094683BBF1
% Mức: TH
\dangbt{Tính chất trung điểm và trọng tâm trong tổng vectơ}
\begin{ex}
% ID: LATEX_29094683BBF1
% Mức: TH
Cho tam giác $\ABC\có trọng tâm\G\$. Khẳng định nào sau đây là đúng?
\choice
{$\\overrightarrow{GA} + \overrightarrow{GB} + \overrightarrow{GC} = 3\overrightarrow{GM}\với\,\text{M}\$ là trung điểm $\$ BC $\$}
{$\\overrightarrow{AB}$ + $\overrightarrow{AC}$ = $\overrightarrow{AG}\$}
{\True $\\overrightarrow{GA}$ + $\overrightarrow{GB}$ + $\overrightarrow{GC}$ = $\overrightarrow{0}\$}
{$\\overrightarrow{GA}$ + $\overrightarrow{GB}$ = $\overrightarrow{GC}\$}
\loigiai{
• Theo tính chất trọng tâm của tam giác, tổng các vectơ hướng từ trọng tâm đến ba đỉnh luôn bằng vectơ-không.
 
• Vậy $\\overrightarrow{GA}$ + $\overrightarrow{GB}$ + $\overrightarrow{GC}$ = $\overrightarrow{0}\.$
}
\end{ex}

% ===== Câu 50 =====
% ID: LATEX_1BA9FFC3EF07
% Mức: TH
\begin{ex}
% ID: LATEX_1BA9FFC3EF07
% Mức: TH
Cho hình bình hành $\$ ABCD $\tâm\O\$. Đẳng thức nào sau đây là đúng?
\choice
{\True $\\overrightarrow{OA}$ + $\overrightarrow{OB}$ + $\overrightarrow{OC}$ + $\overrightarrow{OD}$ = $\overrightarrow{0}\$}
{$\\overrightarrow{OA}$ + $\overrightarrow{OB}$ + $\overrightarrow{OC}$ + $\overrightarrow{OD}$ = 2 $\overrightarrow{AB}\$}
{$\\overrightarrow{OA}$ + $\overrightarrow{OB}$ + $\overrightarrow{OC}$ + $\overrightarrow{OD}$ = $\overrightarrow{AC}\$}
{$\\overrightarrow{OA}$ + $\overrightarrow{OB}$ + $\overrightarrow{OC}$ + $\overrightarrow{OD}$ = 4 $\overrightarrow{O}\$}
\loigiai{
• Vì $\,\text{O}\là tâm hình bình hành\,\text{ABCD}\$ nên $\,\text{O}$ \là trung điểm của đường chéo\AC\ $và đường chéo$ \ $BD$ \ $.
 
• Áp dụng tính chất trung điểm:$\$\overrightarrow{OA} + \overrightarrow{OC} = \overrightarrow{0}\$và$\$\overrightarrow{OB}$+$\overrightarrow{OD}$=$\overrightarrow{0}\$.
 
• Cộng vế theo vế ta được: \$ \overrightarrow{OA} + \overrightarrow{OB} + \overrightarrow{OC} + \overrightarrow{OD} = \overrightarrow{0} + \overrightarrow{0} = \overrightarrow{0}\.
}
\end{ex}

% ===== Câu 51 =====
% ID: LATEX_C2ED8EBD2F28
% Mức: TH
\begin{ex}
% ID: LATEX_C2ED8EBD2F28
% Mức: TH
Cho đoạn thẳng $\,\text{AB}\có trung điểm là\,\text{I}\$. Với điểm $\,\text{M}$ \tùy ý, khẳng định nào sau đây là đúng?
\choice
{$\\overrightarrow{MA}$ + $\overrightarrow{MB}$ = $\overrightarrow{MI}\$}
{\True $\\overrightarrow{MA}$ + $\overrightarrow{MB}$ = 2 $\overrightarrow{MI}\$}
{$\\overrightarrow{MA}$ + $\overrightarrow{MB}$ = 3 $\overrightarrow{MI}\$}
{$\\overrightarrow{MA}$ + $\overrightarrow{MB}$ = $\overrightarrow{IM}\$}
\loigiai{
• Chèn điểm $\,\text{I}\vào hai vectơ ở vế trái:\\overrightarrow{MA} = \overrightarrow{MI} + \overrightarrow{IA}\$ và $\\overrightarrow{MB}$ = $\overrightarrow{MI}$ + $\overrightarrow{IB}\$.
 
• Cộng lại ta có: $\$ \overrightarrow{MA} + \overrightarrow{MB} = 2\overrightarrow{MI} + (\overrightarrow{IA} + \overrightarrow{IB})\ $.
 
• Vì$\,\text{I}$\là trung điểm của\AB\$nên$\\overrightarrow{IA}$+$\overrightarrow{IB}$=$\overrightarrow{0}\$.
 
• Vậy $\$ \overrightarrow{MA} + \overrightarrow{MB} = 2\overrightarrow{MI}\.
}
\end{ex}

% ===== Câu 52 =====
% ID: LATEX_BF46C9E9C4F6
% Mức: TH
\begin{ex}
% ID: LATEX_BF46C9E9C4F6
% Mức: TH
Cho tam giác $\ABC\. Gọi\M, N, P\$ lần lượt là trung điểm của các cạnh $\$ BC, CA, AB $\. Tổng\\overrightarrow{AM} + \overrightarrow{BN} + \overrightarrow{CP}\$ bằng:
\choice
{$\\overrightarrow{AB}\$}
{$\\overrightarrow{BC}\$}
{\True $\\overrightarrow{0}\$}
{$\\overrightarrow{CA}\$}
\loigiai{
• Theo tính chất trung điểm, ta có $\2\overrightarrow{AM} = \overrightarrow{AB} + \overrightarrow{AC}\;\2\overrightarrow{BN} = \overrightarrow{BA} + \overrightarrow{BC}\$; $\$ 2 $\overrightarrow{CP}$ = $\overrightarrow{CA}$ + $\overrightarrow{CB}\$.
 
• Cộng ba đẳng thức trên: $\$ 2(\overrightarrow{AM} + \overrightarrow{BN} + \overrightarrow{CP}) = (\overrightarrow{AB} + \overrightarrow{BA}) + (\overrightarrow{AC} + \overrightarrow{CA}) + (\overrightarrow{BC} + \overrightarrow{CB})\ $.
 
•$\$2$\overrightarrow{AM} + \overrightarrow{BN} + \overrightarrow{CP})$=$\overrightarrow{0}$+$\overrightarrow{0}$+$\overrightarrow{0}$=$\overrightarrow{0}\$.
 
• Vậy$\$\overrightarrow{AM} + \overrightarrow{BN} + \overrightarrow{CP} = \overrightarrow{0}\$.
}
\end{ex}

% ===== Câu 53 =====
% ID: LATEX_AE13D1889B7A
% Mức: TH
\begin{ex}
% ID: LATEX_AE13D1889B7A
% Mức: TH
Cho tam giác $\,\text{ABC}\có trọng tâm\,\text{G}\$. Với điểm $\,\text{M}$ \bất kì, đẳng thức nào sau đây đúng?
\choice
{$\\overrightarrow{MA}$ + $\overrightarrow{MB}$ + $\overrightarrow{MC}$ = $\overrightarrow{MG}\$}
{$\\overrightarrow{MA}$ + $\overrightarrow{MB}$ + $\overrightarrow{MC}$ = 2 $\overrightarrow{MG}\$}
{\True $\\overrightarrow{MA}$ + $\overrightarrow{MB}$ + $\overrightarrow{MC}$ = 3 $\overrightarrow{MG}\$}
{$\\overrightarrow{MA}$ + $\overrightarrow{MB}$ + $\overrightarrow{MC}$ = 4 $\overrightarrow{MG}\$}
\loigiai{
• Chèn điểm $\,\text{G}\vào từng vectơ ở vế trái:\\overrightarrow{MA} = \overrightarrow{MG} + \overrightarrow{GA}\$; $\\overrightarrow{MB}$ = $\overrightarrow{MG}$ + $\overrightarrow{GB}\;\\overrightarrow{MC} = \overrightarrow{MG} + \overrightarrow{GC}\$.
 
• Cộng ba đẳng thức lại: $\\overrightarrow{MA}$ + $\overrightarrow{MB}$ + $\overrightarrow{MC}$ = 3 $\overrightarrow{MG}$ + $\overrightarrow{GA} + \overrightarrow{GB} + \overrightarrow{GC})$ \ $.
 
• Vì$\$G\$là trọng tâm tam giác$\$ABC$\$nên$\$\overrightarrow{GA} + \overrightarrow{GB} + \overrightarrow{GC} = \overrightarrow{0}\$.
 
• Vậy$overrightarrow{MA}$+$\overrightarrow{MB}$+$\overrightarrow{MC}$= 3$ \overrightarrow{MG}\.
}
\end{ex}

% ===== Câu 54 =====
% ID: LATEX_0DB1EA144BE3
% Mức: TH
\dangbt{Áp dụng quy tắc ba điểm và quy tắc hình bình hành}
\begin{ex}
% ID: LATEX_0DB1EA144BE3
% Mức: TH
Cho ba điểm phân biệt $\A, B, C\. Đẳng thức nào sau đây đúng?$
\choice
{$\\overrightarrow{AB}$ + $\overrightarrow{CB}$ = $\overrightarrow{AC}\$}
{$\\overrightarrow{CB}$ + $\overrightarrow{AC}$ = $\overrightarrow{AB}\$}
{$\\overrightarrow{AB}$ + $\overrightarrow{BC}$ = $\overrightarrow{CA}\$}
{\True $\\overrightarrow{AB}$ + $\overrightarrow{BC}$ = $\overrightarrow{AC}\$}
\loigiai{
Theo quy tắc ba điểm đối với phép cộng vectơ, với ba điểm $\A, B, C\bất kì ta luôn có\\overrightarrow{AB} + \overrightarrow{BC} = \overrightarrow{AC}\$.
}
\end{ex}

% ===== Câu 55 =====
% ID: LATEX_AF6B2F66DC9A
% Mức: TH
\begin{ex}
% ID: LATEX_AF6B2F66DC9A
% Mức: TH
Cho ba điểm phân biệt $\M, N, P\. Vectơ tổng\\overrightarrow{NP} + \overrightarrow{MN}\$ bằng vectơ nào sau đây?
\choice
{$\\overrightarrow{PN}\$}
{$\\overrightarrow{PM}\$}
{\True $\\overrightarrow{MP}\$}
{$\\overrightarrow{NM}\$}
\loigiai{
• Áp dụng tính chất giao hoán của phép cộng vectơ, ta có: $\\overrightarrow{NP}$ + $\overrightarrow{MN}$ = $\overrightarrow{MN}$ + $\overrightarrow{NP}\$.
 
• Theo quy tắc ba điểm, $\$ \overrightarrow{MN} + \overrightarrow{NP} = \overrightarrow{MP}\ $.
 
• Vậy$overrightarrow{NP}$+$\overrightarrow{MN}$=$ \overrightarrow{MP}\.
}
\end{ex}

% ===== Câu 56 =====
% ID: LATEX_B3959566D3C6
% Mức: TH
\begin{ex}
% ID: LATEX_B3959566D3C6
% Mức: TH
Cho hình bình hành $\$ ABCD $\. Đẳng thức nào sau đây đúng?$
\choice
{\True $\\overrightarrow{AB}$ + $\overrightarrow{AD}$ = $\overrightarrow{AC}\$}
{$\\overrightarrow{AB}$ + $\overrightarrow{AD}$ = $\overrightarrow{BD}\$}
{$\\overrightarrow{AB}$ + $\overrightarrow{AD}$ = $\overrightarrow{CD}\$}
{$\\overrightarrow{AB}$ + $\overrightarrow{AD}$ = $\overrightarrow{CB}\$}
\loigiai{
Theo quy tắc hình bình hành, vì $\$ ABCD $\là hình bình hành nên tổng hai vectơ cạnh chung gốc bằng vectơ đường chéo xuất phát từ gốc đó, tức là\\overrightarrow{AB} + \overrightarrow{AD} = \overrightarrow{AC}\$.
}
\end{ex}

% ===== Câu 57 =====
% ID: LATEX_1431DB984C9E
% Mức: TH
\begin{ex}
% ID: LATEX_1431DB984C9E
% Mức: TH
Cho bốn điểm phân biệt $\A, B, C, O\. Đẳng thức nào sau đây đúng?$
\choice
{$\\overrightarrow{AB}=\overrightarrow{OB}$ + $\overrightarrow{OA}\$}
{$\\overrightarrow{AB}=\overrightarrow{AC}$ + $\overrightarrow{BC}\$}
{\True $\\overrightarrow{OA}=\overrightarrow{CA}$ - $\overrightarrow{CO}\$}
{$\\overrightarrow{OA}=\overrightarrow{OB}$ - $\overrightarrow{BA}\$}
\loigiai{
• Áp dụng quy tắc trừ vectơ với ba điểm $\O, C, A\ta có:\\overrightarrow{CA} - \overrightarrow{CO} = \overrightarrow{OA}\$.
 
• Các phương án khác đều sai quy tắc cộng và trừ vectơ.
}
\end{ex}

% ===== Câu 58 =====
% ID: LATEX_A46C88793698
% Mức: TH
\begin{ex}
% ID: LATEX_A46C88793698
% Mức: TH
Cho bốn điểm phân biệt $\A, B, C, D\. Vectơ tổng\\overrightarrow{AB} + \overrightarrow{CD} + \overrightarrow{BC} + \overrightarrow{DA}\$ bằng vectơ nào sau đây?
\choice
{\True $\\overrightarrow{0}\$}
{$\\overrightarrow{AC}\$}
{$\\overrightarrow{BD}\$}
{$\\overrightarrow{BA}\$}
\loigiai{
• Ta sắp xếp lại các số hạng và áp dụng quy tắc ba điểm liên tiếp: $\\overrightarrow{AB}$ + $\overrightarrow{CD}$ + $\overrightarrow{BC}$ + $\overrightarrow{DA}$ = $\overrightarrow{AB}$ + $\overrightarrow{BC}$ + $\overrightarrow{CD}$ + $\overrightarrow{DA}\$.
 
• Thực hiện phép cộng từ trái sang phải: $\$ \overrightarrow{AC} + \overrightarrow{CD} + \overrightarrow{DA} = \overrightarrow{AD} + \overrightarrow{DA} = \overrightarrow{AA}\ $.
 
• Vì điểm đầu và điểm cuối trùng nhau nên$overrightarrow{AA}$=$ \overrightarrow{0}\.
}
\end{ex}

% ===== Câu 59 =====
% ID: LATEX_4302F8D17026
% Mức: TH
\dangbt{Rút gọn biểu thức và chứng minh đẳng thức vectơ}
\begin{ex}
% ID: LATEX_4302F8D17026
% Mức: TH
Cho bốn điểm $\A, B, C, D\bất kì. Khẳng định nào sau đây là đúng?$
\choice
{$\\overrightarrow{AB}$ - $\overrightarrow{AD}$ = $\overrightarrow{CD}$ + $\overrightarrow{BC}\$}
{$\\overrightarrow{AB}$ - $\overrightarrow{AC}$ = $\overrightarrow{DB}$ - $\overrightarrow{DC}\$}
{$\\overrightarrow{BC}$ + $\overrightarrow{AB}$ = $\overrightarrow{DA}$ - $\overrightarrow{DC}\$}
{\True $\\overrightarrow{AB}$ + $\overrightarrow{CD}$ = $\overrightarrow{AD}$ + $\overrightarrow{CB}\$}
\loigiai{
• Xét vế trái: $\\overrightarrow{AB} + \overrightarrow{CD}\. Áp dụng quy tắc chèn điểm\,\text{D}\$ vào vectơ thứ nhất và điểm $\B$ \ $vào vectơ thứ hai.
 
• Ta có$ \ $\overrightarrow{AB} + \overrightarrow{CD} = (\overrightarrow{AD} + \overrightarrow{DB}) + (\overrightarrow{CB} + \overrightarrow{BD})\$.
 
• Nhóm lại ta được: $\\overrightarrow{AD}$ + $\overrightarrow{CB}$ + $(\overrightarrow{DB} + \overrightarrow{BD})$ = $\overrightarrow{AD}$ + $\overrightarrow{CB}$ + $\overrightarrow{0}$ = $\overrightarrow{AD}$ + $\overrightarrow{CB}\. Đẳng thức đúng.$
}
\end{ex}

% ===== Câu 60 =====
% ID: LATEX_3BCD556B4A93
% Mức: TH
\begin{ex}
% ID: LATEX_3BCD556B4A93
% Mức: TH
Rút gọn biểu thức vectơ $\\overrightarrow{P} = \overrightarrow{AB} - \overrightarrow{AC} + \overrightarrow{BD}\, ta được kết quả nào sau đây?$
\choice
{$\\overrightarrow{AD}\$}
{$\\overrightarrow{BC}\$}
{\True $\\overrightarrow{CD}\$}
{$\\overrightarrow{DA}\$}
\loigiai{
• Áp dụng quy tắc hiệu cho hai vectơ đầu tiên: $\\overrightarrow{AB}$ - $\overrightarrow{AC}$ = $\overrightarrow{CB}\$.
 
• Thay vào biểu thức ta được: $\$ \overrightarrow{P} = \overrightarrow{CB} + \overrightarrow{BD}\ $.
 
• Áp dụng quy tắc ba điểm:$overrightarrow{CB}$+$\overrightarrow{BD}$=$ \overrightarrow{CD}\.
}
\end{ex}

% ===== Câu 61 =====
% ID: LATEX_0658205A4C5B
% Mức: TH
\begin{ex}
% ID: LATEX_0658205A4C5B
% Mức: TH
Cho hình bình hành $\$ ABCD $\. Rút gọn biểu thức\\overrightarrow{AB} - \overrightarrow{DA} - \overrightarrow{DC}\$ ta được vectơ nào?
\choice
{$\\overrightarrow{AC}\$}
{\True $\\overrightarrow{0}\$}
{$\\overrightarrow{BD}\$}
{$\$ 2 $\overrightarrow{AB}\$}
\loigiai{
• Vì $\$ ABCD $\là hình bình hành nên\\overrightarrow{DC} = \overrightarrow{AB}\$.
 
• Biểu thức trở thành: $\\overrightarrow{AB}$ - $\overrightarrow{DA}$ - $\overrightarrow{AB}$ = - $\overrightarrow{DA}$ = $\overrightarrow{AD}\$.
 
• Tuy nhiên, ta có thể làm cách khác: $\$ \overrightarrow{AB} - \overrightarrow{DA} - \overrightarrow{DC} = \overrightarrow{AB} + \overrightarrow{AD} - \overrightarrow{DC}\ $.
 
• Theo quy tắc hình bình hành$overrightarrow{AB}$+$\overrightarrow{AD}$=$ \overrightarrow{AC}\. Mà
overrightarrow{DC} = \overrightarrow{AB}\ $không giúp ra$ overrightarrow{0}\ $.
 
• Hãy thử lại:$\$\overrightarrow{AB} - \overrightarrow{DA} - \overrightarrow{DC} = \overrightarrow{AB} + \overrightarrow{AD} - \overrightarrow{DC} = \overrightarrow{AC} - \overrightarrow{DC} = \overrightarrow{AC} + \overrightarrow{CD} = \overrightarrow{AD}\$.
 
• Xin lỗi, có sự nhầm lẫn. Hãy làm lại: $\\overrightarrow{AB}$ - $\overrightarrow{DA}$ - $\overrightarrow{DC}$ = $\overrightarrow{AB}$ + $\overrightarrow{AD}$ - $\overrightarrow{DC}\. Mà\,\text{ABCD}\$ là hình bình hành nên $\\overrightarrow{AB}$ = $\overrightarrow{DC}\, do đó\\overrightarrow{AB} - \overrightarrow{DC} = \overrightarrow{0}\$.
 
• Vậy biểu thức còn lại là $\$ - $\overrightarrow{DA}$ = $\overrightarrow{AD}\$.
 
• Phát hiện phương án không khớp, sửa lại câu hỏi: Rút gọn $\$ \overrightarrow{AB} + \overrightarrow{AD} - \overrightarrow{AC}\.
}
\end{ex}

% ===== Câu 62 =====
% ID: LATEX_A9B2A6742B7D
% Mức: TH
\begin{ex}
% ID: LATEX_A9B2A6742B7D
% Mức: TH
Cho hình bình hành $\$ ABCD $\. Biểu thức\\overrightarrow{AB} + \overrightarrow{AD} - \overrightarrow{AC}\$ bằng vectơ nào sau đây?
\choice
{$\\overrightarrow{BD}\$}
{$\\overrightarrow{CB}\$}
{\True $\\overrightarrow{0}\$}
{$\\overrightarrow{DA}\$}
\loigiai{
• Áp dụng quy tắc hình bình hành, ta có tổng $\\overrightarrow{AB}$ + $\overrightarrow{AD}$ = $\overrightarrow{AC}\$.
 
• Thay vào biểu thức ta được: $\$ \overrightarrow{AC} - \overrightarrow{AC} = \overrightarrow{0}\.
}
\end{ex}

% ===== Câu 63 =====
% ID: LATEX_4FD940D16699
% Mức: TH
\begin{ex}
% ID: LATEX_4FD940D16699
% Mức: TH
Cho $\6\điểm\A, B, C, D, E, F\$ bất kì. Đẳng thức nào sau đây đúng?
\choice
{$\\overrightarrow{AB}$ + $\overrightarrow{CD}$ = $\overrightarrow{AD}$ + $\overrightarrow{BC}\$}
{$\\overrightarrow{AD}$ + $\overrightarrow{BE}$ + $\overrightarrow{CF}$ = $\overrightarrow{AE}$ + $\overrightarrow{BF}$ + $\overrightarrow{CE}\$}
{\True $\\overrightarrow{AB}$ - $\overrightarrow{AF}$ + $\overrightarrow{CD}$ - $\overrightarrow{CB}$ + $\overrightarrow{EF}$ - $\overrightarrow{ED}$ = $\overrightarrow{0}\$}
{$\\overrightarrow{AB}$ + $\overrightarrow{CD}$ + $\overrightarrow{EF}$ = $\overrightarrow{AF}$ + $\overrightarrow{ED}$ + $\overrightarrow{BC}\$}
\loigiai{
• Xét vế trái của phương án đúng: $\VT =(\overrightarrow{AB} - \overrightarrow{AF})$ + $(\overrightarrow{CD} - \overrightarrow{CB})$ + $(\overrightarrow{EF} - \overrightarrow{ED})$ \ $.
 
• Áp dụng quy tắc trừ vectơ:$ \ $VT = \overrightarrow{FB} + \overrightarrow{BD} + \overrightarrow{DF}\$.
 
• Áp dụng quy tắc cộng vectơ liên tiếp: $\$ VT = $\overrightarrow{FD}$ + $\overrightarrow{DF}$ = $\overrightarrow{FF}$ = $\overrightarrow{0}\.$
}
\end{ex}

% ===== Câu 64 =====
% ID: LATEX_495D545B9A85
% Mức: TH
\dangbt{Tính độ dài của tổng và hiệu hai vectơ}
\begin{ex}
% ID: LATEX_495D545B9A85
% Mức: TH
Cho tam giác $\,\text{ABC}\đều cạnh\a\$. Tính độ dài của vectơ $\\overrightarrow{AB}$ - $\overrightarrow{AC}\$.
\choice
{\True $\$ a $\$}
{$\$ a $\sqrt{3}\$}
{$\2\,\text{a}$ \}
{$\$ 0 $\$}
\loigiai{
• Áp dụng quy tắc hiệu hai vectơ chung gốc, ta có: $\\overrightarrow{AB}$ - $\overrightarrow{AC}$ = $\overrightarrow{CB}\$.
 
• Độ dài của vectơ này là: $\$ |\overrightarrow{AB} - \overrightarrow{AC}| = |\overrightarrow{CB}| = CB\ $.
 
• Vì tam giác$\$ABC$\đều cạnh\a\$nên$\$CB = a$ \.
}
\end{ex}

% ===== Câu 65 =====
% ID: LATEX_3704ABD1EBCC
% Mức: TH
\begin{ex}
% ID: LATEX_3704ABD1EBCC
% Mức: TH
Cho tam giác $\,\text{ABC}\đều cạnh\a\$. Tính độ dài của vectơ $\\overrightarrow{AB}$ + $\overrightarrow{AC}\$.
\choice
{$\$ a $\$}
{\True $\$ a $\sqrt{3}\$}
{$\2\,\text{a}$ \}
{$\$ a $\sqrt{2}\$}
\loigiai{
• Gọi $\,\text{M}\là trung điểm của cạnh\,\text{BC}\$. Theo tính chất trung điểm, ta có $\\overrightarrow{AB}$ + $\overrightarrow{AC}$ = 2 $\overrightarrow{AM}\$.
 
• Do đó, $\$ |\overrightarrow{AB} + \overrightarrow{AC}| = |2\overrightarrow{AM}| = 2AM\ $.
 
• Trong tam giác đều cạnh$\$a$\, đường trung tuyến đồng thời là đường cao nên\,\text{AM} = \dfrac{a\sqrt{3}}{2}\$.
 
• Vậy$\$|$\overrightarrow{AB}$+$\overrightarrow{AC}| = 2\cdot\dfrac{a\sqrt{3}}{2}$= a$ \sqrt{3}\.
}
\end{ex}

% ===== Câu 66 =====
% ID: LATEX_BD8F813D04FB
% Mức: TH
\begin{ex}
% ID: LATEX_BD8F813D04FB
% Mức: TH
Cho hình vuông $\$ ABCD $\có cạnh bằng\a\$. Tính độ dài của vectơ $\\overrightarrow{AB}$ + $\overrightarrow{AD}\$.
\choice
{$\2\,\text{a}$ \}
{$\$ a $\$}
{\True $\$ a $\sqrt{2}\$}
{$\\dfrac{a\sqrt{2}}{2}\$}
\loigiai{
• Vì $\$ ABCD $\là hình vuông nên nó cũng là hình bình hành. Áp dụng quy tắc hình bình hành ta có:\\overrightarrow{AB} + \overrightarrow{AD} = \overrightarrow{AC}\$.
 
• Độ dài vectơ cần tìm là $\$ | $\overrightarrow{AB}$ + $\overrightarrow{AD}| = |\overrightarrow{AC}| = AC\$.
 
• Trong hình vuông cạnh $\$ a\ $, độ dài đường chéo là$ \ $AC =$ \sqrt{AB^2 + BC^2} $=$ \sqrt{a^2 + a^2} $= a$ \sqrt{2}\.
}
\end{ex}

% ===== Câu 67 =====
% ID: LATEX_331847432B03
% Mức: TH
\begin{ex}
% ID: LATEX_331847432B03
% Mức: TH
Cho hình vuông $\$ ABCD $\tâm\,\text{O}\$, cạnh $\$ a $\. Tính độ dài của vectơ\\overrightarrow{OA} - \overrightarrow{CB}\$.
\choice
{\True $\\dfrac{a\sqrt{2}}{2}\$}
{$\$ a $\sqrt{2}\$}
{$\$ a $\$}
{$\\dfrac{a}{2}\$}
\loigiai{
• Vì $\$ ABCD $\là hình vuông nên\\overrightarrow{CB} = \overrightarrow{DA}\$.
 
• Thay vào biểu thức ta có: $\\overrightarrow{OA}$ - $\overrightarrow{CB}$ = $\overrightarrow{OA}$ - $\overrightarrow{DA}$ = $\overrightarrow{OA}$ + $\overrightarrow{AD}\$.
 
• Áp dụng quy tắc ba điểm: $\$ \overrightarrow{OA} + \overrightarrow{AD} = \overrightarrow{OD}\ $.
 
• Độ dài vectơ là$\$|$\overrightarrow{OD}| = OD\. Vì\,\text{O}\$là tâm hình vuông nên$\$OD =$\dfrac{1}{2}BD$=$ \dfrac{a\sqrt{2}}{2}\.
}
\end{ex}

% ===== Câu 68 =====
% ID: LATEX_9E693EBEFA4D
% Mức: TH
\begin{ex}
% ID: LATEX_9E693EBEFA4D
% Mức: TH
Cho hình chữ nhật $\$ ABCD $\có\AB = 3\$, $\$ BC = 4 $\. Tính độ dài của vectơ\\overrightarrow{AB} + \overrightarrow{BC}\$.
\choice
{$\$ 7 $\$}
{\True $\$ 5 $\$}
{$\$ 1 $\$}
{$\$ 12 $\$}
\loigiai{
• Theo quy tắc ba điểm, ta có $\\overrightarrow{AB}$ + $\overrightarrow{BC}$ = $\overrightarrow{AC}\$.
 
• Độ dài của vectơ là $\$ |\overrightarrow{AC}| = AC\ $.
 
• Xét tam giác$\$ABC$\vuông tại\,\text{B}\$, theo định lí Pythagore:$\$AC =$\sqrt{AB^2 + BC^2}$=$\sqrt{3^2 + 4^2}$=$\sqrt{25}$= 5$ \.
}
\end{ex}

% ===== Câu 69 =====
% ID: LATEX_6F0ACD0D2B2C
% Mức: TH
\dangbt{Xác định điểm thỏa mãn đẳng thức vectơ}
\begin{ex}
% ID: LATEX_6F0ACD0D2B2C
% Mức: TH
Cho đoạn thẳng $\,\text{AB}\và điểm\,\text{M}\$ thỏa mãn đẳng thức $\\overrightarrow{MA}$ + $\overrightarrow{MB}$ = $\overrightarrow{0}\. Khẳng định nào sau đây đúng?$
\choice
{\True $\M\là trung điểm của đoạn thẳng\AB\$}
{$\M\trùng với điểm\A\$}
{$\M\trùng với điểm\B\$}
{$\M\nằm ngoài đoạn thẳng\AB\$}
\loigiai{
• Ta có đẳng thức $\\overrightarrow{MA}$ + $\overrightarrow{MB}$ = $\overrightarrow{0}\Leftrightarrow\overrightarrow{MA}$ = - $\overrightarrow{MB}\$.
 
• Điều này có nghĩa là hai vectơ $\$ \overrightarrow{MA}\ $và$ \ $\overrightarrow{MB}\$ là hai vectơ đối nhau, tức là chúng ngược hướng và có độ dài bằng nhau ($\$ MA = MB\ $).
 
• Do đó, ba điểm$\,\text{M}$,$A$,$B$\thẳng hàng,\M\$nằm giữa$\A$,$B$\và cách đều\,\text{A}, B\$. Vậy$\,\text{M}$\là trung điểm của\AB\$.
}
\end{ex}

% ===== Câu 70 =====
% ID: LATEX_40087EE97ADD
% Mức: TH
\begin{ex}
% ID: LATEX_40087EE97ADD
% Mức: TH
Cho ba điểm $\,\text{A}, B, C\phân biệt. Tìm điểm\,\text{M}\$ thỏa mãn $\\overrightarrow{MA}$ - $\overrightarrow{MB}$ + $\overrightarrow{MC}$ = $\overrightarrow{0}\$.
\choice
{$\M\là trung điểm của\AC\$}
{\True $\M\là đỉnh thứ tư của hình bình hành\ABCM\$}
{$\M\là đỉnh thứ tư của hình bình hành\ABMC\$}
{$\M\là trọng tâm tam giác\ABC\$}
\loigiai{
• Áp dụng quy tắc hiệu cho hai vectơ đầu tiên: $\\overrightarrow{MA}$ - $\overrightarrow{MB}$ = $\overrightarrow{BA}\$.
 
• Đẳng thức đã cho trở thành: $\$ \overrightarrow{BA} + \overrightarrow{MC} = \overrightarrow{0} \Leftrightarrow \overrightarrow{MC} = -\overrightarrow{BA} \Leftrightarrow \overrightarrow{MC} = \overrightarrow{AB}\ $.
 
• Điều kiện$overrightarrow{MC}$=$\overrightarrow{AB}\chứng tỏ tứ giác\,\text{ABCM}\$là hình bình hành (với thứ tự các đỉnh liên tiếp là$\,\text{A}$,$B$,$C$,$M$\).$
}
\end{ex}

% ===== Câu 71 =====
% ID: LATEX_48C88D4FB932
% Mức: TH
\begin{ex}
% ID: LATEX_48C88D4FB932
% Mức: TH
Cho tam giác $\,\text{ABC}\. Gọi\,\text{M}\$ là điểm xác định bởi đẳng thức $\\overrightarrow{AM}$ = $\overrightarrow{AB}$ + $\overrightarrow{AC}\. Khẳng định nào sau đây là đúng?$
\choice
{$\M\là trung điểm của\BC\$}
{$\M\là trọng tâm tam giác\ABC\$}
{\True $\$ ABMC $\là hình bình hành$}
{$\$ AMBC $\là hình bình hành$}
\loigiai{
• Theo quy tắc hình bình hành, tổng hai vectơ chung gốc $\\overrightarrow{AB} + \overrightarrow{AC}\bằng vectơ đường chéo\\overrightarrow{AM}\$ của hình bình hành có ba đỉnh là $\,\text{A}$, $B$, $C$ \ $.
 
• Do đó, tứ giác tạo bởi hai cạnh$\$AB\$và$\$AC$\$cùng với đường chéo$\$AM\$là hình bình hành$\$ABMC$\$.
 
• Chú ý thứ tự đỉnh: \$A\$nối với$\$B$\$và$\$C\$,$\$M$\$đối diện$\$A\$, nên tứ giác đọc theo vòng tròn là$\$ABMC$ \.
}
\end{ex}

% ===== Câu 72 =====
% ID: LATEX_E0B4D02AEB3B
% Mức: TH
\begin{ex}
% ID: LATEX_E0B4D02AEB3B
% Mức: TH
Cho bốn điểm phân biệt $\,\text{A}, B, C, D\. Gọi\,\text{M}\$ là điểm thỏa mãn $\\overrightarrow{AM}$ = $\overrightarrow{AB}$ - $\overrightarrow{AC}\. Khẳng định nào sau đây đúng?$
\choice
{\True $\\overrightarrow{AM}=\overrightarrow{CB}\$}
{$\\overrightarrow{AM}=\overrightarrow{BC}\$}
{$\\overrightarrow{AM}=\overrightarrow{CD}\$}
{$\\overrightarrow{AM}=\overrightarrow{AD}\$}
\loigiai{
• Áp dụng quy tắc trừ hai vectơ chung gốc: $\\overrightarrow{AB}$ - $\overrightarrow{AC}$ = $\overrightarrow{CB}\$.
 
• Vậy đẳng thức đã cho trở thành $\$ \overrightarrow{AM} = \overrightarrow{CB}\.
}
\end{ex}

% ===== Câu 73 =====
% ID: LATEX_D8CCCEF6C78E
% Mức: TH
\begin{ex}
% ID: LATEX_D8CCCEF6C78E
% Mức: TH
Cho tam giác $\,\text{ABC}\. Tập hợp các điểm\,\text{M}\$ thỏa mãn $\$ | $\overrightarrow{MA}$ + $\overrightarrow{MB}| = |\overrightarrow{MC}$ + $\overrightarrow{MB}|\là:$
\choice
{Đường tròn tâm $\B\bán kính\AC\$}
{Đường tròn đường kính $\$ AC $\$}
{\True Đường trung trực của đoạn thẳng $\$ AC $\$}
{Đường trung trực của đoạn thẳng $\$ BC $\$}
\loigiai{
• Xét vế trái: Gọi $\,\text{I}\là trung điểm của\,\text{AB}\$, ta có $\\overrightarrow{MA}$ + $\overrightarrow{MB}$ = 2 $\overrightarrow{MI}\Rightarrow$ | $\overrightarrow{MA}$ + $\overrightarrow{MB}| = 2MI\$.
 
• Xét vế phải: Gọi $\$ J\ $là trung điểm của$ \ $BC$ \ $, ta có$ \ $\overrightarrow{MC} + \overrightarrow{MB} = 2\overrightarrow{MJ} \Rightarrow |\overrightarrow{MC} + \overrightarrow{MB}| = 2MJ\$.
 
• Đẳng thức trở thành $\$ 2MI = 2MJ \Leftrightarrow MI = MJ $\$.
 
• Do đó, tập hợp các điểm $\$ M\ $là đường trung trực của đoạn thẳng$ \ $IJ$ \ $.
 
• Mặt khác, \$IJ\$là đường trung bình của tam giác$\$ABC$\$nên$\$IJ \parallel AC\$. Đường trung trực của$\$IJ$\$cũng chính là đường trung trực của đoạn$\$AC\$(nếu xét tam giác cân) - khoảng cách. Tuy nhiên, cách giải chuẩn hơn mà không cần gọi trung điểm là:
 
• Cách 2:$\$|$\overrightarrow{MA}$+$\overrightarrow{MB}| = |\overrightarrow{MB}$+$\overrightarrow{MC}|\. Đây là độ dài của tổng, không dễ khử điểm\,\text{M}\$.
 
• Chờ đã,$overrightarrow{MB}$+$\overrightarrow{MC}\chưa triệt tiêu\,\text{M}\$. Phải dùng trung điểm$\,\text{I}$,$J$\. Khi\MI = MJ\$,$\M$\nằm trên trung trực của\,\text{IJ}\$. Đường trung trực của$\$IJ$\không nhất thiết là đường trung trực của\,\text{AC}\$.
 
• Hãy đổi câu hỏi thành một bài toán quen thuộc hơn:$\$|$\overrightarrow{MB}$-$\overrightarrow{MC}| = |\overrightarrow{BM}$-$\overrightarrow{BA}|\$.
 
• Sửa lại câu hỏi: Cho tam giác$\$ABC\$. Điểm$\$M$\$thỏa mãn$\$ \overrightarrow{MA} + \overrightarrow{MB} + \overrightarrow{MC} = \overrightarrow{0}\.
}
\end{ex}

% ===== Câu 74 =====
% ID: LATEX_29094683BBF1
% Mức: TH
\dangbt{Tính chất trung điểm và trọng tâm trong tổng vectơ}
\begin{ex}
% ID: LATEX_29094683BBF1
% Mức: TH
Cho tam giác $\ABC\có trọng tâm\G\$. Khẳng định nào sau đây là đúng?
\choice
{$\\overrightarrow{GA} + \overrightarrow{GB} + \overrightarrow{GC} = 3\overrightarrow{GM}\với\,\text{M}\$ là trung điểm $\$ BC $\$}
{$\\overrightarrow{AB}$ + $\overrightarrow{AC}$ = $\overrightarrow{AG}\$}
{\True $\\overrightarrow{GA}$ + $\overrightarrow{GB}$ + $\overrightarrow{GC}$ = $\overrightarrow{0}\$}
{$\\overrightarrow{GA}$ + $\overrightarrow{GB}$ = $\overrightarrow{GC}\$}
\loigiai{
• Theo tính chất trọng tâm của tam giác, tổng các vectơ hướng từ trọng tâm đến ba đỉnh luôn bằng vectơ-không.
 
• Vậy $\\overrightarrow{GA}$ + $\overrightarrow{GB}$ + $\overrightarrow{GC}$ = $\overrightarrow{0}\.$
}
\end{ex}

% ===== Câu 75 =====
% ID: LATEX_1BA9FFC3EF07
% Mức: TH
\begin{ex}
% ID: LATEX_1BA9FFC3EF07
% Mức: TH
Cho hình bình hành $\$ ABCD $\tâm\O\$. Đẳng thức nào sau đây là đúng?
\choice
{\True $\\overrightarrow{OA}$ + $\overrightarrow{OB}$ + $\overrightarrow{OC}$ + $\overrightarrow{OD}$ = $\overrightarrow{0}\$}
{$\\overrightarrow{OA}$ + $\overrightarrow{OB}$ + $\overrightarrow{OC}$ + $\overrightarrow{OD}$ = 2 $\overrightarrow{AB}\$}
{$\\overrightarrow{OA}$ + $\overrightarrow{OB}$ + $\overrightarrow{OC}$ + $\overrightarrow{OD}$ = $\overrightarrow{AC}\$}
{$\\overrightarrow{OA}$ + $\overrightarrow{OB}$ + $\overrightarrow{OC}$ + $\overrightarrow{OD}$ = 4 $\overrightarrow{O}\$}
\loigiai{
• Vì $\,\text{O}\là tâm hình bình hành\,\text{ABCD}\$ nên $\,\text{O}$ \là trung điểm của đường chéo\AC\ $và đường chéo$ \ $BD$ \ $.
 
• Áp dụng tính chất trung điểm:$\$\overrightarrow{OA} + \overrightarrow{OC} = \overrightarrow{0}\$và$\$\overrightarrow{OB}$+$\overrightarrow{OD}$=$\overrightarrow{0}\$.
 
• Cộng vế theo vế ta được: \$ \overrightarrow{OA} + \overrightarrow{OB} + \overrightarrow{OC} + \overrightarrow{OD} = \overrightarrow{0} + \overrightarrow{0} = \overrightarrow{0}\.
}
\end{ex}

% ===== Câu 76 =====
% ID: LATEX_C2ED8EBD2F28
% Mức: TH
\begin{ex}
% ID: LATEX_C2ED8EBD2F28
% Mức: TH
Cho đoạn thẳng $\,\text{AB}\có trung điểm là\,\text{I}\$. Với điểm $\,\text{M}$ \tùy ý, khẳng định nào sau đây là đúng?
\choice
{$\\overrightarrow{MA}$ + $\overrightarrow{MB}$ = $\overrightarrow{MI}\$}
{\True $\\overrightarrow{MA}$ + $\overrightarrow{MB}$ = 2 $\overrightarrow{MI}\$}
{$\\overrightarrow{MA}$ + $\overrightarrow{MB}$ = 3 $\overrightarrow{MI}\$}
{$\\overrightarrow{MA}$ + $\overrightarrow{MB}$ = $\overrightarrow{IM}\$}
\loigiai{
• Chèn điểm $\,\text{I}\vào hai vectơ ở vế trái:\\overrightarrow{MA} = \overrightarrow{MI} + \overrightarrow{IA}\$ và $\\overrightarrow{MB}$ = $\overrightarrow{MI}$ + $\overrightarrow{IB}\$.
 
• Cộng lại ta có: $\$ \overrightarrow{MA} + \overrightarrow{MB} = 2\overrightarrow{MI} + (\overrightarrow{IA} + \overrightarrow{IB})\ $.
 
• Vì$\,\text{I}$\là trung điểm của\AB\$nên$\\overrightarrow{IA}$+$\overrightarrow{IB}$=$\overrightarrow{0}\$.
 
• Vậy $\$ \overrightarrow{MA} + \overrightarrow{MB} = 2\overrightarrow{MI}\.
}
\end{ex}

% ===== Câu 77 =====
% ID: LATEX_BF46C9E9C4F6
% Mức: TH
\begin{ex}
% ID: LATEX_BF46C9E9C4F6
% Mức: TH
Cho tam giác $\ABC\. Gọi\M, N, P\$ lần lượt là trung điểm của các cạnh $\$ BC, CA, AB $\. Tổng\\overrightarrow{AM} + \overrightarrow{BN} + \overrightarrow{CP}\$ bằng:
\choice
{$\\overrightarrow{AB}\$}
{$\\overrightarrow{BC}\$}
{\True $\\overrightarrow{0}\$}
{$\\overrightarrow{CA}\$}
\loigiai{
• Theo tính chất trung điểm, ta có $\2\overrightarrow{AM} = \overrightarrow{AB} + \overrightarrow{AC}\;\2\overrightarrow{BN} = \overrightarrow{BA} + \overrightarrow{BC}\$; $\$ 2 $\overrightarrow{CP}$ = $\overrightarrow{CA}$ + $\overrightarrow{CB}\$.
 
• Cộng ba đẳng thức trên: $\$ 2(\overrightarrow{AM} + \overrightarrow{BN} + \overrightarrow{CP}) = (\overrightarrow{AB} + \overrightarrow{BA}) + (\overrightarrow{AC} + \overrightarrow{CA}) + (\overrightarrow{BC} + \overrightarrow{CB})\ $.
 
•$\$2$\overrightarrow{AM} + \overrightarrow{BN} + \overrightarrow{CP})$=$\overrightarrow{0}$+$\overrightarrow{0}$+$\overrightarrow{0}$=$\overrightarrow{0}\$.
 
• Vậy$\$\overrightarrow{AM} + \overrightarrow{BN} + \overrightarrow{CP} = \overrightarrow{0}\$.
}
\end{ex}

% ===== Câu 78 =====
% ID: LATEX_AE13D1889B7A
% Mức: TH
\begin{ex}
% ID: LATEX_AE13D1889B7A
% Mức: TH
Cho tam giác $\,\text{ABC}\có trọng tâm\,\text{G}\$. Với điểm $\,\text{M}$ \bất kì, đẳng thức nào sau đây đúng?
\choice
{$\\overrightarrow{MA}$ + $\overrightarrow{MB}$ + $\overrightarrow{MC}$ = $\overrightarrow{MG}\$}
{$\\overrightarrow{MA}$ + $\overrightarrow{MB}$ + $\overrightarrow{MC}$ = 2 $\overrightarrow{MG}\$}
{\True $\\overrightarrow{MA}$ + $\overrightarrow{MB}$ + $\overrightarrow{MC}$ = 3 $\overrightarrow{MG}\$}
{$\\overrightarrow{MA}$ + $\overrightarrow{MB}$ + $\overrightarrow{MC}$ = 4 $\overrightarrow{MG}\$}
\loigiai{
• Chèn điểm $\,\text{G}\vào từng vectơ ở vế trái:\\overrightarrow{MA} = \overrightarrow{MG} + \overrightarrow{GA}\$; $\\overrightarrow{MB}$ = $\overrightarrow{MG}$ + $\overrightarrow{GB}\;\\overrightarrow{MC} = \overrightarrow{MG} + \overrightarrow{GC}\$.
 
• Cộng ba đẳng thức lại: $\\overrightarrow{MA}$ + $\overrightarrow{MB}$ + $\overrightarrow{MC}$ = 3 $\overrightarrow{MG}$ + $\overrightarrow{GA} + \overrightarrow{GB} + \overrightarrow{GC})$ \ $.
 
• Vì$\$G\$là trọng tâm tam giác$\$ABC$\$nên$\$\overrightarrow{GA} + \overrightarrow{GB} + \overrightarrow{GC} = \overrightarrow{0}\$.
 
• Vậy$overrightarrow{MA}$+$\overrightarrow{MB}$+$\overrightarrow{MC}$= 3$ \overrightarrow{MG}\.
}
\end{ex}

% ===== Câu 79 =====
% ID: LATEX_0DB1EA144BE3
% Mức: TH
\dangbt{Áp dụng quy tắc ba điểm và quy tắc hình bình hành}
\begin{ex}
% ID: LATEX_0DB1EA144BE3
% Mức: TH
Cho ba điểm phân biệt $\A, B, C\. Đẳng thức nào sau đây đúng?$
\choice
{$\\overrightarrow{AB}$ + $\overrightarrow{CB}$ = $\overrightarrow{AC}\$}
{$\\overrightarrow{CB}$ + $\overrightarrow{AC}$ = $\overrightarrow{AB}\$}
{$\\overrightarrow{AB}$ + $\overrightarrow{BC}$ = $\overrightarrow{CA}\$}
{\True $\\overrightarrow{AB}$ + $\overrightarrow{BC}$ = $\overrightarrow{AC}\$}
\loigiai{
Theo quy tắc ba điểm đối với phép cộng vectơ, với ba điểm $\A, B, C\bất kì ta luôn có\\overrightarrow{AB} + \overrightarrow{BC} = \overrightarrow{AC}\$.
}
\end{ex}

% ===== Câu 80 =====
% ID: LATEX_AF6B2F66DC9A
% Mức: TH
\begin{ex}
% ID: LATEX_AF6B2F66DC9A
% Mức: TH
Cho ba điểm phân biệt $\M, N, P\. Vectơ tổng\\overrightarrow{NP} + \overrightarrow{MN}\$ bằng vectơ nào sau đây?
\choice
{$\\overrightarrow{PN}\$}
{$\\overrightarrow{PM}\$}
{\True $\\overrightarrow{MP}\$}
{$\\overrightarrow{NM}\$}
\loigiai{
• Áp dụng tính chất giao hoán của phép cộng vectơ, ta có: $\\overrightarrow{NP}$ + $\overrightarrow{MN}$ = $\overrightarrow{MN}$ + $\overrightarrow{NP}\$.
 
• Theo quy tắc ba điểm, $\$ \overrightarrow{MN} + \overrightarrow{NP} = \overrightarrow{MP}\ $.
 
• Vậy$overrightarrow{NP}$+$\overrightarrow{MN}$=$ \overrightarrow{MP}\.
}
\end{ex}

% ===== Câu 81 =====
% ID: LATEX_B3959566D3C6
% Mức: TH
\begin{ex}
% ID: LATEX_B3959566D3C6
% Mức: TH
Cho hình bình hành $\$ ABCD $\. Đẳng thức nào sau đây đúng?$
\choice
{\True $\\overrightarrow{AB}$ + $\overrightarrow{AD}$ = $\overrightarrow{AC}\$}
{$\\overrightarrow{AB}$ + $\overrightarrow{AD}$ = $\overrightarrow{BD}\$}
{$\\overrightarrow{AB}$ + $\overrightarrow{AD}$ = $\overrightarrow{CD}\$}
{$\\overrightarrow{AB}$ + $\overrightarrow{AD}$ = $\overrightarrow{CB}\$}
\loigiai{
Theo quy tắc hình bình hành, vì $\$ ABCD $\là hình bình hành nên tổng hai vectơ cạnh chung gốc bằng vectơ đường chéo xuất phát từ gốc đó, tức là\\overrightarrow{AB} + \overrightarrow{AD} = \overrightarrow{AC}\$.
}
\end{ex}

% ===== Câu 82 =====
% ID: LATEX_1431DB984C9E
% Mức: TH
\begin{ex}
% ID: LATEX_1431DB984C9E
% Mức: TH
Cho bốn điểm phân biệt $\A, B, C, O\. Đẳng thức nào sau đây đúng?$
\choice
{$\\overrightarrow{AB}=\overrightarrow{OB}$ + $\overrightarrow{OA}\$}
{$\\overrightarrow{AB}=\overrightarrow{AC}$ + $\overrightarrow{BC}\$}
{\True $\\overrightarrow{OA}=\overrightarrow{CA}$ - $\overrightarrow{CO}\$}
{$\\overrightarrow{OA}=\overrightarrow{OB}$ - $\overrightarrow{BA}\$}
\loigiai{
• Áp dụng quy tắc trừ vectơ với ba điểm $\O, C, A\ta có:\\overrightarrow{CA} - \overrightarrow{CO} = \overrightarrow{OA}\$.
 
• Các phương án khác đều sai quy tắc cộng và trừ vectơ.
}
\end{ex}

% ===== Câu 83 =====
% ID: LATEX_A46C88793698
% Mức: TH
\begin{ex}
% ID: LATEX_A46C88793698
% Mức: TH
Cho bốn điểm phân biệt $\A, B, C, D\. Vectơ tổng\\overrightarrow{AB} + \overrightarrow{CD} + \overrightarrow{BC} + \overrightarrow{DA}\$ bằng vectơ nào sau đây?
\choice
{\True $\\overrightarrow{0}\$}
{$\\overrightarrow{AC}\$}
{$\\overrightarrow{BD}\$}
{$\\overrightarrow{BA}\$}
\loigiai{
• Ta sắp xếp lại các số hạng và áp dụng quy tắc ba điểm liên tiếp: $\\overrightarrow{AB}$ + $\overrightarrow{CD}$ + $\overrightarrow{BC}$ + $\overrightarrow{DA}$ = $\overrightarrow{AB}$ + $\overrightarrow{BC}$ + $\overrightarrow{CD}$ + $\overrightarrow{DA}\$.
 
• Thực hiện phép cộng từ trái sang phải: $\$ \overrightarrow{AC} + \overrightarrow{CD} + \overrightarrow{DA} = \overrightarrow{AD} + \overrightarrow{DA} = \overrightarrow{AA}\ $.
 
• Vì điểm đầu và điểm cuối trùng nhau nên$overrightarrow{AA}$=$ \overrightarrow{0}\.
}
\end{ex}

% ===== Câu 84 =====
% ID: LATEX_4302F8D17026
% Mức: TH
\dangbt{Rút gọn biểu thức và chứng minh đẳng thức vectơ}
\begin{ex}
% ID: LATEX_4302F8D17026
% Mức: TH
Cho bốn điểm $\A, B, C, D\bất kì. Khẳng định nào sau đây là đúng?$
\choice
{$\\overrightarrow{AB}$ - $\overrightarrow{AD}$ = $\overrightarrow{CD}$ + $\overrightarrow{BC}\$}
{$\\overrightarrow{AB}$ - $\overrightarrow{AC}$ = $\overrightarrow{DB}$ - $\overrightarrow{DC}\$}
{$\\overrightarrow{BC}$ + $\overrightarrow{AB}$ = $\overrightarrow{DA}$ - $\overrightarrow{DC}\$}
{\True $\\overrightarrow{AB}$ + $\overrightarrow{CD}$ = $\overrightarrow{AD}$ + $\overrightarrow{CB}\$}
\loigiai{
• Xét vế trái: $\\overrightarrow{AB} + \overrightarrow{CD}\. Áp dụng quy tắc chèn điểm\,\text{D}\$ vào vectơ thứ nhất và điểm $\B$ \ $vào vectơ thứ hai.
 
• Ta có$ \ $\overrightarrow{AB} + \overrightarrow{CD} = (\overrightarrow{AD} + \overrightarrow{DB}) + (\overrightarrow{CB} + \overrightarrow{BD})\$.
 
• Nhóm lại ta được: $\\overrightarrow{AD}$ + $\overrightarrow{CB}$ + $(\overrightarrow{DB} + \overrightarrow{BD})$ = $\overrightarrow{AD}$ + $\overrightarrow{CB}$ + $\overrightarrow{0}$ = $\overrightarrow{AD}$ + $\overrightarrow{CB}\. Đẳng thức đúng.$
}
\end{ex}

% ===== Câu 85 =====
% ID: LATEX_3BCD556B4A93
% Mức: TH
\begin{ex}
% ID: LATEX_3BCD556B4A93
% Mức: TH
Rút gọn biểu thức vectơ $\\overrightarrow{P} = \overrightarrow{AB} - \overrightarrow{AC} + \overrightarrow{BD}\, ta được kết quả nào sau đây?$
\choice
{$\\overrightarrow{AD}\$}
{$\\overrightarrow{BC}\$}
{\True $\\overrightarrow{CD}\$}
{$\\overrightarrow{DA}\$}
\loigiai{
• Áp dụng quy tắc hiệu cho hai vectơ đầu tiên: $\\overrightarrow{AB}$ - $\overrightarrow{AC}$ = $\overrightarrow{CB}\$.
 
• Thay vào biểu thức ta được: $\$ \overrightarrow{P} = \overrightarrow{CB} + \overrightarrow{BD}\ $.
 
• Áp dụng quy tắc ba điểm:$overrightarrow{CB}$+$\overrightarrow{BD}$=$ \overrightarrow{CD}\.
}
\end{ex}

% ===== Câu 86 =====
% ID: LATEX_0658205A4C5B
% Mức: TH
\begin{ex}
% ID: LATEX_0658205A4C5B
% Mức: TH
Cho hình bình hành $\$ ABCD $\. Rút gọn biểu thức\\overrightarrow{AB} - \overrightarrow{DA} - \overrightarrow{DC}\$ ta được vectơ nào?
\choice
{$\\overrightarrow{AC}\$}
{\True $\\overrightarrow{0}\$}
{$\\overrightarrow{BD}\$}
{$\$ 2 $\overrightarrow{AB}\$}
\loigiai{
• Vì $\$ ABCD $\là hình bình hành nên\\overrightarrow{DC} = \overrightarrow{AB}\$.
 
• Biểu thức trở thành: $\\overrightarrow{AB}$ - $\overrightarrow{DA}$ - $\overrightarrow{AB}$ = - $\overrightarrow{DA}$ = $\overrightarrow{AD}\$.
 
• Tuy nhiên, ta có thể làm cách khác: $\$ \overrightarrow{AB} - \overrightarrow{DA} - \overrightarrow{DC} = \overrightarrow{AB} + \overrightarrow{AD} - \overrightarrow{DC}\ $.
 
• Theo quy tắc hình bình hành$overrightarrow{AB}$+$\overrightarrow{AD}$=$ \overrightarrow{AC}\. Mà
overrightarrow{DC} = \overrightarrow{AB}\ $không giúp ra$ overrightarrow{0}\ $.
 
• Hãy thử lại:$\$\overrightarrow{AB} - \overrightarrow{DA} - \overrightarrow{DC} = \overrightarrow{AB} + \overrightarrow{AD} - \overrightarrow{DC} = \overrightarrow{AC} - \overrightarrow{DC} = \overrightarrow{AC} + \overrightarrow{CD} = \overrightarrow{AD}\$.
 
• Xin lỗi, có sự nhầm lẫn. Hãy làm lại: $\\overrightarrow{AB}$ - $\overrightarrow{DA}$ - $\overrightarrow{DC}$ = $\overrightarrow{AB}$ + $\overrightarrow{AD}$ - $\overrightarrow{DC}\. Mà\,\text{ABCD}\$ là hình bình hành nên $\\overrightarrow{AB}$ = $\overrightarrow{DC}\, do đó\\overrightarrow{AB} - \overrightarrow{DC} = \overrightarrow{0}\$.
 
• Vậy biểu thức còn lại là $\$ - $\overrightarrow{DA}$ = $\overrightarrow{AD}\$.
 
• Phát hiện phương án không khớp, sửa lại câu hỏi: Rút gọn $\$ \overrightarrow{AB} + \overrightarrow{AD} - \overrightarrow{AC}\.
}
\end{ex}

% ===== Câu 87 =====
% ID: LATEX_A9B2A6742B7D
% Mức: TH
\begin{ex}
% ID: LATEX_A9B2A6742B7D
% Mức: TH
Cho hình bình hành $\$ ABCD $\. Biểu thức\\overrightarrow{AB} + \overrightarrow{AD} - \overrightarrow{AC}\$ bằng vectơ nào sau đây?
\choice
{$\\overrightarrow{BD}\$}
{$\\overrightarrow{CB}\$}
{\True $\\overrightarrow{0}\$}
{$\\overrightarrow{DA}\$}
\loigiai{
• Áp dụng quy tắc hình bình hành, ta có tổng $\\overrightarrow{AB}$ + $\overrightarrow{AD}$ = $\overrightarrow{AC}\$.
 
• Thay vào biểu thức ta được: $\$ \overrightarrow{AC} - \overrightarrow{AC} = \overrightarrow{0}\.
}
\end{ex}

% ===== Câu 88 =====
% ID: LATEX_4FD940D16699
% Mức: TH
\begin{ex}
% ID: LATEX_4FD940D16699
% Mức: TH
Cho $\6\điểm\A, B, C, D, E, F\$ bất kì. Đẳng thức nào sau đây đúng?
\choice
{$\\overrightarrow{AB}$ + $\overrightarrow{CD}$ = $\overrightarrow{AD}$ + $\overrightarrow{BC}\$}
{$\\overrightarrow{AD}$ + $\overrightarrow{BE}$ + $\overrightarrow{CF}$ = $\overrightarrow{AE}$ + $\overrightarrow{BF}$ + $\overrightarrow{CE}\$}
{\True $\\overrightarrow{AB}$ - $\overrightarrow{AF}$ + $\overrightarrow{CD}$ - $\overrightarrow{CB}$ + $\overrightarrow{EF}$ - $\overrightarrow{ED}$ = $\overrightarrow{0}\$}
{$\\overrightarrow{AB}$ + $\overrightarrow{CD}$ + $\overrightarrow{EF}$ = $\overrightarrow{AF}$ + $\overrightarrow{ED}$ + $\overrightarrow{BC}\$}
\loigiai{
• Xét vế trái của phương án đúng: $\VT =(\overrightarrow{AB} - \overrightarrow{AF})$ + $(\overrightarrow{CD} - \overrightarrow{CB})$ + $(\overrightarrow{EF} - \overrightarrow{ED})$ \ $.
 
• Áp dụng quy tắc trừ vectơ:$ \ $VT = \overrightarrow{FB} + \overrightarrow{BD} + \overrightarrow{DF}\$.
 
• Áp dụng quy tắc cộng vectơ liên tiếp: $\$ VT = $\overrightarrow{FD}$ + $\overrightarrow{DF}$ = $\overrightarrow{FF}$ = $\overrightarrow{0}\.$
}
\end{ex}

% ===== Câu 89 =====
% ID: LATEX_495D545B9A85
% Mức: TH
\dangbt{Tính độ dài của tổng và hiệu hai vectơ}
\begin{ex}
% ID: LATEX_495D545B9A85
% Mức: TH
Cho tam giác $\,\text{ABC}\đều cạnh\a\$. Tính độ dài của vectơ $\\overrightarrow{AB}$ - $\overrightarrow{AC}\$.
\choice
{\True $\$ a $\$}
{$\$ a $\sqrt{3}\$}
{$\2\,\text{a}$ \}
{$\$ 0 $\$}
\loigiai{
• Áp dụng quy tắc hiệu hai vectơ chung gốc, ta có: $\\overrightarrow{AB}$ - $\overrightarrow{AC}$ = $\overrightarrow{CB}\$.
 
• Độ dài của vectơ này là: $\$ |\overrightarrow{AB} - \overrightarrow{AC}| = |\overrightarrow{CB}| = CB\ $.
 
• Vì tam giác$\$ABC$\đều cạnh\a\$nên$\$CB = a$ \.
}
\end{ex}

% ===== Câu 90 =====
% ID: LATEX_3704ABD1EBCC
% Mức: TH
\begin{ex}
% ID: LATEX_3704ABD1EBCC
% Mức: TH
Cho tam giác $\,\text{ABC}\đều cạnh\a\$. Tính độ dài của vectơ $\\overrightarrow{AB}$ + $\overrightarrow{AC}\$.
\choice
{$\$ a $\$}
{\True $\$ a $\sqrt{3}\$}
{$\2\,\text{a}$ \}
{$\$ a $\sqrt{2}\$}
\loigiai{
• Gọi $\,\text{M}\là trung điểm của cạnh\,\text{BC}\$. Theo tính chất trung điểm, ta có $\\overrightarrow{AB}$ + $\overrightarrow{AC}$ = 2 $\overrightarrow{AM}\$.
 
• Do đó, $\$ |\overrightarrow{AB} + \overrightarrow{AC}| = |2\overrightarrow{AM}| = 2AM\ $.
 
• Trong tam giác đều cạnh$\$a$\, đường trung tuyến đồng thời là đường cao nên\,\text{AM} = \dfrac{a\sqrt{3}}{2}\$.
 
• Vậy$\$|$\overrightarrow{AB}$+$\overrightarrow{AC}| = 2\cdot\dfrac{a\sqrt{3}}{2}$= a$ \sqrt{3}\.
}
\end{ex}

% ===== Câu 91 =====
% ID: LATEX_BD8F813D04FB
% Mức: TH
\begin{ex}
% ID: LATEX_BD8F813D04FB
% Mức: TH
Cho hình vuông $\$ ABCD $\có cạnh bằng\a\$. Tính độ dài của vectơ $\\overrightarrow{AB}$ + $\overrightarrow{AD}\$.
\choice
{$\2\,\text{a}$ \}
{$\$ a $\$}
{\True $\$ a $\sqrt{2}\$}
{$\\dfrac{a\sqrt{2}}{2}\$}
\loigiai{
• Vì $\$ ABCD $\là hình vuông nên nó cũng là hình bình hành. Áp dụng quy tắc hình bình hành ta có:\\overrightarrow{AB} + \overrightarrow{AD} = \overrightarrow{AC}\$.
 
• Độ dài vectơ cần tìm là $\$ | $\overrightarrow{AB}$ + $\overrightarrow{AD}| = |\overrightarrow{AC}| = AC\$.
 
• Trong hình vuông cạnh $\$ a\ $, độ dài đường chéo là$ \ $AC =$ \sqrt{AB^2 + BC^2} $=$ \sqrt{a^2 + a^2} $= a$ \sqrt{2}\.
}
\end{ex}

% ===== Câu 92 =====
% ID: LATEX_331847432B03
% Mức: TH
\begin{ex}
% ID: LATEX_331847432B03
% Mức: TH
Cho hình vuông $\$ ABCD $\tâm\,\text{O}\$, cạnh $\$ a $\. Tính độ dài của vectơ\\overrightarrow{OA} - \overrightarrow{CB}\$.
\choice
{\True $\\dfrac{a\sqrt{2}}{2}\$}
{$\$ a $\sqrt{2}\$}
{$\$ a $\$}
{$\\dfrac{a}{2}\$}
\loigiai{
• Vì $\$ ABCD $\là hình vuông nên\\overrightarrow{CB} = \overrightarrow{DA}\$.
 
• Thay vào biểu thức ta có: $\\overrightarrow{OA}$ - $\overrightarrow{CB}$ = $\overrightarrow{OA}$ - $\overrightarrow{DA}$ = $\overrightarrow{OA}$ + $\overrightarrow{AD}\$.
 
• Áp dụng quy tắc ba điểm: $\$ \overrightarrow{OA} + \overrightarrow{AD} = \overrightarrow{OD}\ $.
 
• Độ dài vectơ là$\$|$\overrightarrow{OD}| = OD\. Vì\,\text{O}\$là tâm hình vuông nên$\$OD =$\dfrac{1}{2}BD$=$ \dfrac{a\sqrt{2}}{2}\.
}
\end{ex}

% ===== Câu 93 =====
% ID: LATEX_9E693EBEFA4D
% Mức: TH
\begin{ex}
% ID: LATEX_9E693EBEFA4D
% Mức: TH
Cho hình chữ nhật $\$ ABCD $\có\AB = 3\$, $\$ BC = 4 $\. Tính độ dài của vectơ\\overrightarrow{AB} + \overrightarrow{BC}\$.
\choice
{$\$ 7 $\$}
{\True $\$ 5 $\$}
{$\$ 1 $\$}
{$\$ 12 $\$}
\loigiai{
• Theo quy tắc ba điểm, ta có $\\overrightarrow{AB}$ + $\overrightarrow{BC}$ = $\overrightarrow{AC}\$.
 
• Độ dài của vectơ là $\$ |\overrightarrow{AC}| = AC\ $.
 
• Xét tam giác$\$ABC$\vuông tại\,\text{B}\$, theo định lí Pythagore:$\$AC =$\sqrt{AB^2 + BC^2}$=$\sqrt{3^2 + 4^2}$=$\sqrt{25}$= 5$ \.
}
\end{ex}

% ===== Câu 94 =====
% ID: LATEX_6F0ACD0D2B2C
% Mức: TH
\dangbt{Xác định điểm thỏa mãn đẳng thức vectơ}
\begin{ex}
% ID: LATEX_6F0ACD0D2B2C
% Mức: TH
Cho đoạn thẳng $\,\text{AB}\và điểm\,\text{M}\$ thỏa mãn đẳng thức $\\overrightarrow{MA}$ + $\overrightarrow{MB}$ = $\overrightarrow{0}\. Khẳng định nào sau đây đúng?$
\choice
{\True $\M\là trung điểm của đoạn thẳng\AB\$}
{$\M\trùng với điểm\A\$}
{$\M\trùng với điểm\B\$}
{$\M\nằm ngoài đoạn thẳng\AB\$}
\loigiai{
• Ta có đẳng thức $\\overrightarrow{MA}$ + $\overrightarrow{MB}$ = $\overrightarrow{0}\Leftrightarrow\overrightarrow{MA}$ = - $\overrightarrow{MB}\$.
 
• Điều này có nghĩa là hai vectơ $\$ \overrightarrow{MA}\ $và$ \ $\overrightarrow{MB}\$ là hai vectơ đối nhau, tức là chúng ngược hướng và có độ dài bằng nhau ($\$ MA = MB\ $).
 
• Do đó, ba điểm$\,\text{M}$,$A$,$B$\thẳng hàng,\M\$nằm giữa$\A$,$B$\và cách đều\,\text{A}, B\$. Vậy$\,\text{M}$\là trung điểm của\AB\$.
}
\end{ex}

% ===== Câu 95 =====
% ID: LATEX_40087EE97ADD
% Mức: TH
\begin{ex}
% ID: LATEX_40087EE97ADD
% Mức: TH
Cho ba điểm $\,\text{A}, B, C\phân biệt. Tìm điểm\,\text{M}\$ thỏa mãn $\\overrightarrow{MA}$ - $\overrightarrow{MB}$ + $\overrightarrow{MC}$ = $\overrightarrow{0}\$.
\choice
{$\M\là trung điểm của\AC\$}
{\True $\M\là đỉnh thứ tư của hình bình hành\ABCM\$}
{$\M\là đỉnh thứ tư của hình bình hành\ABMC\$}
{$\M\là trọng tâm tam giác\ABC\$}
\loigiai{
• Áp dụng quy tắc hiệu cho hai vectơ đầu tiên: $\\overrightarrow{MA}$ - $\overrightarrow{MB}$ = $\overrightarrow{BA}\$.
 
• Đẳng thức đã cho trở thành: $\$ \overrightarrow{BA} + \overrightarrow{MC} = \overrightarrow{0} \Leftrightarrow \overrightarrow{MC} = -\overrightarrow{BA} \Leftrightarrow \overrightarrow{MC} = \overrightarrow{AB}\ $.
 
• Điều kiện$overrightarrow{MC}$=$\overrightarrow{AB}\chứng tỏ tứ giác\,\text{ABCM}\$là hình bình hành (với thứ tự các đỉnh liên tiếp là$\,\text{A}$,$B$,$C$,$M$\).$
}
\end{ex}

% ===== Câu 96 =====
% ID: LATEX_48C88D4FB932
% Mức: TH
\begin{ex}
% ID: LATEX_48C88D4FB932
% Mức: TH
Cho tam giác $\,\text{ABC}\. Gọi\,\text{M}\$ là điểm xác định bởi đẳng thức $\\overrightarrow{AM}$ = $\overrightarrow{AB}$ + $\overrightarrow{AC}\. Khẳng định nào sau đây là đúng?$
\choice
{$\M\là trung điểm của\BC\$}
{$\M\là trọng tâm tam giác\ABC\$}
{\True $\$ ABMC $\là hình bình hành$}
{$\$ AMBC $\là hình bình hành$}
\loigiai{
• Theo quy tắc hình bình hành, tổng hai vectơ chung gốc $\\overrightarrow{AB} + \overrightarrow{AC}\bằng vectơ đường chéo\\overrightarrow{AM}\$ của hình bình hành có ba đỉnh là $\,\text{A}$, $B$, $C$ \ $.
 
• Do đó, tứ giác tạo bởi hai cạnh$\$AB\$và$\$AC$\$cùng với đường chéo$\$AM\$là hình bình hành$\$ABMC$\$.
 
• Chú ý thứ tự đỉnh: \$A\$nối với$\$B$\$và$\$C\$,$\$M$\$đối diện$\$A\$, nên tứ giác đọc theo vòng tròn là$\$ABMC$ \.
}
\end{ex}

% ===== Câu 97 =====
% ID: LATEX_E0B4D02AEB3B
% Mức: TH
\begin{ex}
% ID: LATEX_E0B4D02AEB3B
% Mức: TH
Cho bốn điểm phân biệt $\,\text{A}, B, C, D\. Gọi\,\text{M}\$ là điểm thỏa mãn $\\overrightarrow{AM}$ = $\overrightarrow{AB}$ - $\overrightarrow{AC}\. Khẳng định nào sau đây đúng?$
\choice
{\True $\\overrightarrow{AM}=\overrightarrow{CB}\$}
{$\\overrightarrow{AM}=\overrightarrow{BC}\$}
{$\\overrightarrow{AM}=\overrightarrow{CD}\$}
{$\\overrightarrow{AM}=\overrightarrow{AD}\$}
\loigiai{
• Áp dụng quy tắc trừ hai vectơ chung gốc: $\\overrightarrow{AB}$ - $\overrightarrow{AC}$ = $\overrightarrow{CB}\$.
 
• Vậy đẳng thức đã cho trở thành $\$ \overrightarrow{AM} = \overrightarrow{CB}\.
}
\end{ex}

% ===== Câu 98 =====
% ID: LATEX_D8CCCEF6C78E
% Mức: TH
\begin{ex}
% ID: LATEX_D8CCCEF6C78E
% Mức: TH
Cho tam giác $\,\text{ABC}\. Tập hợp các điểm\,\text{M}\$ thỏa mãn $\$ | $\overrightarrow{MA}$ + $\overrightarrow{MB}| = |\overrightarrow{MC}$ + $\overrightarrow{MB}|\là:$
\choice
{Đường tròn tâm $\B\bán kính\AC\$}
{Đường tròn đường kính $\$ AC $\$}
{\True Đường trung trực của đoạn thẳng $\$ AC $\$}
{Đường trung trực của đoạn thẳng $\$ BC $\$}
\loigiai{
• Xét vế trái: Gọi $\,\text{I}\là trung điểm của\,\text{AB}\$, ta có $\\overrightarrow{MA}$ + $\overrightarrow{MB}$ = 2 $\overrightarrow{MI}\Rightarrow$ | $\overrightarrow{MA}$ + $\overrightarrow{MB}| = 2MI\$.
 
• Xét vế phải: Gọi $\$ J\ $là trung điểm của$ \ $BC$ \ $, ta có$ \ $\overrightarrow{MC} + \overrightarrow{MB} = 2\overrightarrow{MJ} \Rightarrow |\overrightarrow{MC} + \overrightarrow{MB}| = 2MJ\$.
 
• Đẳng thức trở thành $\$ 2MI = 2MJ \Leftrightarrow MI = MJ $\$.
 
• Do đó, tập hợp các điểm $\$ M\ $là đường trung trực của đoạn thẳng$ \ $IJ$ \ $.
 
• Mặt khác, \$IJ\$là đường trung bình của tam giác$\$ABC$\$nên$\$IJ \parallel AC\$. Đường trung trực của$\$IJ$\$cũng chính là đường trung trực của đoạn$\$AC\$(nếu xét tam giác cân) - khoảng cách. Tuy nhiên, cách giải chuẩn hơn mà không cần gọi trung điểm là:
 
• Cách 2:$\$|$\overrightarrow{MA}$+$\overrightarrow{MB}| = |\overrightarrow{MB}$+$\overrightarrow{MC}|\. Đây là độ dài của tổng, không dễ khử điểm\,\text{M}\$.
 
• Chờ đã,$overrightarrow{MB}$+$\overrightarrow{MC}\chưa triệt tiêu\,\text{M}\$. Phải dùng trung điểm$\,\text{I}$,$J$\. Khi\MI = MJ\$,$\M$\nằm trên trung trực của\,\text{IJ}\$. Đường trung trực của$\$IJ$\không nhất thiết là đường trung trực của\,\text{AC}\$.
 
• Hãy đổi câu hỏi thành một bài toán quen thuộc hơn:$\$|$\overrightarrow{MB}$-$\overrightarrow{MC}| = |\overrightarrow{BM}$-$\overrightarrow{BA}|\$.
 
• Sửa lại câu hỏi: Cho tam giác$\$ABC\$. Điểm$\$M$\$thỏa mãn$\$ \overrightarrow{MA} + \overrightarrow{MB} + \overrightarrow{MC} = \overrightarrow{0}\.
}
\end{ex}

% ===== Câu 99 =====
% ID: LATEX_29094683BBF1
% Mức: TH
\dangbt{Tính chất trung điểm và trọng tâm trong tổng vectơ}
\begin{ex}
% ID: LATEX_29094683BBF1
% Mức: TH
Cho tam giác $\ABC\có trọng tâm\G\$. Khẳng định nào sau đây là đúng?
\choice
{$\\overrightarrow{GA} + \overrightarrow{GB} + \overrightarrow{GC} = 3\overrightarrow{GM}\với\,\text{M}\$ là trung điểm $\$ BC $\$}
{$\\overrightarrow{AB}$ + $\overrightarrow{AC}$ = $\overrightarrow{AG}\$}
{\True $\\overrightarrow{GA}$ + $\overrightarrow{GB}$ + $\overrightarrow{GC}$ = $\overrightarrow{0}\$}
{$\\overrightarrow{GA}$ + $\overrightarrow{GB}$ = $\overrightarrow{GC}\$}
\loigiai{
• Theo tính chất trọng tâm của tam giác, tổng các vectơ hướng từ trọng tâm đến ba đỉnh luôn bằng vectơ-không.
 
• Vậy $\\overrightarrow{GA}$ + $\overrightarrow{GB}$ + $\overrightarrow{GC}$ = $\overrightarrow{0}\.$
}
\end{ex}

% ===== Câu 100 =====
% ID: LATEX_1BA9FFC3EF07
% Mức: TH
\begin{ex}
% ID: LATEX_1BA9FFC3EF07
% Mức: TH
Cho hình bình hành $\$ ABCD $\tâm\O\$. Đẳng thức nào sau đây là đúng?
\choice
{\True $\\overrightarrow{OA}$ + $\overrightarrow{OB}$ + $\overrightarrow{OC}$ + $\overrightarrow{OD}$ = $\overrightarrow{0}\$}
{$\\overrightarrow{OA}$ + $\overrightarrow{OB}$ + $\overrightarrow{OC}$ + $\overrightarrow{OD}$ = 2 $\overrightarrow{AB}\$}
{$\\overrightarrow{OA}$ + $\overrightarrow{OB}$ + $\overrightarrow{OC}$ + $\overrightarrow{OD}$ = $\overrightarrow{AC}\$}
{$\\overrightarrow{OA}$ + $\overrightarrow{OB}$ + $\overrightarrow{OC}$ + $\overrightarrow{OD}$ = 4 $\overrightarrow{O}\$}
\loigiai{
• Vì $\,\text{O}\là tâm hình bình hành\,\text{ABCD}\$ nên $\,\text{O}$ \là trung điểm của đường chéo\AC\ $và đường chéo$ \ $BD$ \ $.
 
• Áp dụng tính chất trung điểm:$\$\overrightarrow{OA} + \overrightarrow{OC} = \overrightarrow{0}\$và$\$\overrightarrow{OB}$+$\overrightarrow{OD}$=$\overrightarrow{0}\$.
 
• Cộng vế theo vế ta được: \$ \overrightarrow{OA} + \overrightarrow{OB} + \overrightarrow{OC} + \overrightarrow{OD} = \overrightarrow{0} + \overrightarrow{0} = \overrightarrow{0}\.
}
\end{ex}

% ===== Câu 101 =====
% ID: LATEX_C2ED8EBD2F28
% Mức: TH
\begin{ex}
% ID: LATEX_C2ED8EBD2F28
% Mức: TH
Cho đoạn thẳng $\,\text{AB}\có trung điểm là\,\text{I}\$. Với điểm $\,\text{M}$ \tùy ý, khẳng định nào sau đây là đúng?
\choice
{$\\overrightarrow{MA}$ + $\overrightarrow{MB}$ = $\overrightarrow{MI}\$}
{\True $\\overrightarrow{MA}$ + $\overrightarrow{MB}$ = 2 $\overrightarrow{MI}\$}
{$\\overrightarrow{MA}$ + $\overrightarrow{MB}$ = 3 $\overrightarrow{MI}\$}
{$\\overrightarrow{MA}$ + $\overrightarrow{MB}$ = $\overrightarrow{IM}\$}
\loigiai{
• Chèn điểm $\,\text{I}\vào hai vectơ ở vế trái:\\overrightarrow{MA} = \overrightarrow{MI} + \overrightarrow{IA}\$ và $\\overrightarrow{MB}$ = $\overrightarrow{MI}$ + $\overrightarrow{IB}\$.
 
• Cộng lại ta có: $\$ \overrightarrow{MA} + \overrightarrow{MB} = 2\overrightarrow{MI} + (\overrightarrow{IA} + \overrightarrow{IB})\ $.
 
• Vì$\,\text{I}$\là trung điểm của\AB\$nên$\\overrightarrow{IA}$+$\overrightarrow{IB}$=$\overrightarrow{0}\$.
 
• Vậy $\$ \overrightarrow{MA} + \overrightarrow{MB} = 2\overrightarrow{MI}\.
}
\end{ex}

% ===== Câu 102 =====
% ID: LATEX_BF46C9E9C4F6
% Mức: TH
\begin{ex}
% ID: LATEX_BF46C9E9C4F6
% Mức: TH
Cho tam giác $\ABC\. Gọi\M, N, P\$ lần lượt là trung điểm của các cạnh $\$ BC, CA, AB $\. Tổng\\overrightarrow{AM} + \overrightarrow{BN} + \overrightarrow{CP}\$ bằng:
\choice
{$\\overrightarrow{AB}\$}
{$\\overrightarrow{BC}\$}
{\True $\\overrightarrow{0}\$}
{$\\overrightarrow{CA}\$}
\loigiai{
• Theo tính chất trung điểm, ta có $\2\overrightarrow{AM} = \overrightarrow{AB} + \overrightarrow{AC}\;\2\overrightarrow{BN} = \overrightarrow{BA} + \overrightarrow{BC}\$; $\$ 2 $\overrightarrow{CP}$ = $\overrightarrow{CA}$ + $\overrightarrow{CB}\$.
 
• Cộng ba đẳng thức trên: $\$ 2(\overrightarrow{AM} + \overrightarrow{BN} + \overrightarrow{CP}) = (\overrightarrow{AB} + \overrightarrow{BA}) + (\overrightarrow{AC} + \overrightarrow{CA}) + (\overrightarrow{BC} + \overrightarrow{CB})\ $.
 
•$\$2$\overrightarrow{AM} + \overrightarrow{BN} + \overrightarrow{CP})$=$\overrightarrow{0}$+$\overrightarrow{0}$+$\overrightarrow{0}$=$\overrightarrow{0}\$.
 
• Vậy$\$\overrightarrow{AM} + \overrightarrow{BN} + \overrightarrow{CP} = \overrightarrow{0}\$.
}
\end{ex}

% ===== Câu 103 =====
% ID: LATEX_AE13D1889B7A
% Mức: TH
\begin{ex}
% ID: LATEX_AE13D1889B7A
% Mức: TH
Cho tam giác $\,\text{ABC}\có trọng tâm\,\text{G}\$. Với điểm $\,\text{M}$ \bất kì, đẳng thức nào sau đây đúng?
\choice
{$\\overrightarrow{MA}$ + $\overrightarrow{MB}$ + $\overrightarrow{MC}$ = $\overrightarrow{MG}\$}
{$\\overrightarrow{MA}$ + $\overrightarrow{MB}$ + $\overrightarrow{MC}$ = 2 $\overrightarrow{MG}\$}
{\True $\\overrightarrow{MA}$ + $\overrightarrow{MB}$ + $\overrightarrow{MC}$ = 3 $\overrightarrow{MG}\$}
{$\\overrightarrow{MA}$ + $\overrightarrow{MB}$ + $\overrightarrow{MC}$ = 4 $\overrightarrow{MG}\$}
\loigiai{
• Chèn điểm $\,\text{G}\vào từng vectơ ở vế trái:\\overrightarrow{MA} = \overrightarrow{MG} + \overrightarrow{GA}\$; $\\overrightarrow{MB}$ = $\overrightarrow{MG}$ + $\overrightarrow{GB}\;\\overrightarrow{MC} = \overrightarrow{MG} + \overrightarrow{GC}\$.
 
• Cộng ba đẳng thức lại: $\\overrightarrow{MA}$ + $\overrightarrow{MB}$ + $\overrightarrow{MC}$ = 3 $\overrightarrow{MG}$ + $\overrightarrow{GA} + \overrightarrow{GB} + \overrightarrow{GC})$ \ $.
 
• Vì$\$G\$là trọng tâm tam giác$\$ABC$\$nên$\$\overrightarrow{GA} + \overrightarrow{GB} + \overrightarrow{GC} = \overrightarrow{0}\$.
 
• Vậy$overrightarrow{MA}$+$\overrightarrow{MB}$+$\overrightarrow{MC}$= 3$ \overrightarrow{MG}\.
}
\end{ex}

% ===== Câu 104 =====
% ID: LATEX_0DB1EA144BE3
% Mức: TH
\dangbt{Áp dụng quy tắc ba điểm và quy tắc hình bình hành}
\begin{ex}
% ID: LATEX_0DB1EA144BE3
% Mức: TH
Cho ba điểm phân biệt $\A, B, C\. Đẳng thức nào sau đây đúng?$
\choice
{$\\overrightarrow{AB}$ + $\overrightarrow{CB}$ = $\overrightarrow{AC}\$}
{$\\overrightarrow{CB}$ + $\overrightarrow{AC}$ = $\overrightarrow{AB}\$}
{$\\overrightarrow{AB}$ + $\overrightarrow{BC}$ = $\overrightarrow{CA}\$}
{\True $\\overrightarrow{AB}$ + $\overrightarrow{BC}$ = $\overrightarrow{AC}\$}
\loigiai{
Theo quy tắc ba điểm đối với phép cộng vectơ, với ba điểm $\A, B, C\bất kì ta luôn có\\overrightarrow{AB} + \overrightarrow{BC} = \overrightarrow{AC}\$.
}
\end{ex}

% ===== Câu 105 =====
% ID: LATEX_AF6B2F66DC9A
% Mức: TH
\begin{ex}
% ID: LATEX_AF6B2F66DC9A
% Mức: TH
Cho ba điểm phân biệt $\M, N, P\. Vectơ tổng\\overrightarrow{NP} + \overrightarrow{MN}\$ bằng vectơ nào sau đây?
\choice
{$\\overrightarrow{PN}\$}
{$\\overrightarrow{PM}\$}
{\True $\\overrightarrow{MP}\$}
{$\\overrightarrow{NM}\$}
\loigiai{
• Áp dụng tính chất giao hoán của phép cộng vectơ, ta có: $\\overrightarrow{NP}$ + $\overrightarrow{MN}$ = $\overrightarrow{MN}$ + $\overrightarrow{NP}\$.
 
• Theo quy tắc ba điểm, $\$ \overrightarrow{MN} + \overrightarrow{NP} = \overrightarrow{MP}\ $.
 
• Vậy$overrightarrow{NP}$+$\overrightarrow{MN}$=$ \overrightarrow{MP}\.
}
\end{ex}

% ===== Câu 106 =====
% ID: LATEX_B3959566D3C6
% Mức: TH
\begin{ex}
% ID: LATEX_B3959566D3C6
% Mức: TH
Cho hình bình hành $\$ ABCD $\. Đẳng thức nào sau đây đúng?$
\choice
{\True $\\overrightarrow{AB}$ + $\overrightarrow{AD}$ = $\overrightarrow{AC}\$}
{$\\overrightarrow{AB}$ + $\overrightarrow{AD}$ = $\overrightarrow{BD}\$}
{$\\overrightarrow{AB}$ + $\overrightarrow{AD}$ = $\overrightarrow{CD}\$}
{$\\overrightarrow{AB}$ + $\overrightarrow{AD}$ = $\overrightarrow{CB}\$}
\loigiai{
Theo quy tắc hình bình hành, vì $\$ ABCD $\là hình bình hành nên tổng hai vectơ cạnh chung gốc bằng vectơ đường chéo xuất phát từ gốc đó, tức là\\overrightarrow{AB} + \overrightarrow{AD} = \overrightarrow{AC}\$.
}
\end{ex}

% ===== Câu 107 =====
% ID: LATEX_1431DB984C9E
% Mức: TH
\begin{ex}
% ID: LATEX_1431DB984C9E
% Mức: TH
Cho bốn điểm phân biệt $\A, B, C, O\. Đẳng thức nào sau đây đúng?$
\choice
{$\\overrightarrow{AB}=\overrightarrow{OB}$ + $\overrightarrow{OA}\$}
{$\\overrightarrow{AB}=\overrightarrow{AC}$ + $\overrightarrow{BC}\$}
{\True $\\overrightarrow{OA}=\overrightarrow{CA}$ - $\overrightarrow{CO}\$}
{$\\overrightarrow{OA}=\overrightarrow{OB}$ - $\overrightarrow{BA}\$}
\loigiai{
• Áp dụng quy tắc trừ vectơ với ba điểm $\O, C, A\ta có:\\overrightarrow{CA} - \overrightarrow{CO} = \overrightarrow{OA}\$.
 
• Các phương án khác đều sai quy tắc cộng và trừ vectơ.
}
\end{ex}

% ===== Câu 108 =====
% ID: LATEX_A46C88793698
% Mức: TH
\begin{ex}
% ID: LATEX_A46C88793698
% Mức: TH
Cho bốn điểm phân biệt $\A, B, C, D\. Vectơ tổng\\overrightarrow{AB} + \overrightarrow{CD} + \overrightarrow{BC} + \overrightarrow{DA}\$ bằng vectơ nào sau đây?
\choice
{\True $\\overrightarrow{0}\$}
{$\\overrightarrow{AC}\$}
{$\\overrightarrow{BD}\$}
{$\\overrightarrow{BA}\$}
\loigiai{
• Ta sắp xếp lại các số hạng và áp dụng quy tắc ba điểm liên tiếp: $\\overrightarrow{AB}$ + $\overrightarrow{CD}$ + $\overrightarrow{BC}$ + $\overrightarrow{DA}$ = $\overrightarrow{AB}$ + $\overrightarrow{BC}$ + $\overrightarrow{CD}$ + $\overrightarrow{DA}\$.
 
• Thực hiện phép cộng từ trái sang phải: $\$ \overrightarrow{AC} + \overrightarrow{CD} + \overrightarrow{DA} = \overrightarrow{AD} + \overrightarrow{DA} = \overrightarrow{AA}\ $.
 
• Vì điểm đầu và điểm cuối trùng nhau nên$overrightarrow{AA}$=$ \overrightarrow{0}\.
}
\end{ex}

% ===== Câu 109 =====
% ID: LATEX_4302F8D17026
% Mức: TH
\dangbt{Rút gọn biểu thức và chứng minh đẳng thức vectơ}
\begin{ex}
% ID: LATEX_4302F8D17026
% Mức: TH
Cho bốn điểm $\A, B, C, D\bất kì. Khẳng định nào sau đây là đúng?$
\choice
{$\\overrightarrow{AB}$ - $\overrightarrow{AD}$ = $\overrightarrow{CD}$ + $\overrightarrow{BC}\$}
{$\\overrightarrow{AB}$ - $\overrightarrow{AC}$ = $\overrightarrow{DB}$ - $\overrightarrow{DC}\$}
{$\\overrightarrow{BC}$ + $\overrightarrow{AB}$ = $\overrightarrow{DA}$ - $\overrightarrow{DC}\$}
{\True $\\overrightarrow{AB}$ + $\overrightarrow{CD}$ = $\overrightarrow{AD}$ + $\overrightarrow{CB}\$}
\loigiai{
• Xét vế trái: $\\overrightarrow{AB} + \overrightarrow{CD}\. Áp dụng quy tắc chèn điểm\,\text{D}\$ vào vectơ thứ nhất và điểm $\B$ \ $vào vectơ thứ hai.
 
• Ta có$ \ $\overrightarrow{AB} + \overrightarrow{CD} = (\overrightarrow{AD} + \overrightarrow{DB}) + (\overrightarrow{CB} + \overrightarrow{BD})\$.
 
• Nhóm lại ta được: $\\overrightarrow{AD}$ + $\overrightarrow{CB}$ + $(\overrightarrow{DB} + \overrightarrow{BD})$ = $\overrightarrow{AD}$ + $\overrightarrow{CB}$ + $\overrightarrow{0}$ = $\overrightarrow{AD}$ + $\overrightarrow{CB}\. Đẳng thức đúng.$
}
\end{ex}

% ===== Câu 110 =====
% ID: LATEX_3BCD556B4A93
% Mức: TH
\begin{ex}
% ID: LATEX_3BCD556B4A93
% Mức: TH
Rút gọn biểu thức vectơ $\\overrightarrow{P} = \overrightarrow{AB} - \overrightarrow{AC} + \overrightarrow{BD}\, ta được kết quả nào sau đây?$
\choice
{$\\overrightarrow{AD}\$}
{$\\overrightarrow{BC}\$}
{\True $\\overrightarrow{CD}\$}
{$\\overrightarrow{DA}\$}
\loigiai{
• Áp dụng quy tắc hiệu cho hai vectơ đầu tiên: $\\overrightarrow{AB}$ - $\overrightarrow{AC}$ = $\overrightarrow{CB}\$.
 
• Thay vào biểu thức ta được: $\$ \overrightarrow{P} = \overrightarrow{CB} + \overrightarrow{BD}\ $.
 
• Áp dụng quy tắc ba điểm:$overrightarrow{CB}$+$\overrightarrow{BD}$=$ \overrightarrow{CD}\.
}
\end{ex}

% ===== Câu 111 =====
% ID: LATEX_0658205A4C5B
% Mức: TH
\begin{ex}
% ID: LATEX_0658205A4C5B
% Mức: TH
Cho hình bình hành $\$ ABCD $\. Rút gọn biểu thức\\overrightarrow{AB} - \overrightarrow{DA} - \overrightarrow{DC}\$ ta được vectơ nào?
\choice
{$\\overrightarrow{AC}\$}
{\True $\\overrightarrow{0}\$}
{$\\overrightarrow{BD}\$}
{$\$ 2 $\overrightarrow{AB}\$}
\loigiai{
• Vì $\$ ABCD $\là hình bình hành nên\\overrightarrow{DC} = \overrightarrow{AB}\$.
 
• Biểu thức trở thành: $\\overrightarrow{AB}$ - $\overrightarrow{DA}$ - $\overrightarrow{AB}$ = - $\overrightarrow{DA}$ = $\overrightarrow{AD}\$.
 
• Tuy nhiên, ta có thể làm cách khác: $\$ \overrightarrow{AB} - \overrightarrow{DA} - \overrightarrow{DC} = \overrightarrow{AB} + \overrightarrow{AD} - \overrightarrow{DC}\ $.
 
• Theo quy tắc hình bình hành$overrightarrow{AB}$+$\overrightarrow{AD}$=$ \overrightarrow{AC}\. Mà
overrightarrow{DC} = \overrightarrow{AB}\ $không giúp ra$ overrightarrow{0}\ $.
 
• Hãy thử lại:$\$\overrightarrow{AB} - \overrightarrow{DA} - \overrightarrow{DC} = \overrightarrow{AB} + \overrightarrow{AD} - \overrightarrow{DC} = \overrightarrow{AC} - \overrightarrow{DC} = \overrightarrow{AC} + \overrightarrow{CD} = \overrightarrow{AD}\$.
 
• Xin lỗi, có sự nhầm lẫn. Hãy làm lại: $\\overrightarrow{AB}$ - $\overrightarrow{DA}$ - $\overrightarrow{DC}$ = $\overrightarrow{AB}$ + $\overrightarrow{AD}$ - $\overrightarrow{DC}\. Mà\,\text{ABCD}\$ là hình bình hành nên $\\overrightarrow{AB}$ = $\overrightarrow{DC}\, do đó\\overrightarrow{AB} - \overrightarrow{DC} = \overrightarrow{0}\$.
 
• Vậy biểu thức còn lại là $\$ - $\overrightarrow{DA}$ = $\overrightarrow{AD}\$.
 
• Phát hiện phương án không khớp, sửa lại câu hỏi: Rút gọn $\$ \overrightarrow{AB} + \overrightarrow{AD} - \overrightarrow{AC}\.
}
\end{ex}

% ===== Câu 112 =====
% ID: LATEX_A9B2A6742B7D
% Mức: TH
\begin{ex}
% ID: LATEX_A9B2A6742B7D
% Mức: TH
Cho hình bình hành $\$ ABCD $\. Biểu thức\\overrightarrow{AB} + \overrightarrow{AD} - \overrightarrow{AC}\$ bằng vectơ nào sau đây?
\choice
{$\\overrightarrow{BD}\$}
{$\\overrightarrow{CB}\$}
{\True $\\overrightarrow{0}\$}
{$\\overrightarrow{DA}\$}
\loigiai{
• Áp dụng quy tắc hình bình hành, ta có tổng $\\overrightarrow{AB}$ + $\overrightarrow{AD}$ = $\overrightarrow{AC}\$.
 
• Thay vào biểu thức ta được: $\$ \overrightarrow{AC} - \overrightarrow{AC} = \overrightarrow{0}\.
}
\end{ex}

% ===== Câu 113 =====
% ID: LATEX_4FD940D16699
% Mức: TH
\begin{ex}
% ID: LATEX_4FD940D16699
% Mức: TH
Cho $\6\điểm\A, B, C, D, E, F\$ bất kì. Đẳng thức nào sau đây đúng?
\choice
{$\\overrightarrow{AB}$ + $\overrightarrow{CD}$ = $\overrightarrow{AD}$ + $\overrightarrow{BC}\$}
{$\\overrightarrow{AD}$ + $\overrightarrow{BE}$ + $\overrightarrow{CF}$ = $\overrightarrow{AE}$ + $\overrightarrow{BF}$ + $\overrightarrow{CE}\$}
{\True $\\overrightarrow{AB}$ - $\overrightarrow{AF}$ + $\overrightarrow{CD}$ - $\overrightarrow{CB}$ + $\overrightarrow{EF}$ - $\overrightarrow{ED}$ = $\overrightarrow{0}\$}
{$\\overrightarrow{AB}$ + $\overrightarrow{CD}$ + $\overrightarrow{EF}$ = $\overrightarrow{AF}$ + $\overrightarrow{ED}$ + $\overrightarrow{BC}\$}
\loigiai{
• Xét vế trái của phương án đúng: $\VT =(\overrightarrow{AB} - \overrightarrow{AF})$ + $(\overrightarrow{CD} - \overrightarrow{CB})$ + $(\overrightarrow{EF} - \overrightarrow{ED})$ \ $.
 
• Áp dụng quy tắc trừ vectơ:$ \ $VT = \overrightarrow{FB} + \overrightarrow{BD} + \overrightarrow{DF}\$.
 
• Áp dụng quy tắc cộng vectơ liên tiếp: $\$ VT = $\overrightarrow{FD}$ + $\overrightarrow{DF}$ = $\overrightarrow{FF}$ = $\overrightarrow{0}\.$
}
\end{ex}

% ===== Câu 114 =====
% ID: LATEX_495D545B9A85
% Mức: TH
\dangbt{Tính độ dài của tổng và hiệu hai vectơ}
\begin{ex}
% ID: LATEX_495D545B9A85
% Mức: TH
Cho tam giác $\,\text{ABC}\đều cạnh\a\$. Tính độ dài của vectơ $\\overrightarrow{AB}$ - $\overrightarrow{AC}\$.
\choice
{\True $\$ a $\$}
{$\$ a $\sqrt{3}\$}
{$\2\,\text{a}$ \}
{$\$ 0 $\$}
\loigiai{
• Áp dụng quy tắc hiệu hai vectơ chung gốc, ta có: $\\overrightarrow{AB}$ - $\overrightarrow{AC}$ = $\overrightarrow{CB}\$.
 
• Độ dài của vectơ này là: $\$ |\overrightarrow{AB} - \overrightarrow{AC}| = |\overrightarrow{CB}| = CB\ $.
 
• Vì tam giác$\$ABC$\đều cạnh\a\$nên$\$CB = a$ \.
}
\end{ex}

% ===== Câu 115 =====
% ID: LATEX_3704ABD1EBCC
% Mức: TH
\begin{ex}
% ID: LATEX_3704ABD1EBCC
% Mức: TH
Cho tam giác $\,\text{ABC}\đều cạnh\a\$. Tính độ dài của vectơ $\\overrightarrow{AB}$ + $\overrightarrow{AC}\$.
\choice
{$\$ a $\$}
{\True $\$ a $\sqrt{3}\$}
{$\2\,\text{a}$ \}
{$\$ a $\sqrt{2}\$}
\loigiai{
• Gọi $\,\text{M}\là trung điểm của cạnh\,\text{BC}\$. Theo tính chất trung điểm, ta có $\\overrightarrow{AB}$ + $\overrightarrow{AC}$ = 2 $\overrightarrow{AM}\$.
 
• Do đó, $\$ |\overrightarrow{AB} + \overrightarrow{AC}| = |2\overrightarrow{AM}| = 2AM\ $.
 
• Trong tam giác đều cạnh$\$a$\, đường trung tuyến đồng thời là đường cao nên\,\text{AM} = \dfrac{a\sqrt{3}}{2}\$.
 
• Vậy$\$|$\overrightarrow{AB}$+$\overrightarrow{AC}| = 2\cdot\dfrac{a\sqrt{3}}{2}$= a$ \sqrt{3}\.
}
\end{ex}

% ===== Câu 116 =====
% ID: LATEX_BD8F813D04FB
% Mức: TH
\begin{ex}
% ID: LATEX_BD8F813D04FB
% Mức: TH
Cho hình vuông $\$ ABCD $\có cạnh bằng\a\$. Tính độ dài của vectơ $\\overrightarrow{AB}$ + $\overrightarrow{AD}\$.
\choice
{$\2\,\text{a}$ \}
{$\$ a $\$}
{\True $\$ a $\sqrt{2}\$}
{$\\dfrac{a\sqrt{2}}{2}\$}
\loigiai{
• Vì $\$ ABCD $\là hình vuông nên nó cũng là hình bình hành. Áp dụng quy tắc hình bình hành ta có:\\overrightarrow{AB} + \overrightarrow{AD} = \overrightarrow{AC}\$.
 
• Độ dài vectơ cần tìm là $\$ | $\overrightarrow{AB}$ + $\overrightarrow{AD}| = |\overrightarrow{AC}| = AC\$.
 
• Trong hình vuông cạnh $\$ a\ $, độ dài đường chéo là$ \ $AC =$ \sqrt{AB^2 + BC^2} $=$ \sqrt{a^2 + a^2} $= a$ \sqrt{2}\.
}
\end{ex}

% ===== Câu 117 =====
% ID: LATEX_331847432B03
% Mức: TH
\begin{ex}
% ID: LATEX_331847432B03
% Mức: TH
Cho hình vuông $\$ ABCD $\tâm\,\text{O}\$, cạnh $\$ a $\. Tính độ dài của vectơ\\overrightarrow{OA} - \overrightarrow{CB}\$.
\choice
{\True $\\dfrac{a\sqrt{2}}{2}\$}
{$\$ a $\sqrt{2}\$}
{$\$ a $\$}
{$\\dfrac{a}{2}\$}
\loigiai{
• Vì $\$ ABCD $\là hình vuông nên\\overrightarrow{CB} = \overrightarrow{DA}\$.
 
• Thay vào biểu thức ta có: $\\overrightarrow{OA}$ - $\overrightarrow{CB}$ = $\overrightarrow{OA}$ - $\overrightarrow{DA}$ = $\overrightarrow{OA}$ + $\overrightarrow{AD}\$.
 
• Áp dụng quy tắc ba điểm: $\$ \overrightarrow{OA} + \overrightarrow{AD} = \overrightarrow{OD}\ $.
 
• Độ dài vectơ là$\$|$\overrightarrow{OD}| = OD\. Vì\,\text{O}\$là tâm hình vuông nên$\$OD =$\dfrac{1}{2}BD$=$ \dfrac{a\sqrt{2}}{2}\.
}
\end{ex}

% ===== Câu 118 =====
% ID: LATEX_9E693EBEFA4D
% Mức: TH
\begin{ex}
% ID: LATEX_9E693EBEFA4D
% Mức: TH
Cho hình chữ nhật $\$ ABCD $\có\AB = 3\$, $\$ BC = 4 $\. Tính độ dài của vectơ\\overrightarrow{AB} + \overrightarrow{BC}\$.
\choice
{$\$ 7 $\$}
{\True $\$ 5 $\$}
{$\$ 1 $\$}
{$\$ 12 $\$}
\loigiai{
• Theo quy tắc ba điểm, ta có $\\overrightarrow{AB}$ + $\overrightarrow{BC}$ = $\overrightarrow{AC}\$.
 
• Độ dài của vectơ là $\$ |\overrightarrow{AC}| = AC\ $.
 
• Xét tam giác$\$ABC$\vuông tại\,\text{B}\$, theo định lí Pythagore:$\$AC =$\sqrt{AB^2 + BC^2}$=$\sqrt{3^2 + 4^2}$=$\sqrt{25}$= 5$ \.
}
\end{ex}

% ===== Câu 119 =====
% ID: LATEX_6F0ACD0D2B2C
% Mức: TH
\dangbt{Xác định điểm thỏa mãn đẳng thức vectơ}
\begin{ex}
% ID: LATEX_6F0ACD0D2B2C
% Mức: TH
Cho đoạn thẳng $\,\text{AB}\và điểm\,\text{M}\$ thỏa mãn đẳng thức $\\overrightarrow{MA}$ + $\overrightarrow{MB}$ = $\overrightarrow{0}\. Khẳng định nào sau đây đúng?$
\choice
{\True $\M\là trung điểm của đoạn thẳng\AB\$}
{$\M\trùng với điểm\A\$}
{$\M\trùng với điểm\B\$}
{$\M\nằm ngoài đoạn thẳng\AB\$}
\loigiai{
• Ta có đẳng thức $\\overrightarrow{MA}$ + $\overrightarrow{MB}$ = $\overrightarrow{0}\Leftrightarrow\overrightarrow{MA}$ = - $\overrightarrow{MB}\$.
 
• Điều này có nghĩa là hai vectơ $\$ \overrightarrow{MA}\ $và$ \ $\overrightarrow{MB}\$ là hai vectơ đối nhau, tức là chúng ngược hướng và có độ dài bằng nhau ($\$ MA = MB\ $).
 
• Do đó, ba điểm$\,\text{M}$,$A$,$B$\thẳng hàng,\M\$nằm giữa$\A$,$B$\và cách đều\,\text{A}, B\$. Vậy$\,\text{M}$\là trung điểm của\AB\$.
}
\end{ex}

% ===== Câu 120 =====
% ID: LATEX_40087EE97ADD
% Mức: TH
\begin{ex}
% ID: LATEX_40087EE97ADD
% Mức: TH
Cho ba điểm $\,\text{A}, B, C\phân biệt. Tìm điểm\,\text{M}\$ thỏa mãn $\\overrightarrow{MA}$ - $\overrightarrow{MB}$ + $\overrightarrow{MC}$ = $\overrightarrow{0}\$.
\choice
{$\M\là trung điểm của\AC\$}
{\True $\M\là đỉnh thứ tư của hình bình hành\ABCM\$}
{$\M\là đỉnh thứ tư của hình bình hành\ABMC\$}
{$\M\là trọng tâm tam giác\ABC\$}
\loigiai{
• Áp dụng quy tắc hiệu cho hai vectơ đầu tiên: $\\overrightarrow{MA}$ - $\overrightarrow{MB}$ = $\overrightarrow{BA}\$.
 
• Đẳng thức đã cho trở thành: $\$ \overrightarrow{BA} + \overrightarrow{MC} = \overrightarrow{0} \Leftrightarrow \overrightarrow{MC} = -\overrightarrow{BA} \Leftrightarrow \overrightarrow{MC} = \overrightarrow{AB}\ $.
 
• Điều kiện$overrightarrow{MC}$=$\overrightarrow{AB}\chứng tỏ tứ giác\,\text{ABCM}\$là hình bình hành (với thứ tự các đỉnh liên tiếp là$\,\text{A}$,$B$,$C$,$M$\).$
}
\end{ex}

% ===== Câu 121 =====
% ID: LATEX_48C88D4FB932
% Mức: TH
\begin{ex}
% ID: LATEX_48C88D4FB932
% Mức: TH
Cho tam giác $\,\text{ABC}\. Gọi\,\text{M}\$ là điểm xác định bởi đẳng thức $\\overrightarrow{AM}$ = $\overrightarrow{AB}$ + $\overrightarrow{AC}\. Khẳng định nào sau đây là đúng?$
\choice
{$\M\là trung điểm của\BC\$}
{$\M\là trọng tâm tam giác\ABC\$}
{\True $\$ ABMC $\là hình bình hành$}
{$\$ AMBC $\là hình bình hành$}
\loigiai{
• Theo quy tắc hình bình hành, tổng hai vectơ chung gốc $\\overrightarrow{AB} + \overrightarrow{AC}\bằng vectơ đường chéo\\overrightarrow{AM}\$ của hình bình hành có ba đỉnh là $\,\text{A}$, $B$, $C$ \ $.
 
• Do đó, tứ giác tạo bởi hai cạnh$\$AB\$và$\$AC$\$cùng với đường chéo$\$AM\$là hình bình hành$\$ABMC$\$.
 
• Chú ý thứ tự đỉnh: \$A\$nối với$\$B$\$và$\$C\$,$\$M$\$đối diện$\$A\$, nên tứ giác đọc theo vòng tròn là$\$ABMC$ \.
}
\end{ex}

% ===== Câu 122 =====
% ID: LATEX_E0B4D02AEB3B
% Mức: TH
\begin{ex}
% ID: LATEX_E0B4D02AEB3B
% Mức: TH
Cho bốn điểm phân biệt $\,\text{A}, B, C, D\. Gọi\,\text{M}\$ là điểm thỏa mãn $\\overrightarrow{AM}$ = $\overrightarrow{AB}$ - $\overrightarrow{AC}\. Khẳng định nào sau đây đúng?$
\choice
{\True $\\overrightarrow{AM}=\overrightarrow{CB}\$}
{$\\overrightarrow{AM}=\overrightarrow{BC}\$}
{$\\overrightarrow{AM}=\overrightarrow{CD}\$}
{$\\overrightarrow{AM}=\overrightarrow{AD}\$}
\loigiai{
• Áp dụng quy tắc trừ hai vectơ chung gốc: $\\overrightarrow{AB}$ - $\overrightarrow{AC}$ = $\overrightarrow{CB}\$.
 
• Vậy đẳng thức đã cho trở thành $\$ \overrightarrow{AM} = \overrightarrow{CB}\.
}
\end{ex}

% ===== Câu 123 =====
% ID: LATEX_D8CCCEF6C78E
% Mức: TH
\begin{ex}
% ID: LATEX_D8CCCEF6C78E
% Mức: TH
Cho tam giác $\,\text{ABC}\. Tập hợp các điểm\,\text{M}\$ thỏa mãn $\$ | $\overrightarrow{MA}$ + $\overrightarrow{MB}| = |\overrightarrow{MC}$ + $\overrightarrow{MB}|\là:$
\choice
{Đường tròn tâm $\B\bán kính\AC\$}
{Đường tròn đường kính $\$ AC $\$}
{\True Đường trung trực của đoạn thẳng $\$ AC $\$}
{Đường trung trực của đoạn thẳng $\$ BC $\$}
\loigiai{
• Xét vế trái: Gọi $\,\text{I}\là trung điểm của\,\text{AB}\$, ta có $\\overrightarrow{MA}$ + $\overrightarrow{MB}$ = 2 $\overrightarrow{MI}\Rightarrow$ | $\overrightarrow{MA}$ + $\overrightarrow{MB}| = 2MI\$.
 
• Xét vế phải: Gọi $\$ J\ $là trung điểm của$ \ $BC$ \ $, ta có$ \ $\overrightarrow{MC} + \overrightarrow{MB} = 2\overrightarrow{MJ} \Rightarrow |\overrightarrow{MC} + \overrightarrow{MB}| = 2MJ\$.
 
• Đẳng thức trở thành $\$ 2MI = 2MJ \Leftrightarrow MI = MJ $\$.
 
• Do đó, tập hợp các điểm $\$ M\ $là đường trung trực của đoạn thẳng$ \ $IJ$ \ $.
 
• Mặt khác, \$IJ\$là đường trung bình của tam giác$\$ABC$\$nên$\$IJ \parallel AC\$. Đường trung trực của$\$IJ$\$cũng chính là đường trung trực của đoạn$\$AC\$(nếu xét tam giác cân) - khoảng cách. Tuy nhiên, cách giải chuẩn hơn mà không cần gọi trung điểm là:
 
• Cách 2:$\$|$\overrightarrow{MA}$+$\overrightarrow{MB}| = |\overrightarrow{MB}$+$\overrightarrow{MC}|\. Đây là độ dài của tổng, không dễ khử điểm\,\text{M}\$.
 
• Chờ đã,$overrightarrow{MB}$+$\overrightarrow{MC}\chưa triệt tiêu\,\text{M}\$. Phải dùng trung điểm$\,\text{I}$,$J$\. Khi\MI = MJ\$,$\M$\nằm trên trung trực của\,\text{IJ}\$. Đường trung trực của$\$IJ$\không nhất thiết là đường trung trực của\,\text{AC}\$.
 
• Hãy đổi câu hỏi thành một bài toán quen thuộc hơn:$\$|$\overrightarrow{MB}$-$\overrightarrow{MC}| = |\overrightarrow{BM}$-$\overrightarrow{BA}|\$.
 
• Sửa lại câu hỏi: Cho tam giác$\$ABC\$. Điểm$\$M$\$thỏa mãn$\$ \overrightarrow{MA} + \overrightarrow{MB} + \overrightarrow{MC} = \overrightarrow{0}\.
}
\end{ex}

% ===== Câu 124 =====
% ID: LATEX_29094683BBF1
% Mức: TH
\dangbt{Tính chất trung điểm và trọng tâm trong tổng vectơ}
\begin{ex}
% ID: LATEX_29094683BBF1
% Mức: TH
Cho tam giác $\ABC\có trọng tâm\G\$. Khẳng định nào sau đây là đúng?
\choice
{$\\overrightarrow{GA} + \overrightarrow{GB} + \overrightarrow{GC} = 3\overrightarrow{GM}\với\,\text{M}\$ là trung điểm $\$ BC $\$}
{$\\overrightarrow{AB}$ + $\overrightarrow{AC}$ = $\overrightarrow{AG}\$}
{\True $\\overrightarrow{GA}$ + $\overrightarrow{GB}$ + $\overrightarrow{GC}$ = $\overrightarrow{0}\$}
{$\\overrightarrow{GA}$ + $\overrightarrow{GB}$ = $\overrightarrow{GC}\$}
\loigiai{
• Theo tính chất trọng tâm của tam giác, tổng các vectơ hướng từ trọng tâm đến ba đỉnh luôn bằng vectơ-không.
 
• Vậy $\\overrightarrow{GA}$ + $\overrightarrow{GB}$ + $\overrightarrow{GC}$ = $\overrightarrow{0}\.$
}
\end{ex}

% ===== Câu 125 =====
% ID: LATEX_1BA9FFC3EF07
% Mức: TH
\begin{ex}
% ID: LATEX_1BA9FFC3EF07
% Mức: TH
Cho hình bình hành $\$ ABCD $\tâm\O\$. Đẳng thức nào sau đây là đúng?
\choice
{\True $\\overrightarrow{OA}$ + $\overrightarrow{OB}$ + $\overrightarrow{OC}$ + $\overrightarrow{OD}$ = $\overrightarrow{0}\$}
{$\\overrightarrow{OA}$ + $\overrightarrow{OB}$ + $\overrightarrow{OC}$ + $\overrightarrow{OD}$ = 2 $\overrightarrow{AB}\$}
{$\\overrightarrow{OA}$ + $\overrightarrow{OB}$ + $\overrightarrow{OC}$ + $\overrightarrow{OD}$ = $\overrightarrow{AC}\$}
{$\\overrightarrow{OA}$ + $\overrightarrow{OB}$ + $\overrightarrow{OC}$ + $\overrightarrow{OD}$ = 4 $\overrightarrow{O}\$}
\loigiai{
• Vì $\,\text{O}\là tâm hình bình hành\,\text{ABCD}\$ nên $\,\text{O}$ \là trung điểm của đường chéo\AC\ $và đường chéo$ \ $BD$ \ $.
 
• Áp dụng tính chất trung điểm:$\$\overrightarrow{OA} + \overrightarrow{OC} = \overrightarrow{0}\$và$\$\overrightarrow{OB}$+$\overrightarrow{OD}$=$\overrightarrow{0}\$.
 
• Cộng vế theo vế ta được: \$ \overrightarrow{OA} + \overrightarrow{OB} + \overrightarrow{OC} + \overrightarrow{OD} = \overrightarrow{0} + \overrightarrow{0} = \overrightarrow{0}\.
}
\end{ex}

% ===== Câu 126 =====
% ID: LATEX_C2ED8EBD2F28
% Mức: TH
\begin{ex}
% ID: LATEX_C2ED8EBD2F28
% Mức: TH
Cho đoạn thẳng $\,\text{AB}\có trung điểm là\,\text{I}\$. Với điểm $\,\text{M}$ \tùy ý, khẳng định nào sau đây là đúng?
\choice
{$\\overrightarrow{MA}$ + $\overrightarrow{MB}$ = $\overrightarrow{MI}\$}
{\True $\\overrightarrow{MA}$ + $\overrightarrow{MB}$ = 2 $\overrightarrow{MI}\$}
{$\\overrightarrow{MA}$ + $\overrightarrow{MB}$ = 3 $\overrightarrow{MI}\$}
{$\\overrightarrow{MA}$ + $\overrightarrow{MB}$ = $\overrightarrow{IM}\$}
\loigiai{
• Chèn điểm $\,\text{I}\vào hai vectơ ở vế trái:\\overrightarrow{MA} = \overrightarrow{MI} + \overrightarrow{IA}\$ và $\\overrightarrow{MB}$ = $\overrightarrow{MI}$ + $\overrightarrow{IB}\$.
 
• Cộng lại ta có: $\$ \overrightarrow{MA} + \overrightarrow{MB} = 2\overrightarrow{MI} + (\overrightarrow{IA} + \overrightarrow{IB})\ $.
 
• Vì$\,\text{I}$\là trung điểm của\AB\$nên$\\overrightarrow{IA}$+$\overrightarrow{IB}$=$\overrightarrow{0}\$.
 
• Vậy $\$ \overrightarrow{MA} + \overrightarrow{MB} = 2\overrightarrow{MI}\.
}
\end{ex}

% ===== Câu 127 =====
% ID: LATEX_BF46C9E9C4F6
% Mức: TH
\begin{ex}
% ID: LATEX_BF46C9E9C4F6
% Mức: TH
Cho tam giác $\ABC\. Gọi\M, N, P\$ lần lượt là trung điểm của các cạnh $\$ BC, CA, AB $\. Tổng\\overrightarrow{AM} + \overrightarrow{BN} + \overrightarrow{CP}\$ bằng:
\choice
{$\\overrightarrow{AB}\$}
{$\\overrightarrow{BC}\$}
{\True $\\overrightarrow{0}\$}
{$\\overrightarrow{CA}\$}
\loigiai{
• Theo tính chất trung điểm, ta có $\2\overrightarrow{AM} = \overrightarrow{AB} + \overrightarrow{AC}\;\2\overrightarrow{BN} = \overrightarrow{BA} + \overrightarrow{BC}\$; $\$ 2 $\overrightarrow{CP}$ = $\overrightarrow{CA}$ + $\overrightarrow{CB}\$.
 
• Cộng ba đẳng thức trên: $\$ 2(\overrightarrow{AM} + \overrightarrow{BN} + \overrightarrow{CP}) = (\overrightarrow{AB} + \overrightarrow{BA}) + (\overrightarrow{AC} + \overrightarrow{CA}) + (\overrightarrow{BC} + \overrightarrow{CB})\ $.
 
•$\$2$\overrightarrow{AM} + \overrightarrow{BN} + \overrightarrow{CP})$=$\overrightarrow{0}$+$\overrightarrow{0}$+$\overrightarrow{0}$=$\overrightarrow{0}\$.
 
• Vậy$\$\overrightarrow{AM} + \overrightarrow{BN} + \overrightarrow{CP} = \overrightarrow{0}\$.
}
\end{ex}

% ===== Câu 128 =====
% ID: LATEX_AE13D1889B7A
% Mức: TH
\begin{ex}
% ID: LATEX_AE13D1889B7A
% Mức: TH
Cho tam giác $\,\text{ABC}\có trọng tâm\,\text{G}\$. Với điểm $\,\text{M}$ \bất kì, đẳng thức nào sau đây đúng?
\choice
{$\\overrightarrow{MA}$ + $\overrightarrow{MB}$ + $\overrightarrow{MC}$ = $\overrightarrow{MG}\$}
{$\\overrightarrow{MA}$ + $\overrightarrow{MB}$ + $\overrightarrow{MC}$ = 2 $\overrightarrow{MG}\$}
{\True $\\overrightarrow{MA}$ + $\overrightarrow{MB}$ + $\overrightarrow{MC}$ = 3 $\overrightarrow{MG}\$}
{$\\overrightarrow{MA}$ + $\overrightarrow{MB}$ + $\overrightarrow{MC}$ = 4 $\overrightarrow{MG}\$}
\loigiai{
• Chèn điểm $\,\text{G}\vào từng vectơ ở vế trái:\\overrightarrow{MA} = \overrightarrow{MG} + \overrightarrow{GA}\$; $\\overrightarrow{MB}$ = $\overrightarrow{MG}$ + $\overrightarrow{GB}\;\\overrightarrow{MC} = \overrightarrow{MG} + \overrightarrow{GC}\$.
 
• Cộng ba đẳng thức lại: $\\overrightarrow{MA}$ + $\overrightarrow{MB}$ + $\overrightarrow{MC}$ = 3 $\overrightarrow{MG}$ + $\overrightarrow{GA} + \overrightarrow{GB} + \overrightarrow{GC})$ \ $.
 
• Vì$\$G\$là trọng tâm tam giác$\$ABC$\$nên$\$\overrightarrow{GA} + \overrightarrow{GB} + \overrightarrow{GC} = \overrightarrow{0}\$.
 
• Vậy$overrightarrow{MA}$+$\overrightarrow{MB}$+$\overrightarrow{MC}$= 3$ \overrightarrow{MG}\.
}
\end{ex}

% ===== Câu 129 =====
% ID: LATEX_0DB1EA144BE3
% Mức: TH
\dangbt{Áp dụng quy tắc ba điểm và quy tắc hình bình hành}
\begin{ex}
% ID: LATEX_0DB1EA144BE3
% Mức: TH
Cho ba điểm phân biệt $\A, B, C\. Đẳng thức nào sau đây đúng?$
\choice
{$\\overrightarrow{AB}$ + $\overrightarrow{CB}$ = $\overrightarrow{AC}\$}
{$\\overrightarrow{CB}$ + $\overrightarrow{AC}$ = $\overrightarrow{AB}\$}
{$\\overrightarrow{AB}$ + $\overrightarrow{BC}$ = $\overrightarrow{CA}\$}
{\True $\\overrightarrow{AB}$ + $\overrightarrow{BC}$ = $\overrightarrow{AC}\$}
\loigiai{
Theo quy tắc ba điểm đối với phép cộng vectơ, với ba điểm $\A, B, C\bất kì ta luôn có\\overrightarrow{AB} + \overrightarrow{BC} = \overrightarrow{AC}\$.
}
\end{ex}

% ===== Câu 130 =====
% ID: LATEX_AF6B2F66DC9A
% Mức: TH
\begin{ex}
% ID: LATEX_AF6B2F66DC9A
% Mức: TH
Cho ba điểm phân biệt $\M, N, P\. Vectơ tổng\\overrightarrow{NP} + \overrightarrow{MN}\$ bằng vectơ nào sau đây?
\choice
{$\\overrightarrow{PN}\$}
{$\\overrightarrow{PM}\$}
{\True $\\overrightarrow{MP}\$}
{$\\overrightarrow{NM}\$}
\loigiai{
• Áp dụng tính chất giao hoán của phép cộng vectơ, ta có: $\\overrightarrow{NP}$ + $\overrightarrow{MN}$ = $\overrightarrow{MN}$ + $\overrightarrow{NP}\$.
 
• Theo quy tắc ba điểm, $\$ \overrightarrow{MN} + \overrightarrow{NP} = \overrightarrow{MP}\ $.
 
• Vậy$overrightarrow{NP}$+$\overrightarrow{MN}$=$ \overrightarrow{MP}\.
}
\end{ex}

% ===== Câu 131 =====
% ID: LATEX_B3959566D3C6
% Mức: TH
\begin{ex}
% ID: LATEX_B3959566D3C6
% Mức: TH
Cho hình bình hành $\$ ABCD $\. Đẳng thức nào sau đây đúng?$
\choice
{\True $\\overrightarrow{AB}$ + $\overrightarrow{AD}$ = $\overrightarrow{AC}\$}
{$\\overrightarrow{AB}$ + $\overrightarrow{AD}$ = $\overrightarrow{BD}\$}
{$\\overrightarrow{AB}$ + $\overrightarrow{AD}$ = $\overrightarrow{CD}\$}
{$\\overrightarrow{AB}$ + $\overrightarrow{AD}$ = $\overrightarrow{CB}\$}
\loigiai{
Theo quy tắc hình bình hành, vì $\$ ABCD $\là hình bình hành nên tổng hai vectơ cạnh chung gốc bằng vectơ đường chéo xuất phát từ gốc đó, tức là\\overrightarrow{AB} + \overrightarrow{AD} = \overrightarrow{AC}\$.
}
\end{ex}

% ===== Câu 132 =====
% ID: LATEX_1431DB984C9E
% Mức: TH
\begin{ex}
% ID: LATEX_1431DB984C9E
% Mức: TH
Cho bốn điểm phân biệt $\A, B, C, O\. Đẳng thức nào sau đây đúng?$
\choice
{$\\overrightarrow{AB}=\overrightarrow{OB}$ + $\overrightarrow{OA}\$}
{$\\overrightarrow{AB}=\overrightarrow{AC}$ + $\overrightarrow{BC}\$}
{\True $\\overrightarrow{OA}=\overrightarrow{CA}$ - $\overrightarrow{CO}\$}
{$\\overrightarrow{OA}=\overrightarrow{OB}$ - $\overrightarrow{BA}\$}
\loigiai{
• Áp dụng quy tắc trừ vectơ với ba điểm $\O, C, A\ta có:\\overrightarrow{CA} - \overrightarrow{CO} = \overrightarrow{OA}\$.
 
• Các phương án khác đều sai quy tắc cộng và trừ vectơ.
}
\end{ex}

% ===== Câu 133 =====
% ID: LATEX_A46C88793698
% Mức: TH
\begin{ex}
% ID: LATEX_A46C88793698
% Mức: TH
Cho bốn điểm phân biệt $\A, B, C, D\. Vectơ tổng\\overrightarrow{AB} + \overrightarrow{CD} + \overrightarrow{BC} + \overrightarrow{DA}\$ bằng vectơ nào sau đây?
\choice
{\True $\\overrightarrow{0}\$}
{$\\overrightarrow{AC}\$}
{$\\overrightarrow{BD}\$}
{$\\overrightarrow{BA}\$}
\loigiai{
• Ta sắp xếp lại các số hạng và áp dụng quy tắc ba điểm liên tiếp: $\\overrightarrow{AB}$ + $\overrightarrow{CD}$ + $\overrightarrow{BC}$ + $\overrightarrow{DA}$ = $\overrightarrow{AB}$ + $\overrightarrow{BC}$ + $\overrightarrow{CD}$ + $\overrightarrow{DA}\$.
 
• Thực hiện phép cộng từ trái sang phải: $\$ \overrightarrow{AC} + \overrightarrow{CD} + \overrightarrow{DA} = \overrightarrow{AD} + \overrightarrow{DA} = \overrightarrow{AA}\ $.
 
• Vì điểm đầu và điểm cuối trùng nhau nên$overrightarrow{AA}$=$ \overrightarrow{0}\.
}
\end{ex}

% ===== Câu 134 =====
% ID: LATEX_4302F8D17026
% Mức: TH
\dangbt{Rút gọn biểu thức và chứng minh đẳng thức vectơ}
\begin{ex}
% ID: LATEX_4302F8D17026
% Mức: TH
Cho bốn điểm $\A, B, C, D\bất kì. Khẳng định nào sau đây là đúng?$
\choice
{$\\overrightarrow{AB}$ - $\overrightarrow{AD}$ = $\overrightarrow{CD}$ + $\overrightarrow{BC}\$}
{$\\overrightarrow{AB}$ - $\overrightarrow{AC}$ = $\overrightarrow{DB}$ - $\overrightarrow{DC}\$}
{$\\overrightarrow{BC}$ + $\overrightarrow{AB}$ = $\overrightarrow{DA}$ - $\overrightarrow{DC}\$}
{\True $\\overrightarrow{AB}$ + $\overrightarrow{CD}$ = $\overrightarrow{AD}$ + $\overrightarrow{CB}\$}
\loigiai{
• Xét vế trái: $\\overrightarrow{AB} + \overrightarrow{CD}\. Áp dụng quy tắc chèn điểm\,\text{D}\$ vào vectơ thứ nhất và điểm $\B$ \ $vào vectơ thứ hai.
 
• Ta có$ \ $\overrightarrow{AB} + \overrightarrow{CD} = (\overrightarrow{AD} + \overrightarrow{DB}) + (\overrightarrow{CB} + \overrightarrow{BD})\$.
 
• Nhóm lại ta được: $\\overrightarrow{AD}$ + $\overrightarrow{CB}$ + $(\overrightarrow{DB} + \overrightarrow{BD})$ = $\overrightarrow{AD}$ + $\overrightarrow{CB}$ + $\overrightarrow{0}$ = $\overrightarrow{AD}$ + $\overrightarrow{CB}\. Đẳng thức đúng.$
}
\end{ex}

% ===== Câu 135 =====
% ID: LATEX_3BCD556B4A93
% Mức: TH
\begin{ex}
% ID: LATEX_3BCD556B4A93
% Mức: TH
Rút gọn biểu thức vectơ $\\overrightarrow{P} = \overrightarrow{AB} - \overrightarrow{AC} + \overrightarrow{BD}\, ta được kết quả nào sau đây?$
\choice
{$\\overrightarrow{AD}\$}
{$\\overrightarrow{BC}\$}
{\True $\\overrightarrow{CD}\$}
{$\\overrightarrow{DA}\$}
\loigiai{
• Áp dụng quy tắc hiệu cho hai vectơ đầu tiên: $\\overrightarrow{AB}$ - $\overrightarrow{AC}$ = $\overrightarrow{CB}\$.
 
• Thay vào biểu thức ta được: $\$ \overrightarrow{P} = \overrightarrow{CB} + \overrightarrow{BD}\ $.
 
• Áp dụng quy tắc ba điểm:$overrightarrow{CB}$+$\overrightarrow{BD}$=$ \overrightarrow{CD}\.
}
\end{ex}

% ===== Câu 136 =====
% ID: LATEX_0658205A4C5B
% Mức: TH
\begin{ex}
% ID: LATEX_0658205A4C5B
% Mức: TH
Cho hình bình hành $\$ ABCD $\. Rút gọn biểu thức\\overrightarrow{AB} - \overrightarrow{DA} - \overrightarrow{DC}\$ ta được vectơ nào?
\choice
{$\\overrightarrow{AC}\$}
{\True $\\overrightarrow{0}\$}
{$\\overrightarrow{BD}\$}
{$\$ 2 $\overrightarrow{AB}\$}
\loigiai{
• Vì $\$ ABCD $\là hình bình hành nên\\overrightarrow{DC} = \overrightarrow{AB}\$.
 
• Biểu thức trở thành: $\\overrightarrow{AB}$ - $\overrightarrow{DA}$ - $\overrightarrow{AB}$ = - $\overrightarrow{DA}$ = $\overrightarrow{AD}\$.
 
• Tuy nhiên, ta có thể làm cách khác: $\$ \overrightarrow{AB} - \overrightarrow{DA} - \overrightarrow{DC} = \overrightarrow{AB} + \overrightarrow{AD} - \overrightarrow{DC}\ $.
 
• Theo quy tắc hình bình hành$overrightarrow{AB}$+$\overrightarrow{AD}$=$ \overrightarrow{AC}\. Mà
overrightarrow{DC} = \overrightarrow{AB}\ $không giúp ra$ overrightarrow{0}\ $.
 
• Hãy thử lại:$\$\overrightarrow{AB} - \overrightarrow{DA} - \overrightarrow{DC} = \overrightarrow{AB} + \overrightarrow{AD} - \overrightarrow{DC} = \overrightarrow{AC} - \overrightarrow{DC} = \overrightarrow{AC} + \overrightarrow{CD} = \overrightarrow{AD}\$.
 
• Xin lỗi, có sự nhầm lẫn. Hãy làm lại: $\\overrightarrow{AB}$ - $\overrightarrow{DA}$ - $\overrightarrow{DC}$ = $\overrightarrow{AB}$ + $\overrightarrow{AD}$ - $\overrightarrow{DC}\. Mà\,\text{ABCD}\$ là hình bình hành nên $\\overrightarrow{AB}$ = $\overrightarrow{DC}\, do đó\\overrightarrow{AB} - \overrightarrow{DC} = \overrightarrow{0}\$.
 
• Vậy biểu thức còn lại là $\$ - $\overrightarrow{DA}$ = $\overrightarrow{AD}\$.
 
• Phát hiện phương án không khớp, sửa lại câu hỏi: Rút gọn $\$ \overrightarrow{AB} + \overrightarrow{AD} - \overrightarrow{AC}\.
}
\end{ex}

% ===== Câu 137 =====
% ID: LATEX_A9B2A6742B7D
% Mức: TH
\begin{ex}
% ID: LATEX_A9B2A6742B7D
% Mức: TH
Cho hình bình hành $\$ ABCD $\. Biểu thức\\overrightarrow{AB} + \overrightarrow{AD} - \overrightarrow{AC}\$ bằng vectơ nào sau đây?
\choice
{$\\overrightarrow{BD}\$}
{$\\overrightarrow{CB}\$}
{\True $\\overrightarrow{0}\$}
{$\\overrightarrow{DA}\$}
\loigiai{
• Áp dụng quy tắc hình bình hành, ta có tổng $\\overrightarrow{AB}$ + $\overrightarrow{AD}$ = $\overrightarrow{AC}\$.
 
• Thay vào biểu thức ta được: $\$ \overrightarrow{AC} - \overrightarrow{AC} = \overrightarrow{0}\.
}
\end{ex}

% ===== Câu 138 =====
% ID: LATEX_4FD940D16699
% Mức: TH
\begin{ex}
% ID: LATEX_4FD940D16699
% Mức: TH
Cho $\6\điểm\A, B, C, D, E, F\$ bất kì. Đẳng thức nào sau đây đúng?
\choice
{$\\overrightarrow{AB}$ + $\overrightarrow{CD}$ = $\overrightarrow{AD}$ + $\overrightarrow{BC}\$}
{$\\overrightarrow{AD}$ + $\overrightarrow{BE}$ + $\overrightarrow{CF}$ = $\overrightarrow{AE}$ + $\overrightarrow{BF}$ + $\overrightarrow{CE}\$}
{\True $\\overrightarrow{AB}$ - $\overrightarrow{AF}$ + $\overrightarrow{CD}$ - $\overrightarrow{CB}$ + $\overrightarrow{EF}$ - $\overrightarrow{ED}$ = $\overrightarrow{0}\$}
{$\\overrightarrow{AB}$ + $\overrightarrow{CD}$ + $\overrightarrow{EF}$ = $\overrightarrow{AF}$ + $\overrightarrow{ED}$ + $\overrightarrow{BC}\$}
\loigiai{
• Xét vế trái của phương án đúng: $\VT =(\overrightarrow{AB} - \overrightarrow{AF})$ + $(\overrightarrow{CD} - \overrightarrow{CB})$ + $(\overrightarrow{EF} - \overrightarrow{ED})$ \ $.
 
• Áp dụng quy tắc trừ vectơ:$ \ $VT = \overrightarrow{FB} + \overrightarrow{BD} + \overrightarrow{DF}\$.
 
• Áp dụng quy tắc cộng vectơ liên tiếp: $\$ VT = $\overrightarrow{FD}$ + $\overrightarrow{DF}$ = $\overrightarrow{FF}$ = $\overrightarrow{0}\.$
}
\end{ex}

% ===== Câu 139 =====
% ID: LATEX_495D545B9A85
% Mức: TH
\dangbt{Tính độ dài của tổng và hiệu hai vectơ}
\begin{ex}
% ID: LATEX_495D545B9A85
% Mức: TH
Cho tam giác $\,\text{ABC}\đều cạnh\a\$. Tính độ dài của vectơ $\\overrightarrow{AB}$ - $\overrightarrow{AC}\$.
\choice
{\True $\$ a $\$}
{$\$ a $\sqrt{3}\$}
{$\2\,\text{a}$ \}
{$\$ 0 $\$}
\loigiai{
• Áp dụng quy tắc hiệu hai vectơ chung gốc, ta có: $\\overrightarrow{AB}$ - $\overrightarrow{AC}$ = $\overrightarrow{CB}\$.
 
• Độ dài của vectơ này là: $\$ |\overrightarrow{AB} - \overrightarrow{AC}| = |\overrightarrow{CB}| = CB\ $.
 
• Vì tam giác$\$ABC$\đều cạnh\a\$nên$\$CB = a$ \.
}
\end{ex}

% ===== Câu 140 =====
% ID: LATEX_3704ABD1EBCC
% Mức: TH
\begin{ex}
% ID: LATEX_3704ABD1EBCC
% Mức: TH
Cho tam giác $\,\text{ABC}\đều cạnh\a\$. Tính độ dài của vectơ $\\overrightarrow{AB}$ + $\overrightarrow{AC}\$.
\choice
{$\$ a $\$}
{\True $\$ a $\sqrt{3}\$}
{$\2\,\text{a}$ \}
{$\$ a $\sqrt{2}\$}
\loigiai{
• Gọi $\,\text{M}\là trung điểm của cạnh\,\text{BC}\$. Theo tính chất trung điểm, ta có $\\overrightarrow{AB}$ + $\overrightarrow{AC}$ = 2 $\overrightarrow{AM}\$.
 
• Do đó, $\$ |\overrightarrow{AB} + \overrightarrow{AC}| = |2\overrightarrow{AM}| = 2AM\ $.
 
• Trong tam giác đều cạnh$\$a$\, đường trung tuyến đồng thời là đường cao nên\,\text{AM} = \dfrac{a\sqrt{3}}{2}\$.
 
• Vậy$\$|$\overrightarrow{AB}$+$\overrightarrow{AC}| = 2\cdot\dfrac{a\sqrt{3}}{2}$= a$ \sqrt{3}\.
}
\end{ex}

% ===== Câu 141 =====
% ID: LATEX_BD8F813D04FB
% Mức: TH
\begin{ex}
% ID: LATEX_BD8F813D04FB
% Mức: TH
Cho hình vuông $\$ ABCD $\có cạnh bằng\a\$. Tính độ dài của vectơ $\\overrightarrow{AB}$ + $\overrightarrow{AD}\$.
\choice
{$\2\,\text{a}$ \}
{$\$ a $\$}
{\True $\$ a $\sqrt{2}\$}
{$\\dfrac{a\sqrt{2}}{2}\$}
\loigiai{
• Vì $\$ ABCD $\là hình vuông nên nó cũng là hình bình hành. Áp dụng quy tắc hình bình hành ta có:\\overrightarrow{AB} + \overrightarrow{AD} = \overrightarrow{AC}\$.
 
• Độ dài vectơ cần tìm là $\$ | $\overrightarrow{AB}$ + $\overrightarrow{AD}| = |\overrightarrow{AC}| = AC\$.
 
• Trong hình vuông cạnh $\$ a\ $, độ dài đường chéo là$ \ $AC =$ \sqrt{AB^2 + BC^2} $=$ \sqrt{a^2 + a^2} $= a$ \sqrt{2}\.
}
\end{ex}

% ===== Câu 142 =====
% ID: LATEX_331847432B03
% Mức: TH
\begin{ex}
% ID: LATEX_331847432B03
% Mức: TH
Cho hình vuông $\$ ABCD $\tâm\,\text{O}\$, cạnh $\$ a $\. Tính độ dài của vectơ\\overrightarrow{OA} - \overrightarrow{CB}\$.
\choice
{\True $\\dfrac{a\sqrt{2}}{2}\$}
{$\$ a $\sqrt{2}\$}
{$\$ a $\$}
{$\\dfrac{a}{2}\$}
\loigiai{
• Vì $\$ ABCD $\là hình vuông nên\\overrightarrow{CB} = \overrightarrow{DA}\$.
 
• Thay vào biểu thức ta có: $\\overrightarrow{OA}$ - $\overrightarrow{CB}$ = $\overrightarrow{OA}$ - $\overrightarrow{DA}$ = $\overrightarrow{OA}$ + $\overrightarrow{AD}\$.
 
• Áp dụng quy tắc ba điểm: $\$ \overrightarrow{OA} + \overrightarrow{AD} = \overrightarrow{OD}\ $.
 
• Độ dài vectơ là$\$|$\overrightarrow{OD}| = OD\. Vì\,\text{O}\$là tâm hình vuông nên$\$OD =$\dfrac{1}{2}BD$=$ \dfrac{a\sqrt{2}}{2}\.
}
\end{ex}

% ===== Câu 143 =====
% ID: LATEX_9E693EBEFA4D
% Mức: TH
\begin{ex}
% ID: LATEX_9E693EBEFA4D
% Mức: TH
Cho hình chữ nhật $\$ ABCD $\có\AB = 3\$, $\$ BC = 4 $\. Tính độ dài của vectơ\\overrightarrow{AB} + \overrightarrow{BC}\$.
\choice
{$\$ 7 $\$}
{\True $\$ 5 $\$}
{$\$ 1 $\$}
{$\$ 12 $\$}
\loigiai{
• Theo quy tắc ba điểm, ta có $\\overrightarrow{AB}$ + $\overrightarrow{BC}$ = $\overrightarrow{AC}\$.
 
• Độ dài của vectơ là $\$ |\overrightarrow{AC}| = AC\ $.
 
• Xét tam giác$\$ABC$\vuông tại\,\text{B}\$, theo định lí Pythagore:$\$AC =$\sqrt{AB^2 + BC^2}$=$\sqrt{3^2 + 4^2}$=$\sqrt{25}$= 5$ \.
}
\end{ex}

% ===== Câu 144 =====
% ID: LATEX_6F0ACD0D2B2C
% Mức: TH
\dangbt{Xác định điểm thỏa mãn đẳng thức vectơ}
\begin{ex}
% ID: LATEX_6F0ACD0D2B2C
% Mức: TH
Cho đoạn thẳng $\,\text{AB}\và điểm\,\text{M}\$ thỏa mãn đẳng thức $\\overrightarrow{MA}$ + $\overrightarrow{MB}$ = $\overrightarrow{0}\. Khẳng định nào sau đây đúng?$
\choice
{\True $\M\là trung điểm của đoạn thẳng\AB\$}
{$\M\trùng với điểm\A\$}
{$\M\trùng với điểm\B\$}
{$\M\nằm ngoài đoạn thẳng\AB\$}
\loigiai{
• Ta có đẳng thức $\\overrightarrow{MA}$ + $\overrightarrow{MB}$ = $\overrightarrow{0}\Leftrightarrow\overrightarrow{MA}$ = - $\overrightarrow{MB}\$.
 
• Điều này có nghĩa là hai vectơ $\$ \overrightarrow{MA}\ $và$ \ $\overrightarrow{MB}\$ là hai vectơ đối nhau, tức là chúng ngược hướng và có độ dài bằng nhau ($\$ MA = MB\ $).
 
• Do đó, ba điểm$\,\text{M}$,$A$,$B$\thẳng hàng,\M\$nằm giữa$\A$,$B$\và cách đều\,\text{A}, B\$. Vậy$\,\text{M}$\là trung điểm của\AB\$.
}
\end{ex}

% ===== Câu 145 =====
% ID: LATEX_40087EE97ADD
% Mức: TH
\begin{ex}
% ID: LATEX_40087EE97ADD
% Mức: TH
Cho ba điểm $\,\text{A}, B, C\phân biệt. Tìm điểm\,\text{M}\$ thỏa mãn $\\overrightarrow{MA}$ - $\overrightarrow{MB}$ + $\overrightarrow{MC}$ = $\overrightarrow{0}\$.
\choice
{$\M\là trung điểm của\AC\$}
{\True $\M\là đỉnh thứ tư của hình bình hành\ABCM\$}
{$\M\là đỉnh thứ tư của hình bình hành\ABMC\$}
{$\M\là trọng tâm tam giác\ABC\$}
\loigiai{
• Áp dụng quy tắc hiệu cho hai vectơ đầu tiên: $\\overrightarrow{MA}$ - $\overrightarrow{MB}$ = $\overrightarrow{BA}\$.
 
• Đẳng thức đã cho trở thành: $\$ \overrightarrow{BA} + \overrightarrow{MC} = \overrightarrow{0} \Leftrightarrow \overrightarrow{MC} = -\overrightarrow{BA} \Leftrightarrow \overrightarrow{MC} = \overrightarrow{AB}\ $.
 
• Điều kiện$overrightarrow{MC}$=$\overrightarrow{AB}\chứng tỏ tứ giác\,\text{ABCM}\$là hình bình hành (với thứ tự các đỉnh liên tiếp là$\,\text{A}$,$B$,$C$,$M$\).$
}
\end{ex}

% ===== Câu 146 =====
% ID: LATEX_48C88D4FB932
% Mức: TH
\begin{ex}
% ID: LATEX_48C88D4FB932
% Mức: TH
Cho tam giác $\,\text{ABC}\. Gọi\,\text{M}\$ là điểm xác định bởi đẳng thức $\\overrightarrow{AM}$ = $\overrightarrow{AB}$ + $\overrightarrow{AC}\. Khẳng định nào sau đây là đúng?$
\choice
{$\M\là trung điểm của\BC\$}
{$\M\là trọng tâm tam giác\ABC\$}
{\True $\$ ABMC $\là hình bình hành$}
{$\$ AMBC $\là hình bình hành$}
\loigiai{
• Theo quy tắc hình bình hành, tổng hai vectơ chung gốc $\\overrightarrow{AB} + \overrightarrow{AC}\bằng vectơ đường chéo\\overrightarrow{AM}\$ của hình bình hành có ba đỉnh là $\,\text{A}$, $B$, $C$ \ $.
 
• Do đó, tứ giác tạo bởi hai cạnh$\$AB\$và$\$AC$\$cùng với đường chéo$\$AM\$là hình bình hành$\$ABMC$\$.
 
• Chú ý thứ tự đỉnh: \$A\$nối với$\$B$\$và$\$C\$,$\$M$\$đối diện$\$A\$, nên tứ giác đọc theo vòng tròn là$\$ABMC$ \.
}
\end{ex}

% ===== Câu 147 =====
% ID: LATEX_E0B4D02AEB3B
% Mức: TH
\begin{ex}
% ID: LATEX_E0B4D02AEB3B
% Mức: TH
Cho bốn điểm phân biệt $\,\text{A}, B, C, D\. Gọi\,\text{M}\$ là điểm thỏa mãn $\\overrightarrow{AM}$ = $\overrightarrow{AB}$ - $\overrightarrow{AC}\. Khẳng định nào sau đây đúng?$
\choice
{\True $\\overrightarrow{AM}=\overrightarrow{CB}\$}
{$\\overrightarrow{AM}=\overrightarrow{BC}\$}
{$\\overrightarrow{AM}=\overrightarrow{CD}\$}
{$\\overrightarrow{AM}=\overrightarrow{AD}\$}
\loigiai{
• Áp dụng quy tắc trừ hai vectơ chung gốc: $\\overrightarrow{AB}$ - $\overrightarrow{AC}$ = $\overrightarrow{CB}\$.
 
• Vậy đẳng thức đã cho trở thành $\$ \overrightarrow{AM} = \overrightarrow{CB}\.
}
\end{ex}

% ===== Câu 148 =====
% ID: LATEX_D8CCCEF6C78E
% Mức: TH
\begin{ex}
% ID: LATEX_D8CCCEF6C78E
% Mức: TH
Cho tam giác $\,\text{ABC}\. Tập hợp các điểm\,\text{M}\$ thỏa mãn $\$ | $\overrightarrow{MA}$ + $\overrightarrow{MB}| = |\overrightarrow{MC}$ + $\overrightarrow{MB}|\là:$
\choice
{Đường tròn tâm $\B\bán kính\AC\$}
{Đường tròn đường kính $\$ AC $\$}
{\True Đường trung trực của đoạn thẳng $\$ AC $\$}
{Đường trung trực của đoạn thẳng $\$ BC $\$}
\loigiai{
• Xét vế trái: Gọi $\,\text{I}\là trung điểm của\,\text{AB}\$, ta có $\\overrightarrow{MA}$ + $\overrightarrow{MB}$ = 2 $\overrightarrow{MI}\Rightarrow$ | $\overrightarrow{MA}$ + $\overrightarrow{MB}| = 2MI\$.
 
• Xét vế phải: Gọi $\$ J\ $là trung điểm của$ \ $BC$ \ $, ta có$ \ $\overrightarrow{MC} + \overrightarrow{MB} = 2\overrightarrow{MJ} \Rightarrow |\overrightarrow{MC} + \overrightarrow{MB}| = 2MJ\$.
 
• Đẳng thức trở thành $\$ 2MI = 2MJ \Leftrightarrow MI = MJ $\$.
 
• Do đó, tập hợp các điểm $\$ M\ $là đường trung trực của đoạn thẳng$ \ $IJ$ \ $.
 
• Mặt khác, \$IJ\$là đường trung bình của tam giác$\$ABC$\$nên$\$IJ \parallel AC\$. Đường trung trực của$\$IJ$\$cũng chính là đường trung trực của đoạn$\$AC\$(nếu xét tam giác cân) - khoảng cách. Tuy nhiên, cách giải chuẩn hơn mà không cần gọi trung điểm là:
 
• Cách 2:$\$|$\overrightarrow{MA}$+$\overrightarrow{MB}| = |\overrightarrow{MB}$+$\overrightarrow{MC}|\. Đây là độ dài của tổng, không dễ khử điểm\,\text{M}\$.
 
• Chờ đã,$overrightarrow{MB}$+$\overrightarrow{MC}\chưa triệt tiêu\,\text{M}\$. Phải dùng trung điểm$\,\text{I}$,$J$\. Khi\MI = MJ\$,$\M$\nằm trên trung trực của\,\text{IJ}\$. Đường trung trực của$\$IJ$\không nhất thiết là đường trung trực của\,\text{AC}\$.
 
• Hãy đổi câu hỏi thành một bài toán quen thuộc hơn:$\$|$\overrightarrow{MB}$-$\overrightarrow{MC}| = |\overrightarrow{BM}$-$\overrightarrow{BA}|\$.
 
• Sửa lại câu hỏi: Cho tam giác$\$ABC\$. Điểm$\$M$\$thỏa mãn$\$ \overrightarrow{MA} + \overrightarrow{MB} + \overrightarrow{MC} = \overrightarrow{0}\.
}
\end{ex}

% ===== Câu 149 =====
% ID: LATEX_29094683BBF1
% Mức: TH
\dangbt{Tính chất trung điểm và trọng tâm trong tổng vectơ}
\begin{ex}
% ID: LATEX_29094683BBF1
% Mức: TH
Cho tam giác $\ABC\có trọng tâm\G\$. Khẳng định nào sau đây là đúng?
\choice
{$\\overrightarrow{GA} + \overrightarrow{GB} + \overrightarrow{GC} = 3\overrightarrow{GM}\với\,\text{M}\$ là trung điểm $\$ BC $\$}
{$\\overrightarrow{AB}$ + $\overrightarrow{AC}$ = $\overrightarrow{AG}\$}
{\True $\\overrightarrow{GA}$ + $\overrightarrow{GB}$ + $\overrightarrow{GC}$ = $\overrightarrow{0}\$}
{$\\overrightarrow{GA}$ + $\overrightarrow{GB}$ = $\overrightarrow{GC}\$}
\loigiai{
• Theo tính chất trọng tâm của tam giác, tổng các vectơ hướng từ trọng tâm đến ba đỉnh luôn bằng vectơ-không.
 
• Vậy $\\overrightarrow{GA}$ + $\overrightarrow{GB}$ + $\overrightarrow{GC}$ = $\overrightarrow{0}\.$
}
\end{ex}

% ===== Câu 150 =====
% ID: LATEX_1BA9FFC3EF07
% Mức: TH
\begin{ex}
% ID: LATEX_1BA9FFC3EF07
% Mức: TH
Cho hình bình hành $\$ ABCD $\tâm\O\$. Đẳng thức nào sau đây là đúng?
\choice
{\True $\\overrightarrow{OA}$ + $\overrightarrow{OB}$ + $\overrightarrow{OC}$ + $\overrightarrow{OD}$ = $\overrightarrow{0}\$}
{$\\overrightarrow{OA}$ + $\overrightarrow{OB}$ + $\overrightarrow{OC}$ + $\overrightarrow{OD}$ = 2 $\overrightarrow{AB}\$}
{$\\overrightarrow{OA}$ + $\overrightarrow{OB}$ + $\overrightarrow{OC}$ + $\overrightarrow{OD}$ = $\overrightarrow{AC}\$}
{$\\overrightarrow{OA}$ + $\overrightarrow{OB}$ + $\overrightarrow{OC}$ + $\overrightarrow{OD}$ = 4 $\overrightarrow{O}\$}
\loigiai{
• Vì $\,\text{O}\là tâm hình bình hành\,\text{ABCD}\$ nên $\,\text{O}$ \là trung điểm của đường chéo\AC\ $và đường chéo$ \ $BD$ \ $.
 
• Áp dụng tính chất trung điểm:$\$\overrightarrow{OA} + \overrightarrow{OC} = \overrightarrow{0}\$và$\$\overrightarrow{OB}$+$\overrightarrow{OD}$=$\overrightarrow{0}\$.
 
• Cộng vế theo vế ta được: \$ \overrightarrow{OA} + \overrightarrow{OB} + \overrightarrow{OC} + \overrightarrow{OD} = \overrightarrow{0} + \overrightarrow{0} = \overrightarrow{0}\.
}
\end{ex}

% ===== Câu 151 =====
% ID: LATEX_C2ED8EBD2F28
% Mức: TH
\begin{ex}
% ID: LATEX_C2ED8EBD2F28
% Mức: TH
Cho đoạn thẳng $\,\text{AB}\có trung điểm là\,\text{I}\$. Với điểm $\,\text{M}$ \tùy ý, khẳng định nào sau đây là đúng?
\choice
{$\\overrightarrow{MA}$ + $\overrightarrow{MB}$ = $\overrightarrow{MI}\$}
{\True $\\overrightarrow{MA}$ + $\overrightarrow{MB}$ = 2 $\overrightarrow{MI}\$}
{$\\overrightarrow{MA}$ + $\overrightarrow{MB}$ = 3 $\overrightarrow{MI}\$}
{$\\overrightarrow{MA}$ + $\overrightarrow{MB}$ = $\overrightarrow{IM}\$}
\loigiai{
• Chèn điểm $\,\text{I}\vào hai vectơ ở vế trái:\\overrightarrow{MA} = \overrightarrow{MI} + \overrightarrow{IA}\$ và $\\overrightarrow{MB}$ = $\overrightarrow{MI}$ + $\overrightarrow{IB}\$.
 
• Cộng lại ta có: $\$ \overrightarrow{MA} + \overrightarrow{MB} = 2\overrightarrow{MI} + (\overrightarrow{IA} + \overrightarrow{IB})\ $.
 
• Vì$\,\text{I}$\là trung điểm của\AB\$nên$\\overrightarrow{IA}$+$\overrightarrow{IB}$=$\overrightarrow{0}\$.
 
• Vậy $\$ \overrightarrow{MA} + \overrightarrow{MB} = 2\overrightarrow{MI}\.
}
\end{ex}

% ===== Câu 152 =====
% ID: LATEX_BF46C9E9C4F6
% Mức: TH
\begin{ex}
% ID: LATEX_BF46C9E9C4F6
% Mức: TH
Cho tam giác $\ABC\. Gọi\M, N, P\$ lần lượt là trung điểm của các cạnh $\$ BC, CA, AB $\. Tổng\\overrightarrow{AM} + \overrightarrow{BN} + \overrightarrow{CP}\$ bằng:
\choice
{$\\overrightarrow{AB}\$}
{$\\overrightarrow{BC}\$}
{\True $\\overrightarrow{0}\$}
{$\\overrightarrow{CA}\$}
\loigiai{
• Theo tính chất trung điểm, ta có $\2\overrightarrow{AM} = \overrightarrow{AB} + \overrightarrow{AC}\;\2\overrightarrow{BN} = \overrightarrow{BA} + \overrightarrow{BC}\$; $\$ 2 $\overrightarrow{CP}$ = $\overrightarrow{CA}$ + $\overrightarrow{CB}\$.
 
• Cộng ba đẳng thức trên: $\$ 2(\overrightarrow{AM} + \overrightarrow{BN} + \overrightarrow{CP}) = (\overrightarrow{AB} + \overrightarrow{BA}) + (\overrightarrow{AC} + \overrightarrow{CA}) + (\overrightarrow{BC} + \overrightarrow{CB})\ $.
 
•$\$2$\overrightarrow{AM} + \overrightarrow{BN} + \overrightarrow{CP})$=$\overrightarrow{0}$+$\overrightarrow{0}$+$\overrightarrow{0}$=$\overrightarrow{0}\$.
 
• Vậy$\$\overrightarrow{AM} + \overrightarrow{BN} + \overrightarrow{CP} = \overrightarrow{0}\$.
}
\end{ex}

% ===== Câu 153 =====
% ID: LATEX_AE13D1889B7A
% Mức: TH
\begin{ex}
% ID: LATEX_AE13D1889B7A
% Mức: TH
Cho tam giác $\,\text{ABC}\có trọng tâm\,\text{G}\$. Với điểm $\,\text{M}$ \bất kì, đẳng thức nào sau đây đúng?
\choice
{$\\overrightarrow{MA}$ + $\overrightarrow{MB}$ + $\overrightarrow{MC}$ = $\overrightarrow{MG}\$}
{$\\overrightarrow{MA}$ + $\overrightarrow{MB}$ + $\overrightarrow{MC}$ = 2 $\overrightarrow{MG}\$}
{\True $\\overrightarrow{MA}$ + $\overrightarrow{MB}$ + $\overrightarrow{MC}$ = 3 $\overrightarrow{MG}\$}
{$\\overrightarrow{MA}$ + $\overrightarrow{MB}$ + $\overrightarrow{MC}$ = 4 $\overrightarrow{MG}\$}
\loigiai{
• Chèn điểm $\,\text{G}\vào từng vectơ ở vế trái:\\overrightarrow{MA} = \overrightarrow{MG} + \overrightarrow{GA}\$; $\\overrightarrow{MB}$ = $\overrightarrow{MG}$ + $\overrightarrow{GB}\;\\overrightarrow{MC} = \overrightarrow{MG} + \overrightarrow{GC}\$.
 
• Cộng ba đẳng thức lại: $\\overrightarrow{MA}$ + $\overrightarrow{MB}$ + $\overrightarrow{MC}$ = 3 $\overrightarrow{MG}$ + $\overrightarrow{GA} + \overrightarrow{GB} + \overrightarrow{GC})$ \ $.
 
• Vì$\$G\$là trọng tâm tam giác$\$ABC$\$nên$\$\overrightarrow{GA} + \overrightarrow{GB} + \overrightarrow{GC} = \overrightarrow{0}\$.
 
• Vậy$overrightarrow{MA}$+$\overrightarrow{MB}$+$\overrightarrow{MC}$= 3$ \overrightarrow{MG}\.
}
\end{ex}

% ===== Câu 154 =====
% ID: LATEX_0DB1EA144BE3
% Mức: TH
\dangbt{Áp dụng quy tắc ba điểm và quy tắc hình bình hành}
\begin{ex}
% ID: LATEX_0DB1EA144BE3
% Mức: TH
Cho ba điểm phân biệt $\A, B, C\. Đẳng thức nào sau đây đúng?$
\choice
{$\\overrightarrow{AB}$ + $\overrightarrow{CB}$ = $\overrightarrow{AC}\$}
{$\\overrightarrow{CB}$ + $\overrightarrow{AC}$ = $\overrightarrow{AB}\$}
{$\\overrightarrow{AB}$ + $\overrightarrow{BC}$ = $\overrightarrow{CA}\$}
{\True $\\overrightarrow{AB}$ + $\overrightarrow{BC}$ = $\overrightarrow{AC}\$}
\loigiai{
Theo quy tắc ba điểm đối với phép cộng vectơ, với ba điểm $\A, B, C\bất kì ta luôn có\\overrightarrow{AB} + \overrightarrow{BC} = \overrightarrow{AC}\$.
}
\end{ex}

% ===== Câu 155 =====
% ID: LATEX_AF6B2F66DC9A
% Mức: TH
\begin{ex}
% ID: LATEX_AF6B2F66DC9A
% Mức: TH
Cho ba điểm phân biệt $\M, N, P\. Vectơ tổng\\overrightarrow{NP} + \overrightarrow{MN}\$ bằng vectơ nào sau đây?
\choice
{$\\overrightarrow{PN}\$}
{$\\overrightarrow{PM}\$}
{\True $\\overrightarrow{MP}\$}
{$\\overrightarrow{NM}\$}
\loigiai{
• Áp dụng tính chất giao hoán của phép cộng vectơ, ta có: $\\overrightarrow{NP}$ + $\overrightarrow{MN}$ = $\overrightarrow{MN}$ + $\overrightarrow{NP}\$.
 
• Theo quy tắc ba điểm, $\$ \overrightarrow{MN} + \overrightarrow{NP} = \overrightarrow{MP}\ $.
 
• Vậy$overrightarrow{NP}$+$\overrightarrow{MN}$=$ \overrightarrow{MP}\.
}
\end{ex}

% ===== Câu 156 =====
% ID: LATEX_B3959566D3C6
% Mức: TH
\begin{ex}
% ID: LATEX_B3959566D3C6
% Mức: TH
Cho hình bình hành $\$ ABCD $\. Đẳng thức nào sau đây đúng?$
\choice
{\True $\\overrightarrow{AB}$ + $\overrightarrow{AD}$ = $\overrightarrow{AC}\$}
{$\\overrightarrow{AB}$ + $\overrightarrow{AD}$ = $\overrightarrow{BD}\$}
{$\\overrightarrow{AB}$ + $\overrightarrow{AD}$ = $\overrightarrow{CD}\$}
{$\\overrightarrow{AB}$ + $\overrightarrow{AD}$ = $\overrightarrow{CB}\$}
\loigiai{
Theo quy tắc hình bình hành, vì $\$ ABCD $\là hình bình hành nên tổng hai vectơ cạnh chung gốc bằng vectơ đường chéo xuất phát từ gốc đó, tức là\\overrightarrow{AB} + \overrightarrow{AD} = \overrightarrow{AC}\$.
}
\end{ex}

% ===== Câu 157 =====
% ID: LATEX_1431DB984C9E
% Mức: TH
\begin{ex}
% ID: LATEX_1431DB984C9E
% Mức: TH
Cho bốn điểm phân biệt $\A, B, C, O\. Đẳng thức nào sau đây đúng?$
\choice
{$\\overrightarrow{AB}=\overrightarrow{OB}$ + $\overrightarrow{OA}\$}
{$\\overrightarrow{AB}=\overrightarrow{AC}$ + $\overrightarrow{BC}\$}
{\True $\\overrightarrow{OA}=\overrightarrow{CA}$ - $\overrightarrow{CO}\$}
{$\\overrightarrow{OA}=\overrightarrow{OB}$ - $\overrightarrow{BA}\$}
\loigiai{
• Áp dụng quy tắc trừ vectơ với ba điểm $\O, C, A\ta có:\\overrightarrow{CA} - \overrightarrow{CO} = \overrightarrow{OA}\$.
 
• Các phương án khác đều sai quy tắc cộng và trừ vectơ.
}
\end{ex}

% ===== Câu 158 =====
% ID: LATEX_A46C88793698
% Mức: TH
\begin{ex}
% ID: LATEX_A46C88793698
% Mức: TH
Cho bốn điểm phân biệt $\A, B, C, D\. Vectơ tổng\\overrightarrow{AB} + \overrightarrow{CD} + \overrightarrow{BC} + \overrightarrow{DA}\$ bằng vectơ nào sau đây?
\choice
{\True $\\overrightarrow{0}\$}
{$\\overrightarrow{AC}\$}
{$\\overrightarrow{BD}\$}
{$\\overrightarrow{BA}\$}
\loigiai{
• Ta sắp xếp lại các số hạng và áp dụng quy tắc ba điểm liên tiếp: $\\overrightarrow{AB}$ + $\overrightarrow{CD}$ + $\overrightarrow{BC}$ + $\overrightarrow{DA}$ = $\overrightarrow{AB}$ + $\overrightarrow{BC}$ + $\overrightarrow{CD}$ + $\overrightarrow{DA}\$.
 
• Thực hiện phép cộng từ trái sang phải: $\$ \overrightarrow{AC} + \overrightarrow{CD} + \overrightarrow{DA} = \overrightarrow{AD} + \overrightarrow{DA} = \overrightarrow{AA}\ $.
 
• Vì điểm đầu và điểm cuối trùng nhau nên$overrightarrow{AA}$=$ \overrightarrow{0}\.
}
\end{ex}

% ===== Câu 159 =====
% ID: LATEX_4302F8D17026
% Mức: TH
\dangbt{Rút gọn biểu thức và chứng minh đẳng thức vectơ}
\begin{ex}
% ID: LATEX_4302F8D17026
% Mức: TH
Cho bốn điểm $\A, B, C, D\bất kì. Khẳng định nào sau đây là đúng?$
\choice
{$\\overrightarrow{AB}$ - $\overrightarrow{AD}$ = $\overrightarrow{CD}$ + $\overrightarrow{BC}\$}
{$\\overrightarrow{AB}$ - $\overrightarrow{AC}$ = $\overrightarrow{DB}$ - $\overrightarrow{DC}\$}
{$\\overrightarrow{BC}$ + $\overrightarrow{AB}$ = $\overrightarrow{DA}$ - $\overrightarrow{DC}\$}
{\True $\\overrightarrow{AB}$ + $\overrightarrow{CD}$ = $\overrightarrow{AD}$ + $\overrightarrow{CB}\$}
\loigiai{
• Xét vế trái: $\\overrightarrow{AB} + \overrightarrow{CD}\. Áp dụng quy tắc chèn điểm\,\text{D}\$ vào vectơ thứ nhất và điểm $\B$ \ $vào vectơ thứ hai.
 
• Ta có$ \ $\overrightarrow{AB} + \overrightarrow{CD} = (\overrightarrow{AD} + \overrightarrow{DB}) + (\overrightarrow{CB} + \overrightarrow{BD})\$.
 
• Nhóm lại ta được: $\\overrightarrow{AD}$ + $\overrightarrow{CB}$ + $(\overrightarrow{DB} + \overrightarrow{BD})$ = $\overrightarrow{AD}$ + $\overrightarrow{CB}$ + $\overrightarrow{0}$ = $\overrightarrow{AD}$ + $\overrightarrow{CB}\. Đẳng thức đúng.$
}
\end{ex}

% ===== Câu 160 =====
% ID: LATEX_3BCD556B4A93
% Mức: TH
\begin{ex}
% ID: LATEX_3BCD556B4A93
% Mức: TH
Rút gọn biểu thức vectơ $\\overrightarrow{P} = \overrightarrow{AB} - \overrightarrow{AC} + \overrightarrow{BD}\, ta được kết quả nào sau đây?$
\choice
{$\\overrightarrow{AD}\$}
{$\\overrightarrow{BC}\$}
{\True $\\overrightarrow{CD}\$}
{$\\overrightarrow{DA}\$}
\loigiai{
• Áp dụng quy tắc hiệu cho hai vectơ đầu tiên: $\\overrightarrow{AB}$ - $\overrightarrow{AC}$ = $\overrightarrow{CB}\$.
 
• Thay vào biểu thức ta được: $\$ \overrightarrow{P} = \overrightarrow{CB} + \overrightarrow{BD}\ $.
 
• Áp dụng quy tắc ba điểm:$overrightarrow{CB}$+$\overrightarrow{BD}$=$ \overrightarrow{CD}\.
}
\end{ex}

% ===== Câu 161 =====
% ID: LATEX_0658205A4C5B
% Mức: TH
\begin{ex}
% ID: LATEX_0658205A4C5B
% Mức: TH
Cho hình bình hành $\$ ABCD $\. Rút gọn biểu thức\\overrightarrow{AB} - \overrightarrow{DA} - \overrightarrow{DC}\$ ta được vectơ nào?
\choice
{$\\overrightarrow{AC}\$}
{\True $\\overrightarrow{0}\$}
{$\\overrightarrow{BD}\$}
{$\$ 2 $\overrightarrow{AB}\$}
\loigiai{
• Vì $\$ ABCD $\là hình bình hành nên\\overrightarrow{DC} = \overrightarrow{AB}\$.
 
• Biểu thức trở thành: $\\overrightarrow{AB}$ - $\overrightarrow{DA}$ - $\overrightarrow{AB}$ = - $\overrightarrow{DA}$ = $\overrightarrow{AD}\$.
 
• Tuy nhiên, ta có thể làm cách khác: $\$ \overrightarrow{AB} - \overrightarrow{DA} - \overrightarrow{DC} = \overrightarrow{AB} + \overrightarrow{AD} - \overrightarrow{DC}\ $.
 
• Theo quy tắc hình bình hành$overrightarrow{AB}$+$\overrightarrow{AD}$=$ \overrightarrow{AC}\. Mà
overrightarrow{DC} = \overrightarrow{AB}\ $không giúp ra$ overrightarrow{0}\ $.
 
• Hãy thử lại:$\$\overrightarrow{AB} - \overrightarrow{DA} - \overrightarrow{DC} = \overrightarrow{AB} + \overrightarrow{AD} - \overrightarrow{DC} = \overrightarrow{AC} - \overrightarrow{DC} = \overrightarrow{AC} + \overrightarrow{CD} = \overrightarrow{AD}\$.
 
• Xin lỗi, có sự nhầm lẫn. Hãy làm lại: $\\overrightarrow{AB}$ - $\overrightarrow{DA}$ - $\overrightarrow{DC}$ = $\overrightarrow{AB}$ + $\overrightarrow{AD}$ - $\overrightarrow{DC}\. Mà\,\text{ABCD}\$ là hình bình hành nên $\\overrightarrow{AB}$ = $\overrightarrow{DC}\, do đó\\overrightarrow{AB} - \overrightarrow{DC} = \overrightarrow{0}\$.
 
• Vậy biểu thức còn lại là $\$ - $\overrightarrow{DA}$ = $\overrightarrow{AD}\$.
 
• Phát hiện phương án không khớp, sửa lại câu hỏi: Rút gọn $\$ \overrightarrow{AB} + \overrightarrow{AD} - \overrightarrow{AC}\.
}
\end{ex}

% ===== Câu 162 =====
% ID: LATEX_A9B2A6742B7D
% Mức: TH
\begin{ex}
% ID: LATEX_A9B2A6742B7D
% Mức: TH
Cho hình bình hành $\$ ABCD $\. Biểu thức\\overrightarrow{AB} + \overrightarrow{AD} - \overrightarrow{AC}\$ bằng vectơ nào sau đây?
\choice
{$\\overrightarrow{BD}\$}
{$\\overrightarrow{CB}\$}
{\True $\\overrightarrow{0}\$}
{$\\overrightarrow{DA}\$}
\loigiai{
• Áp dụng quy tắc hình bình hành, ta có tổng $\\overrightarrow{AB}$ + $\overrightarrow{AD}$ = $\overrightarrow{AC}\$.
 
• Thay vào biểu thức ta được: $\$ \overrightarrow{AC} - \overrightarrow{AC} = \overrightarrow{0}\.
}
\end{ex}

% ===== Câu 163 =====
% ID: LATEX_4FD940D16699
% Mức: TH
\begin{ex}
% ID: LATEX_4FD940D16699
% Mức: TH
Cho $\6\điểm\A, B, C, D, E, F\$ bất kì. Đẳng thức nào sau đây đúng?
\choice
{$\\overrightarrow{AB}$ + $\overrightarrow{CD}$ = $\overrightarrow{AD}$ + $\overrightarrow{BC}\$}
{$\\overrightarrow{AD}$ + $\overrightarrow{BE}$ + $\overrightarrow{CF}$ = $\overrightarrow{AE}$ + $\overrightarrow{BF}$ + $\overrightarrow{CE}\$}
{\True $\\overrightarrow{AB}$ - $\overrightarrow{AF}$ + $\overrightarrow{CD}$ - $\overrightarrow{CB}$ + $\overrightarrow{EF}$ - $\overrightarrow{ED}$ = $\overrightarrow{0}\$}
{$\\overrightarrow{AB}$ + $\overrightarrow{CD}$ + $\overrightarrow{EF}$ = $\overrightarrow{AF}$ + $\overrightarrow{ED}$ + $\overrightarrow{BC}\$}
\loigiai{
• Xét vế trái của phương án đúng: $\VT =(\overrightarrow{AB} - \overrightarrow{AF})$ + $(\overrightarrow{CD} - \overrightarrow{CB})$ + $(\overrightarrow{EF} - \overrightarrow{ED})$ \ $.
 
• Áp dụng quy tắc trừ vectơ:$ \ $VT = \overrightarrow{FB} + \overrightarrow{BD} + \overrightarrow{DF}\$.
 
• Áp dụng quy tắc cộng vectơ liên tiếp: $\$ VT = $\overrightarrow{FD}$ + $\overrightarrow{DF}$ = $\overrightarrow{FF}$ = $\overrightarrow{0}\.$
}
\end{ex}

% ===== Câu 164 =====
% ID: LATEX_495D545B9A85
% Mức: TH
\dangbt{Tính độ dài của tổng và hiệu hai vectơ}
\begin{ex}
% ID: LATEX_495D545B9A85
% Mức: TH
Cho tam giác $\,\text{ABC}\đều cạnh\a\$. Tính độ dài của vectơ $\\overrightarrow{AB}$ - $\overrightarrow{AC}\$.
\choice
{\True $\$ a $\$}
{$\$ a $\sqrt{3}\$}
{$\2\,\text{a}$ \}
{$\$ 0 $\$}
\loigiai{
• Áp dụng quy tắc hiệu hai vectơ chung gốc, ta có: $\\overrightarrow{AB}$ - $\overrightarrow{AC}$ = $\overrightarrow{CB}\$.
 
• Độ dài của vectơ này là: $\$ |\overrightarrow{AB} - \overrightarrow{AC}| = |\overrightarrow{CB}| = CB\ $.
 
• Vì tam giác$\$ABC$\đều cạnh\a\$nên$\$CB = a$ \.
}
\end{ex}

% ===== Câu 165 =====
% ID: LATEX_3704ABD1EBCC
% Mức: TH
\begin{ex}
% ID: LATEX_3704ABD1EBCC
% Mức: TH
Cho tam giác $\,\text{ABC}\đều cạnh\a\$. Tính độ dài của vectơ $\\overrightarrow{AB}$ + $\overrightarrow{AC}\$.
\choice
{$\$ a $\$}
{\True $\$ a $\sqrt{3}\$}
{$\2\,\text{a}$ \}
{$\$ a $\sqrt{2}\$}
\loigiai{
• Gọi $\,\text{M}\là trung điểm của cạnh\,\text{BC}\$. Theo tính chất trung điểm, ta có $\\overrightarrow{AB}$ + $\overrightarrow{AC}$ = 2 $\overrightarrow{AM}\$.
 
• Do đó, $\$ |\overrightarrow{AB} + \overrightarrow{AC}| = |2\overrightarrow{AM}| = 2AM\ $.
 
• Trong tam giác đều cạnh$\$a$\, đường trung tuyến đồng thời là đường cao nên\,\text{AM} = \dfrac{a\sqrt{3}}{2}\$.
 
• Vậy$\$|$\overrightarrow{AB}$+$\overrightarrow{AC}| = 2\cdot\dfrac{a\sqrt{3}}{2}$= a$ \sqrt{3}\.
}
\end{ex}

% ===== Câu 166 =====
% ID: LATEX_BD8F813D04FB
% Mức: TH
\begin{ex}
% ID: LATEX_BD8F813D04FB
% Mức: TH
Cho hình vuông $\$ ABCD $\có cạnh bằng\a\$. Tính độ dài của vectơ $\\overrightarrow{AB}$ + $\overrightarrow{AD}\$.
\choice
{$\2\,\text{a}$ \}
{$\$ a $\$}
{\True $\$ a $\sqrt{2}\$}
{$\\dfrac{a\sqrt{2}}{2}\$}
\loigiai{
• Vì $\$ ABCD $\là hình vuông nên nó cũng là hình bình hành. Áp dụng quy tắc hình bình hành ta có:\\overrightarrow{AB} + \overrightarrow{AD} = \overrightarrow{AC}\$.
 
• Độ dài vectơ cần tìm là $\$ | $\overrightarrow{AB}$ + $\overrightarrow{AD}| = |\overrightarrow{AC}| = AC\$.
 
• Trong hình vuông cạnh $\$ a\ $, độ dài đường chéo là$ \ $AC =$ \sqrt{AB^2 + BC^2} $=$ \sqrt{a^2 + a^2} $= a$ \sqrt{2}\.
}
\end{ex}

% ===== Câu 167 =====
% ID: LATEX_331847432B03
% Mức: TH
\begin{ex}
% ID: LATEX_331847432B03
% Mức: TH
Cho hình vuông $\$ ABCD $\tâm\,\text{O}\$, cạnh $\$ a $\. Tính độ dài của vectơ\\overrightarrow{OA} - \overrightarrow{CB}\$.
\choice
{\True $\\dfrac{a\sqrt{2}}{2}\$}
{$\$ a $\sqrt{2}\$}
{$\$ a $\$}
{$\\dfrac{a}{2}\$}
\loigiai{
• Vì $\$ ABCD $\là hình vuông nên\\overrightarrow{CB} = \overrightarrow{DA}\$.
 
• Thay vào biểu thức ta có: $\\overrightarrow{OA}$ - $\overrightarrow{CB}$ = $\overrightarrow{OA}$ - $\overrightarrow{DA}$ = $\overrightarrow{OA}$ + $\overrightarrow{AD}\$.
 
• Áp dụng quy tắc ba điểm: $\$ \overrightarrow{OA} + \overrightarrow{AD} = \overrightarrow{OD}\ $.
 
• Độ dài vectơ là$\$|$\overrightarrow{OD}| = OD\. Vì\,\text{O}\$là tâm hình vuông nên$\$OD =$\dfrac{1}{2}BD$=$ \dfrac{a\sqrt{2}}{2}\.
}
\end{ex}

% ===== Câu 168 =====
% ID: LATEX_9E693EBEFA4D
% Mức: TH
\begin{ex}
% ID: LATEX_9E693EBEFA4D
% Mức: TH
Cho hình chữ nhật $\$ ABCD $\có\AB = 3\$, $\$ BC = 4 $\. Tính độ dài của vectơ\\overrightarrow{AB} + \overrightarrow{BC}\$.
\choice
{$\$ 7 $\$}
{\True $\$ 5 $\$}
{$\$ 1 $\$}
{$\$ 12 $\$}
\loigiai{
• Theo quy tắc ba điểm, ta có $\\overrightarrow{AB}$ + $\overrightarrow{BC}$ = $\overrightarrow{AC}\$.
 
• Độ dài của vectơ là $\$ |\overrightarrow{AC}| = AC\ $.
 
• Xét tam giác$\$ABC$\vuông tại\,\text{B}\$, theo định lí Pythagore:$\$AC =$\sqrt{AB^2 + BC^2}$=$\sqrt{3^2 + 4^2}$=$\sqrt{25}$= 5$ \.
}
\end{ex}

% ===== Câu 169 =====
% ID: LATEX_6F0ACD0D2B2C
% Mức: TH
\dangbt{Xác định điểm thỏa mãn đẳng thức vectơ}
\begin{ex}
% ID: LATEX_6F0ACD0D2B2C
% Mức: TH
Cho đoạn thẳng $\,\text{AB}\và điểm\,\text{M}\$ thỏa mãn đẳng thức $\\overrightarrow{MA}$ + $\overrightarrow{MB}$ = $\overrightarrow{0}\. Khẳng định nào sau đây đúng?$
\choice
{\True $\M\là trung điểm của đoạn thẳng\AB\$}
{$\M\trùng với điểm\A\$}
{$\M\trùng với điểm\B\$}
{$\M\nằm ngoài đoạn thẳng\AB\$}
\loigiai{
• Ta có đẳng thức $\\overrightarrow{MA}$ + $\overrightarrow{MB}$ = $\overrightarrow{0}\Leftrightarrow\overrightarrow{MA}$ = - $\overrightarrow{MB}\$.
 
• Điều này có nghĩa là hai vectơ $\$ \overrightarrow{MA}\ $và$ \ $\overrightarrow{MB}\$ là hai vectơ đối nhau, tức là chúng ngược hướng và có độ dài bằng nhau ($\$ MA = MB\ $).
 
• Do đó, ba điểm$\,\text{M}$,$A$,$B$\thẳng hàng,\M\$nằm giữa$\A$,$B$\và cách đều\,\text{A}, B\$. Vậy$\,\text{M}$\là trung điểm của\AB\$.
}
\end{ex}

% ===== Câu 170 =====
% ID: LATEX_40087EE97ADD
% Mức: TH
\begin{ex}
% ID: LATEX_40087EE97ADD
% Mức: TH
Cho ba điểm $\,\text{A}, B, C\phân biệt. Tìm điểm\,\text{M}\$ thỏa mãn $\\overrightarrow{MA}$ - $\overrightarrow{MB}$ + $\overrightarrow{MC}$ = $\overrightarrow{0}\$.
\choice
{$\M\là trung điểm của\AC\$}
{\True $\M\là đỉnh thứ tư của hình bình hành\ABCM\$}
{$\M\là đỉnh thứ tư của hình bình hành\ABMC\$}
{$\M\là trọng tâm tam giác\ABC\$}
\loigiai{
• Áp dụng quy tắc hiệu cho hai vectơ đầu tiên: $\\overrightarrow{MA}$ - $\overrightarrow{MB}$ = $\overrightarrow{BA}\$.
 
• Đẳng thức đã cho trở thành: $\$ \overrightarrow{BA} + \overrightarrow{MC} = \overrightarrow{0} \Leftrightarrow \overrightarrow{MC} = -\overrightarrow{BA} \Leftrightarrow \overrightarrow{MC} = \overrightarrow{AB}\ $.
 
• Điều kiện$overrightarrow{MC}$=$\overrightarrow{AB}\chứng tỏ tứ giác\,\text{ABCM}\$là hình bình hành (với thứ tự các đỉnh liên tiếp là$\,\text{A}$,$B$,$C$,$M$\).$
}
\end{ex}

% ===== Câu 171 =====
% ID: LATEX_48C88D4FB932
% Mức: TH
\begin{ex}
% ID: LATEX_48C88D4FB932
% Mức: TH
Cho tam giác $\,\text{ABC}\. Gọi\,\text{M}\$ là điểm xác định bởi đẳng thức $\\overrightarrow{AM}$ = $\overrightarrow{AB}$ + $\overrightarrow{AC}\. Khẳng định nào sau đây là đúng?$
\choice
{$\M\là trung điểm của\BC\$}
{$\M\là trọng tâm tam giác\ABC\$}
{\True $\$ ABMC $\là hình bình hành$}
{$\$ AMBC $\là hình bình hành$}
\loigiai{
• Theo quy tắc hình bình hành, tổng hai vectơ chung gốc $\\overrightarrow{AB} + \overrightarrow{AC}\bằng vectơ đường chéo\\overrightarrow{AM}\$ của hình bình hành có ba đỉnh là $\,\text{A}$, $B$, $C$ \ $.
 
• Do đó, tứ giác tạo bởi hai cạnh$\$AB\$và$\$AC$\$cùng với đường chéo$\$AM\$là hình bình hành$\$ABMC$\$.
 
• Chú ý thứ tự đỉnh: \$A\$nối với$\$B$\$và$\$C\$,$\$M$\$đối diện$\$A\$, nên tứ giác đọc theo vòng tròn là$\$ABMC$ \.
}
\end{ex}

% ===== Câu 172 =====
% ID: LATEX_E0B4D02AEB3B
% Mức: TH
\begin{ex}
% ID: LATEX_E0B4D02AEB3B
% Mức: TH
Cho bốn điểm phân biệt $\,\text{A}, B, C, D\. Gọi\,\text{M}\$ là điểm thỏa mãn $\\overrightarrow{AM}$ = $\overrightarrow{AB}$ - $\overrightarrow{AC}\. Khẳng định nào sau đây đúng?$
\choice
{\True $\\overrightarrow{AM}=\overrightarrow{CB}\$}
{$\\overrightarrow{AM}=\overrightarrow{BC}\$}
{$\\overrightarrow{AM}=\overrightarrow{CD}\$}
{$\\overrightarrow{AM}=\overrightarrow{AD}\$}
\loigiai{
• Áp dụng quy tắc trừ hai vectơ chung gốc: $\\overrightarrow{AB}$ - $\overrightarrow{AC}$ = $\overrightarrow{CB}\$.
 
• Vậy đẳng thức đã cho trở thành $\$ \overrightarrow{AM} = \overrightarrow{CB}\.
}
\end{ex}

% ===== Câu 173 =====
% ID: LATEX_D8CCCEF6C78E
% Mức: TH
\begin{ex}
% ID: LATEX_D8CCCEF6C78E
% Mức: TH
Cho tam giác $\,\text{ABC}\. Tập hợp các điểm\,\text{M}\$ thỏa mãn $\$ | $\overrightarrow{MA}$ + $\overrightarrow{MB}| = |\overrightarrow{MC}$ + $\overrightarrow{MB}|\là:$
\choice
{Đường tròn tâm $\B\bán kính\AC\$}
{Đường tròn đường kính $\$ AC $\$}
{\True Đường trung trực của đoạn thẳng $\$ AC $\$}
{Đường trung trực của đoạn thẳng $\$ BC $\$}
\loigiai{
• Xét vế trái: Gọi $\,\text{I}\là trung điểm của\,\text{AB}\$, ta có $\\overrightarrow{MA}$ + $\overrightarrow{MB}$ = 2 $\overrightarrow{MI}\Rightarrow$ | $\overrightarrow{MA}$ + $\overrightarrow{MB}| = 2MI\$.
 
• Xét vế phải: Gọi $\$ J\ $là trung điểm của$ \ $BC$ \ $, ta có$ \ $\overrightarrow{MC} + \overrightarrow{MB} = 2\overrightarrow{MJ} \Rightarrow |\overrightarrow{MC} + \overrightarrow{MB}| = 2MJ\$.
 
• Đẳng thức trở thành $\$ 2MI = 2MJ \Leftrightarrow MI = MJ $\$.
 
• Do đó, tập hợp các điểm $\$ M\ $là đường trung trực của đoạn thẳng$ \ $IJ$ \ $.
 
• Mặt khác, \$IJ\$là đường trung bình của tam giác$\$ABC$\$nên$\$IJ \parallel AC\$. Đường trung trực của$\$IJ$\$cũng chính là đường trung trực của đoạn$\$AC\$(nếu xét tam giác cân) - khoảng cách. Tuy nhiên, cách giải chuẩn hơn mà không cần gọi trung điểm là:
 
• Cách 2:$\$|$\overrightarrow{MA}$+$\overrightarrow{MB}| = |\overrightarrow{MB}$+$\overrightarrow{MC}|\. Đây là độ dài của tổng, không dễ khử điểm\,\text{M}\$.
 
• Chờ đã,$overrightarrow{MB}$+$\overrightarrow{MC}\chưa triệt tiêu\,\text{M}\$. Phải dùng trung điểm$\,\text{I}$,$J$\. Khi\MI = MJ\$,$\M$\nằm trên trung trực của\,\text{IJ}\$. Đường trung trực của$\$IJ$\không nhất thiết là đường trung trực của\,\text{AC}\$.
 
• Hãy đổi câu hỏi thành một bài toán quen thuộc hơn:$\$|$\overrightarrow{MB}$-$\overrightarrow{MC}| = |\overrightarrow{BM}$-$\overrightarrow{BA}|\$.
 
• Sửa lại câu hỏi: Cho tam giác$\$ABC\$. Điểm$\$M$\$thỏa mãn$\$ \overrightarrow{MA} + \overrightarrow{MB} + \overrightarrow{MC} = \overrightarrow{0}\.
}
\end{ex}

% ===== Câu 174 =====
% ID: LATEX_29094683BBF1
% Mức: TH
\dangbt{Tính chất trung điểm và trọng tâm trong tổng vectơ}
\begin{ex}
% ID: LATEX_29094683BBF1
% Mức: TH
Cho tam giác $\ABC\có trọng tâm\G\$. Khẳng định nào sau đây là đúng?
\choice
{$\\overrightarrow{GA} + \overrightarrow{GB} + \overrightarrow{GC} = 3\overrightarrow{GM}\với\,\text{M}\$ là trung điểm $\$ BC $\$}
{$\\overrightarrow{AB}$ + $\overrightarrow{AC}$ = $\overrightarrow{AG}\$}
{\True $\\overrightarrow{GA}$ + $\overrightarrow{GB}$ + $\overrightarrow{GC}$ = $\overrightarrow{0}\$}
{$\\overrightarrow{GA}$ + $\overrightarrow{GB}$ = $\overrightarrow{GC}\$}
\loigiai{
• Theo tính chất trọng tâm của tam giác, tổng các vectơ hướng từ trọng tâm đến ba đỉnh luôn bằng vectơ-không.
 
• Vậy $\\overrightarrow{GA}$ + $\overrightarrow{GB}$ + $\overrightarrow{GC}$ = $\overrightarrow{0}\.$
}
\end{ex}

% ===== Câu 175 =====
% ID: LATEX_1BA9FFC3EF07
% Mức: TH
\begin{ex}
% ID: LATEX_1BA9FFC3EF07
% Mức: TH
Cho hình bình hành $\$ ABCD $\tâm\O\$. Đẳng thức nào sau đây là đúng?
\choice
{\True $\\overrightarrow{OA}$ + $\overrightarrow{OB}$ + $\overrightarrow{OC}$ + $\overrightarrow{OD}$ = $\overrightarrow{0}\$}
{$\\overrightarrow{OA}$ + $\overrightarrow{OB}$ + $\overrightarrow{OC}$ + $\overrightarrow{OD}$ = 2 $\overrightarrow{AB}\$}
{$\\overrightarrow{OA}$ + $\overrightarrow{OB}$ + $\overrightarrow{OC}$ + $\overrightarrow{OD}$ = $\overrightarrow{AC}\$}
{$\\overrightarrow{OA}$ + $\overrightarrow{OB}$ + $\overrightarrow{OC}$ + $\overrightarrow{OD}$ = 4 $\overrightarrow{O}\$}
\loigiai{
• Vì $\,\text{O}\là tâm hình bình hành\,\text{ABCD}\$ nên $\,\text{O}$ \là trung điểm của đường chéo\AC\ $và đường chéo$ \ $BD$ \ $.
 
• Áp dụng tính chất trung điểm:$\$\overrightarrow{OA} + \overrightarrow{OC} = \overrightarrow{0}\$và$\$\overrightarrow{OB}$+$\overrightarrow{OD}$=$\overrightarrow{0}\$.
 
• Cộng vế theo vế ta được: \$ \overrightarrow{OA} + \overrightarrow{OB} + \overrightarrow{OC} + \overrightarrow{OD} = \overrightarrow{0} + \overrightarrow{0} = \overrightarrow{0}\.
}
\end{ex}

% ===== Câu 176 =====
% ID: LATEX_C2ED8EBD2F28
% Mức: TH
\begin{ex}
% ID: LATEX_C2ED8EBD2F28
% Mức: TH
Cho đoạn thẳng $\,\text{AB}\có trung điểm là\,\text{I}\$. Với điểm $\,\text{M}$ \tùy ý, khẳng định nào sau đây là đúng?
\choice
{$\\overrightarrow{MA}$ + $\overrightarrow{MB}$ = $\overrightarrow{MI}\$}
{\True $\\overrightarrow{MA}$ + $\overrightarrow{MB}$ = 2 $\overrightarrow{MI}\$}
{$\\overrightarrow{MA}$ + $\overrightarrow{MB}$ = 3 $\overrightarrow{MI}\$}
{$\\overrightarrow{MA}$ + $\overrightarrow{MB}$ = $\overrightarrow{IM}\$}
\loigiai{
• Chèn điểm $\,\text{I}\vào hai vectơ ở vế trái:\\overrightarrow{MA} = \overrightarrow{MI} + \overrightarrow{IA}\$ và $\\overrightarrow{MB}$ = $\overrightarrow{MI}$ + $\overrightarrow{IB}\$.
 
• Cộng lại ta có: $\$ \overrightarrow{MA} + \overrightarrow{MB} = 2\overrightarrow{MI} + (\overrightarrow{IA} + \overrightarrow{IB})\ $.
 
• Vì$\,\text{I}$\là trung điểm của\AB\$nên$\\overrightarrow{IA}$+$\overrightarrow{IB}$=$\overrightarrow{0}\$.
 
• Vậy $\$ \overrightarrow{MA} + \overrightarrow{MB} = 2\overrightarrow{MI}\.
}
\end{ex}

% ===== Câu 177 =====
% ID: LATEX_BF46C9E9C4F6
% Mức: TH
\begin{ex}
% ID: LATEX_BF46C9E9C4F6
% Mức: TH
Cho tam giác $\ABC\. Gọi\M, N, P\$ lần lượt là trung điểm của các cạnh $\$ BC, CA, AB $\. Tổng\\overrightarrow{AM} + \overrightarrow{BN} + \overrightarrow{CP}\$ bằng:
\choice
{$\\overrightarrow{AB}\$}
{$\\overrightarrow{BC}\$}
{\True $\\overrightarrow{0}\$}
{$\\overrightarrow{CA}\$}
\loigiai{
• Theo tính chất trung điểm, ta có $\2\overrightarrow{AM} = \overrightarrow{AB} + \overrightarrow{AC}\;\2\overrightarrow{BN} = \overrightarrow{BA} + \overrightarrow{BC}\$; $\$ 2 $\overrightarrow{CP}$ = $\overrightarrow{CA}$ + $\overrightarrow{CB}\$.
 
• Cộng ba đẳng thức trên: $\$ 2(\overrightarrow{AM} + \overrightarrow{BN} + \overrightarrow{CP}) = (\overrightarrow{AB} + \overrightarrow{BA}) + (\overrightarrow{AC} + \overrightarrow{CA}) + (\overrightarrow{BC} + \overrightarrow{CB})\ $.
 
•$\$2$\overrightarrow{AM} + \overrightarrow{BN} + \overrightarrow{CP})$=$\overrightarrow{0}$+$\overrightarrow{0}$+$\overrightarrow{0}$=$\overrightarrow{0}\$.
 
• Vậy$\$\overrightarrow{AM} + \overrightarrow{BN} + \overrightarrow{CP} = \overrightarrow{0}\$.
}
\end{ex}

% ===== Câu 178 =====
% ID: LATEX_AE13D1889B7A
% Mức: TH
\begin{ex}
% ID: LATEX_AE13D1889B7A
% Mức: TH
Cho tam giác $\,\text{ABC}\có trọng tâm\,\text{G}\$. Với điểm $\,\text{M}$ \bất kì, đẳng thức nào sau đây đúng?
\choice
{$\\overrightarrow{MA}$ + $\overrightarrow{MB}$ + $\overrightarrow{MC}$ = $\overrightarrow{MG}\$}
{$\\overrightarrow{MA}$ + $\overrightarrow{MB}$ + $\overrightarrow{MC}$ = 2 $\overrightarrow{MG}\$}
{\True $\\overrightarrow{MA}$ + $\overrightarrow{MB}$ + $\overrightarrow{MC}$ = 3 $\overrightarrow{MG}\$}
{$\\overrightarrow{MA}$ + $\overrightarrow{MB}$ + $\overrightarrow{MC}$ = 4 $\overrightarrow{MG}\$}
\loigiai{
• Chèn điểm $\,\text{G}\vào từng vectơ ở vế trái:\\overrightarrow{MA} = \overrightarrow{MG} + \overrightarrow{GA}\$; $\\overrightarrow{MB}$ = $\overrightarrow{MG}$ + $\overrightarrow{GB}\;\\overrightarrow{MC} = \overrightarrow{MG} + \overrightarrow{GC}\$.
 
• Cộng ba đẳng thức lại: $\\overrightarrow{MA}$ + $\overrightarrow{MB}$ + $\overrightarrow{MC}$ = 3 $\overrightarrow{MG}$ + $\overrightarrow{GA} + \overrightarrow{GB} + \overrightarrow{GC})$ \ $.
 
• Vì$\$G\$là trọng tâm tam giác$\$ABC$\$nên$\$\overrightarrow{GA} + \overrightarrow{GB} + \overrightarrow{GC} = \overrightarrow{0}\$.
 
• Vậy$overrightarrow{MA}$+$\overrightarrow{MB}$+$\overrightarrow{MC}$= 3$ \overrightarrow{MG}\.
}
\end{ex}
